\chapter{Representations of finite groups and the mod \texorpdfstring{$\ell$}{l} world}\label{ch:reps}
\chaptermark{Representations and the mod $\ell$ world}

The objects of the paper are representations: $V=V_\ell E$, $T=T_\ell E$, $\rho=\rho_{E,\ell}$ and its reduction $\rhobar$ are
representations of the absolute Galois group $G_\Q$, and the adjoint $A_f=\ad V$ is built from them by linear algebra. Most of
what the paper needs about them is not Galois theory at all but the representation theory of a group, often a finite one, over a
field of characteristic $0$ or $\ell$, or over $\Zl$. This chapter supplies that theory: irreducibility and absolute
irreducibility, Schur's lemma, Maschke's theorem and its failure in characteristic $\ell$, the theorem of Brauer and Nesbitt,
induced representations and Clifford theory for a subgroup of index $2$, twisting, the adjoint representation and its trace form,
the representations of $\GL_2(\F_3)$ and of its subgroups on $\F_3^2$ and $\F_9^2$, the action of finite cyclic groups and of
unipotent elements on rank-two spaces and lattices, and lattices over a discrete valuation ring. Every statement is proved for an
abstract group; the Galois input (continuity, Chebotarev, complex conjugation, the cyclotomic character) is added in
\Cref{ch:galois} and \Cref{ch:galrep}, which cite the results proved here. The chapter uses \Cref{ch:groups} for the group theory
of $\GL_2(\F_3)$, \Cref{ch:rings} for finite fields, and \Cref{ch:modules} for modules over a discrete valuation ring and the
Jordan-H\"older theorem.

\section*{Core results at a glance}

Throughout, $G$ is a group, $k$ a field, $\bar k$ an algebraic closure, $R$ a commutative ring, and representations are
finite-dimensional (\Cref{conv:reps:conventions}). $\mathcal{O}$ is a discrete valuation ring with uniformiser $\varpi$, fraction field
$K$ and residue field $\mathcal{O}/\varpi$. The matrices $J=\sm01{-1}0$, $\sigma=\mathrm{diag}(1,-1)$, $w=\sm0110$, $u=\sm1101$,
$a=\sm1{-1}11$ and the subgroups of $\GL_2(\F_3)$ are the standard models of \Cref{ch:groups}; $\F_9=\F_3[i]$.

\begin{enumerate}[label=\textbf{\thechapter.\arabic*}]
\item \textbf{Conventions.} Matrices act on column vectors on the left; $V^\vee=\Hom_k(V,k)$ with $(g\lambda)(v)=\lambda(g^{-1}v)$ (the
symbol ${}^*$ is reserved for Pontryagin duals); $\End(V)=V\otimes V^\vee$ with $g\cdot x=\rho(g)x\rho(g)^{-1}$, $\ad V$ its trace-zero
part, with basis $e=e_1\otimes e_2^\vee$, $h=e_1\otimes e_1^\vee-e_2\otimes e_2^\vee$, $f=e_2\otimes e_1^\vee$ when $\dim V=2$; a character is a
homomorphism $\chi\colon G\to R^\times$, $M(\chi)=M\otimes_RR(\chi)$; $\tr_V$ is the trace function. \Cref{conv:reps:conventions}.
\item \textbf{Representations} (definition): a homomorphism $\rho\colon G\to\GL(V)$, $V$ a finite-dimensional $k$-vector space, equivalently
a $k[G]$-module of finite dimension; matrix representations are isomorphic iff conjugate by one $P\in\GL_n(k)$; examples: characters,
permutation and regular representations, the standard representation of $H\le\GL_2(k)$, inflation. \Cref{def:reps:representation}.
\item \textbf{Constructions} (definition): subrepresentations, quotients, $V\oplus W$, $V\otimes W$, $V^\vee$, $\Hom_k(V,W)$ with
$(g\phi)(v)=g\phi(g^{-1}v)$, invariants $V^G$, $\det\rho$. \Cref{def:reps:constructions}.
\item \textbf{Basic facts.} $\tr(g\mid V^\vee)=\tr(g^{-1}\mid V)$, $\tr(g\mid V\otimes W)=\tr(g\mid V)\tr(g\mid W)$, $\det(V\otimes\chi)=\chi^n\det V$;
$V^{\vee\vee}\simeq V$, $(V\otimes W)^\vee\simeq V^\vee\otimes W^\vee$; if $\dim V=2$ (or $V$ free of rank $2$ over $R$), $v\mapsto\det[v\,|\,{-}]$ is a
$G$-isomorphism $V\simeq V^\vee\otimes\det$, so $V^\vee\simeq V\otimes\det^{-1}$. \emph{(proved here)} \Cref{prop:reps:basic-facts}.
\item \textbf{$\Hom_G$ as invariants.} $V^\vee\otimes W\to\Hom(V,W)$, $\lambda\otimes w\mapsto\lambda(\cdot)w$, is a $G$-isomorphism (also for free
modules of finite rank over $R$), and $\Hom_G(V,W)=\Hom(V,W)^G=(V^\vee\otimes W)^G$; for a character $\chi$,
$(\End(V)\otimes\chi)^G=\Hom_G(V,V\otimes\chi)$. \emph{(proved here)} \Cref{prop:reps:homG}.
\item \textbf{Extension of scalars.} For a field extension $K/k$: (i) $\Hom_{K[G]}(V_K,W_K)=K\otimes_k\Hom_{k[G]}(V,W)$ and
$(V_K)^G=K\otimes_kV^G$; (ii) if $V,W$ are irreducible and $V_K\simeq W_K$ then $V\simeq W$; (iii) if $V_K$ is irreducible and $W_K\simeq V_K$
then $W\simeq V$; (iv) (Noether-Deuring) $V_K\simeq W_K$ implies $V\simeq W$ for all $V,W$. \emph{(proved here for (i) to (iii); (iv)
stated; reference)} \Cref{prop:reps:scalars}.
\item \textbf{Irreducible, indecomposable} (definition): $V\ne0$ with no subrepresentations but $0$, $V$; not a direct sum of two nonzero
subrepresentations; for $\dim V=2$, irreducible iff no line is stable. \Cref{def:reps:irreducible}.
\item \textbf{Schur's lemma.} For irreducible $V,W$ over any field: (i) a $G$-map $V\to W$ is $0$ or an isomorphism, and $\End_G(V)$ is a
division algebra; (ii) if $k=\bar k$, $\End_G(V)=k$ and $\dim\Hom_G(V,W)\in\{0,1\}$, with $1$ iff $V\simeq W$; (iii) the same over any $k$
for absolutely irreducible $V,W$; (iv) the centre acts on an absolutely irreducible $V$ by scalars, absolutely irreducible
representations of abelian groups are one-dimensional, and $\End_G(V)$ is a field for $V$ irreducible and $G$ abelian.
\emph{(proved here)} \Cref{thm:reps:schur}.
\item \textbf{Absolutely irreducible} (definition): $V_{\bar k}$ is irreducible. \Cref{def:reps:absirr}.
\item \textbf{Criteria for absolute irreducibility.} (i) $V$ is absolutely irreducible iff $k[G]\to\End_k(V)$ is surjective, and then $V_K$
is irreducible for every extension $K/k$; (ii) absolutely irreducible implies irreducible with $\End_G(V)=k$; (iii) for $\dim V=2$:
iff no line of $\bar k^2$ is stable; (iv) over any field, an irreducible $V$ is absolutely irreducible iff $\End_G(V)=k$.
\emph{(proved here)} \Cref{prop:reps:absirr}.
\item \textbf{Distinct eigenvalues.} If some $g\in G$ acts on $V$ with $\dim V$ distinct eigenvalues in $k$, every $G$-stable subspace of
$V_{\bar k}$ is defined over $k$; so $V$ is irreducible iff it is absolutely irreducible. In particular this holds for $\dim V=2$,
$\operatorname{char}k\ne2$, if some $c\in G$ has $\rho(c)^2=1$, $\det\rho(c)=-1$. \emph{(proved here)} \Cref{prop:reps:distinct-eigenvalues}.
\item \textbf{$C_4$ on $\F_3^2$.} $C_4=\langle J\rangle$ acts irreducibly but not absolutely irreducibly on $\F_3^2$, with $\End_{C_4}(\F_3^2)=\F_3[J]\simeq\F_9$;
over $\F_9$, $\F_9^2=\F_9v\oplus\F_9\bar v$ with $v=(1,i)^t$, $\bar v=(1,-i)^t$, $Jv=iv$, $J\bar v=-i\bar v$: the sum $\psi\oplus\psi^{-1}$, $\psi(J)=i$ of order
$4$. \emph{(proved here)} \Cref{ex:reps:c4}.
\item \textbf{Semisimple} (definition): a sum of irreducible subrepresentations. \Cref{def:reps:semisimple}.
\item \textbf{Semisimple modules.} For finite-dimensional $V$ and any $G$: $V$ is a sum of irreducibles iff a direct sum of irreducibles
iff every subrepresentation has a stable complement; subrepresentations and quotients of semisimple representations are
semisimple; $V=\bigoplus_SV_S$ with $V_S\simeq S^{m_S}$, and $m_S=\dim\Hom_G(S,V)$ if $S$ is absolutely irreducible. \emph{(proved here)}
\Cref{prop:reps:semisimple}.
\item \textbf{Maschke.} (i) If $G$ is finite and $\operatorname{char}k\nmid\#G$, every representation of $G$ over $k$ is semisimple.
(ii) If $\#G\in R^\times$ and $M$ is an $R[G]$-module, $e_G=\frac1{\#G}\sum_gg$ is an idempotent with $e_GM=M^G$, $M=M^G\oplus(1-e_G)M$,
and $(M/IM)^G=M^G/IM^G$ for every ideal $I$: invariants commute with reduction. \emph{(proved here)} \Cref{thm:reps:maschke}.
\item \textbf{Failures of semisimplicity.} $\mathrm U=\langle u\rangle$ on $\F_3^2$ is indecomposable and not semisimple; $\F_p[C_p]$ is
uniserial; $\Z$ acting on $\Ql^2$ through $\sm1{n'}01$, $n'\ne0$, is not semisimple; $\mathrm B(\F_3)$ on $\F_3^2$ has $\End_{\mathrm B}=\F_3$ but is
reducible. \emph{(proved here)} \Cref{ex:reps:maschke-fails}.
\item \textbf{Jordan-H\"older.} A module of finite length over any ring (here $k[G]$ or $\mathcal{O}[G]$) has a composition series, and the
composition factors with multiplicity, and the length, do not depend on the series. \emph{(stated; reference: proved in
\Cref{ch:modules})} \Cref{thm:reps:jordan-holder}.
\item \textbf{Semisimplification} (definition): $V^{ss}=\bigoplus_iM_i/M_{i-1}$ for a composition series $(M_i)$.
\Cref{def:reps:semisimplification}.
\item \textbf{Properties of $V^{ss}$.} (i) $V^{ss}$ is well defined, $V^{ss}\simeq V$ iff $V$ is semisimple, and $(\cdot)^{ss}$ is additive in
short exact sequences; (ii) $\det(X-g\mid V)=\det(X-g\mid V^{ss})$; (iii) over a perfect $k$, $(V^{ss})_K\simeq(V_K)^{ss}$ for algebraic $K/k$.
\emph{(proved here for (i), (ii); (iii) stated; reference)} \Cref{prop:reps:semisimplification}.
\item \textbf{Extensions of two characters.} Let $\dim V=2$ with a stable line $L$ on which $G$ acts by $\eta_1$, and $\eta_2$ on $V/L$,
$\rho=\sm{\eta_1}b0{\eta_2}$. (i) $V$ is semisimple iff $V\simeq\eta_1\oplus\eta_2$ iff $V$ has a second stable line iff $b(g)=c(\eta_2(g)-\eta_1(g))$
for some $c\in k$; (ii) if $V$ is not semisimple, $L$ is its only stable line, and an isomorphism of two such extensions matches
the sub-characters and the quotient characters; (iii) if $\eta_1\ne\eta_2$ and $V$ is semisimple, $V$ has exactly two stable lines.
\emph{(proved here)} \Cref{prop:reps:nonsplit-line}.
\item \textbf{Absolute irreducibility in dimension $2$.} (i) If $\rho(G)$ is abelian, $\rho$ is absolutely reducible; (ii) if $\rho(G)$ is finite of
order prime to $\operatorname{char}k$ and $\rho$ is absolutely reducible, then $\rho(G)$ is abelian; so a finite non-abelian subgroup of $\GL_2(k)$
of order prime to $\operatorname{char}k$ acts absolutely irreducibly; (iii) if $\Res_H\rho$ is absolutely irreducible, so is $\rho$; (iv) $Q_8$,
$\SL_2(\F_3)$, $N_{ns}(\F_3)$, $\GL_2(\F_3)$ and $N_s(\F_3)$ act absolutely irreducibly on $\F_3^2$; $C_4$, $C_{ns}(\F_3)$, $\mathrm B(\F_3)$ do not.
\emph{(proved here)} \Cref{prop:reps:dim2-absirr}.
\item \textbf{Trace function} (definition): $\tr_V(g)=\tr\rho(g)$, a class function (the ``character'' of Riley-Hobson-Bence).
\Cref{def:reps:trace-character}.
\item \textbf{Orthogonality.} $G$ finite, $k=\C$, $\langle\phi,\psi\rangle=\frac1{\#G}\sum_g\phi(g)\overline{\psi(g)}$: (i) for irreducible $V,W$,
$\langle\tr_V,\tr_W\rangle$ is $1$ if $V\simeq W$ and $0$ otherwise; (ii) $\dim\Hom_G(V,W)=\langle\tr_W,\tr_V\rangle$, $V$ is irreducible iff
$\langle\tr_V,\tr_V\rangle=1$, and the trace function determines $V$; (iii) the irreducible trace functions are an orthonormal basis of the
class functions, their number is the number of conjugacy classes, and $\sum_i(\dim V_i)^2=\#G$; (iv) $\dim V_i$ divides $\#G$.
\emph{(proved here for (i) to (iii); (iv) stated; reference)} \Cref{thm:reps:orthogonality}.
\item \textbf{Traces in characteristic $0$.} If $\operatorname{char}k=0$, two semisimple representations of any group with the same trace
function are isomorphic. \emph{(proved here)} \Cref{cor:reps:traces-char0}.
\item \textbf{$D_4$ and $Q_8$.} They have the same character table and are not isomorphic; the two-dimensional representation of $D_4$ is
defined over $\Q$ and reduces modulo $3$ to the standard representation of $N_s(\F_3)$. \emph{(proved here)} \Cref{ex:reps:tables}.
\item \textbf{Traces fail in characteristic $p$.} Over $\F_3$, $\mathbf 1^{\oplus3}$ and $\sgn^{\oplus3}$ of $C_2$ have the same traces and are not
isomorphic; traces suffice in dimension $<p$. \emph{(proved here)} \Cref{ex:reps:traces-fail}.
\item \textbf{Density theorem (Jacobson, Burnside).} For pairwise non-isomorphic irreducible $S_1,\dots,S_r$ with $D_i=\End_G(S_i)$:
(a) the image of $k[G]\to\prod_i\End_k(S_i)$ is $\prod_i\End_{D_i}(S_i)$; (b) if $k=\bar k$, it is $\prod_i\End_k(S_i)$.
\emph{(proved here)} \Cref{thm:reps:burnside}.
\item \textbf{Brauer-Nesbitt.} If $\det(X-\rho_1(g))=\det(X-\rho_2(g))$ in $k[X]$ for all $g\in G$ (any group, any field), then
$V_1^{ss}\simeq V_2^{ss}$. Equal traces suffice if $\operatorname{char}k=0$ or $\dim V_1,\dim V_2<\operatorname{char}k$. \emph{(proved here)}
\Cref{thm:reps:brauer-nesbitt}.
\item \textbf{Consequences.} (i) In dimension $2$, equal traces and equal determinants give $V_1^{ss}\simeq V_2^{ss}$; (ii) if $V_1$ is over
$k$, $V_2$ over $k'\supseteq k$, they have the same characteristic polynomials, and $V_1\otimes_kk'$ is irreducible, then
$V_2\simeq V_1\otimes_kk'$; this fails when $V_1$ is only irreducible over $k$. \emph{(proved here)} \Cref{cor:reps:bn-consequences}.
\item \textbf{Induction} (definition): $\Ind_H^GW=k[G]\otimes_{k[H]}W$, $\Res_H^GV$. \Cref{def:reps:induction}.
\item \textbf{Induced representations.} For $G=\bigsqcup_is_iH$ of finite index $m$: $\Ind W=\bigoplus_is_i\otimes W$, $\dim\Ind W=m\dim W$,
$\tr(g\mid\Ind W)=\sum_{s_i^{-1}gs_i\in H}\tr(s_i^{-1}gs_i\mid W)$; for $[G:H]=2$, $s\notin H$ and a character $\psi$ of $H$:
$h\mapsto\mathrm{diag}(\psi(h),\psi^s(h))$, $s\mapsto\sm0{\psi(s^2)}10$, with $\psi^s(h)=\psi(s^{-1}hs)$. \emph{(proved here)} \Cref{prop:reps:induction}.
\item \textbf{Frobenius reciprocity.} $\Hom_G(\Ind_H^GW,V)\simeq\Hom_H(W,\Res_HV)$, and for finite index $\Hom_G(V,\Ind_H^GW)\simeq\Hom_H(\Res_HV,W)$,
over a field or a commutative ring. \emph{(proved here)} \Cref{thm:reps:frobenius-reciprocity}.
\item \textbf{Projection formula.} $\Ind(W\otimes\Res U)\simeq\Ind W\otimes U$; so $\Ind W\otimes\chi\simeq\Ind W$ if $\chi|_H=1$; $\Ind_H^G\mathbf1=k[G/H]$,
$\simeq\mathbf1\oplus\chi_H$ for index $2$ and $\operatorname{char}k\ne2$; transitivity, additivity, extension of scalars. \emph{(proved here)}
\Cref{prop:reps:projection-formula}.
\item \textbf{Mackey.} (a) $\Res_K\Ind_H^GW\simeq\bigoplus_{s\in K\backslash G/H}\Ind^K_{K\cap sHs^{-1}}W^s$; (b) in characteristic $0$, $G$ finite,
$\Ind W$ is irreducible iff $W$ is irreducible and $W^s$, $W$ have no common constituent on $sHs^{-1}\cap H$ for $s\notin H$; (c) for
$[G:H]=2$, any characteristic, $W$ absolutely irreducible: $\Res_H\Ind W\simeq W\oplus W^s$, and $\Ind W$ is absolutely irreducible iff
$W_{\bar k}\not\simeq W^s_{\bar k}$. \emph{(stated; reference for (a); proof sketched here for (b); proved here for (c))} \Cref{thm:reps:mackey}.
\item \textbf{Clifford.} For $N\trianglelefteq G$ of finite index and $V$ irreducible over any field, $\Res_NV$ is semisimple, its constituents
form one $G$-orbit and have equal multiplicities. \emph{(proved here)} \Cref{thm:reps:clifford}.
\item \textbf{Clifford theory for index $2$.} $[G:H]=2$, $\chi_H$ the character with kernel $H$, $\operatorname{char}k\ne2$, $V$ absolutely
irreducible of dimension $n$. Either (A) $\Res_HV$ is absolutely irreducible and $V\not\simeq V\otimes\chi_H$; or (B) $n=2d$,
$\Res_HV_{\bar k}\simeq W\oplus W^s$ with $W\not\simeq W^s$ absolutely irreducible, $V_{\bar k}\simeq\Ind_H^GW$, and $V\simeq V\otimes\chi_H$ over $k$, with
$\Hom_G(V,V\otimes\chi_H)$ a line spanned by an isomorphism of trace $0$ with eigenspaces $W$, $W^s$. For $n=2$: (B) iff
$V_{\bar k}\simeq\Ind\psi$ with $\psi^s\ne\psi$. \emph{(proved here)} \Cref{thm:reps:clifford-index-two}.
\item \textbf{Dihedral representations.} For $\rho\colon G\to\GL_2(k)$ absolutely irreducible, $\operatorname{char}k\ne2$, with finite projective image
$P$: $P$ is not cyclic; $P$ is dihedral (including $V_4$) iff $\rho\simeq\rho\otimes\chi$ for a character of order $2$ iff $\rho_{\bar k}$ is induced
from an index-two subgroup; the characters $\chi\ne1$ with $\rho\simeq\rho\otimes\chi$ correspond to the cyclic subgroups of index $2$ of $P$
(one if $\#P\ge6$, three if $P\simeq V_4$). \emph{(proved here)} \Cref{prop:reps:dihedral}.
\item \textbf{Twisting.} $\tr(\rho\otimes\chi)=\chi\tr\rho$, $\det(\rho\otimes\chi)=\chi^n\det\rho$; (absolute) irreducibility and $\Hom_G$ are
unchanged; $\End(V\otimes\chi)=\End(V)$ and $\ad(V\otimes\chi)=\ad V$; $V\simeq V\otimes\chi$ forces $\chi^n=1$. \emph{(proved here)}
\Cref{prop:reps:twist}.
\item \textbf{Trace-zero lemma.} If $\phi\colon V\to V\otimes\chi$ is a $G$-map then $(1-\chi(g))\tr\phi=0$ for all $g$; so $\tr\phi=0$ if $\chi\ne1$ (over
a field) or if $1-\chi(g)\in R^\times$ for some $g$. \emph{(proved here)} \Cref{lem:reps:trace-zero}.
\item \textbf{Adjoint} (definition): $\End(V)$ with $g\cdot x=\rho(g)x\rho(g)^{-1}$, and $\ad V=\ker\tr$. \Cref{def:reps:adjoint}.
\item \textbf{The adjoint.} (i) $\End(V)=k\oplus\ad V$ if $n\in k^\times$; (ii) $\Ad\rho(g)$ has eigenvalues $\lambda_i/\lambda_j$, and on $\ad V$ for $n=2$:
$\lambda_1/\lambda_2,1,\lambda_2/\lambda_1$; (iii) in the basis $e,h,f$: $\Ad\mathrm{diag}(x,y)=\mathrm{diag}(x/y,1,y/x)$; for $t=\sm1{n'}01$, $(t-1)e=0$,
$(t-1)h=-2n'e$, $(t-1)f=n'h-n'^2e$; for $F=\sm{\varepsilon q}{c'}0\varepsilon$, $FeF^{-1}=qe$, $FhF^{-1}=h-2c'\varepsilon^{-1}e$,
$FfF^{-1}=q^{-1}f+c'\varepsilon^{-1}q^{-1}h-c'^2\varepsilon^{-2}q^{-1}e$. \emph{(proved here)} \Cref{prop:reps:adjoint}.
\item \textbf{Invariants of twisted adjoints.} (i) $(\End V\otimes\chi)^G=\Hom_G(V,V\otimes\chi)$; (ii) if $n\in R^\times$, $(\ad V\otimes\chi)^G$ is its
trace-zero part; (iii) if $\chi$ (over a local $R$: its reduction) is nontrivial, every such map has trace $0$; (iv) for $V$
absolutely irreducible, $(\ad V)^G=0$, and for $\chi\ne1$, $(\ad V\otimes\chi)^G\ne0$ iff $V\simeq V\otimes\chi$, and then it is a line.
\emph{(proved here)} \Cref{prop:reps:H0-twisted-adjoint}.
\item \textbf{Symmetric square.} For $\dim V=2$, $\Xi(xy)(v)=\omega(v,x)y+\omega(v,y)x$ is a $G$-map $\Sym V\otimes\det^{-1}\to\ad V$ with
$\Xi(e_1^2)=-2e$, $\Xi(e_1e_2)=h$, $\Xi(e_2^2)=2f$, an isomorphism iff $2$ is invertible (index $4$ over $\Z$); so $\ad V\simeq\Sym V\otimes\det^{-1}$;
for $V=V_\ell E$, $\Sym(V^\vee)\simeq\ad V(-1)$ and $L_p(\Sym f,s)=L_p(A_f,s-1)$. \emph{(proved here)} \Cref{prop:reps:sym2}.
\item \textbf{Trace form.} $B(x,y)=\tr(xy)$ on $\End_R(T)$ is $G$-invariant and perfect; on $\ad T$ ($n=2$) its Gram matrix in $e,h,f$ has
determinant $-2$, so for $\ell$ odd $\ad T\simeq(\ad T)^\vee$ and $\Hom_{\Zl}(\ad T,\Zl(\chi))\simeq\ad T(\chi)$; for $\ell=2$ the index is $2$.
\emph{(proved here)} \Cref{prop:reps:trace-form}.
\item \textbf{The adjoint ignores twists, duals and homotheties.} $\ad(V\otimes\chi)=\ad V$; $x\mapsto-x^t$ is a $G$-isomorphism
$\ad V\simeq\ad(V^\vee)$; for lattices, $\ad(cL)=\ad L$ and $\ad(L^\vee)$ corresponds to $\ad L$. \emph{(proved here)}
\Cref{prop:reps:adjoint-invariance}.
\item \textbf{The subgroups of $\GL_2(\F_3)$ on $\F_3^2$.} The table of irreducibility, absolute irreducibility, $\dim\End_H$, $\dim(\ad\rho)^H$,
$\dim\Hom_H(\rho,\rho\otimes\det)$, the twist $\rho\simeq\rho\otimes\det$, and projective images, for $\GL_2,\SL_2,N_{ns},Q_8,N_s,C_{ns},C_4,C_s,\mathrm B,\mathrm U$.
\emph{(proved here; checked by computer)} \Cref{prop:reps:gl23-table}.
\item \textbf{$N_s(\F_3)$ is induced over $\F_9$.} The standard representation $\rho$ of $N_s(\F_3)$ is absolutely irreducible; in the basis
$v=(1,i)^t$, $\sigma v=(1,-i)^t$ of $\F_9^2$, $\rho(J)=\mathrm{diag}(i,-i)$ and $\rho(\sigma)=w$, so $\rho\otimes\F_9\simeq\Ind_{C_4}^{N_s}\psi$, $\psi(J)=i$; over $\F_3$, $\rho$
is not induced from $C_4$; $\rho\simeq\rho\otimes\chi$ for the three characters with kernels $C_4$, $C_s$, $\{\pm1,\pm w\}$, with intertwiners
$J$, $\sigma$, $w$; $(\ad\rho\otimes\det)^{N_s}=\F_3J$. \emph{(proved here)} \Cref{prop:reps:ns3-induced}.
\item \textbf{$N_{ns}(\F_3)$.} Its standard representation and its restriction to $Q_8$ are absolutely irreducible; $\rho\not\simeq\rho\otimes\det$ and
$(\ad\rho\otimes\det)^{N_{ns}}=0$; $\rho\otimes\F_9\simeq\Ind_{C_{ns}}^{N_{ns}}\theta$, $\theta(a)=1+i$, and $\rho\simeq\rho\otimes\chi$ only for the character with
kernel $C_{ns}(\F_3)$. \emph{(proved here)} \Cref{prop:reps:nns3}.
\item \textbf{$\SL_2(\F_3)$ and $\GL_2(\F_3)$.} Their standard representations are absolutely irreducible; for $\GL_2(\F_3)$, $\rho\not\simeq\rho\otimes\det$
and $(\ad\rho\otimes\det)^{\GL_2(\F_3)}=0$; projective images $A_4$, $S_4$. \emph{(proved here)} \Cref{prop:reps:gl23-sl23}.
\item \textbf{The three images.} For $H$ conjugate to $\GL_2(\F_3)$, $N_{ns}(\F_3)$ or $N_s(\F_3)$ and $H_1=H\cap\SL_2(\F_3)$, equivalent: (a) $H$ is
conjugate to $N_s(\F_3)$; (b) $\Res_{H_1}\rho$ is not absolutely irreducible; (c) $\rho\simeq\rho\otimes\det$; (d) $(\ad\rho\otimes\det)^H\ne0$;
(e) $\rho\otimes\F_9\simeq\Ind_{H_1}^H\psi$ with $\psi$ of order $4$; (f) $H_1$ is abelian. Then $(\ad\rho\otimes\det)^H$ is a line spanned by a
trace-zero isomorphism; otherwise it is $0$. \emph{(proved here)} \Cref{thm:reps:three-images}.
\item \textbf{The letter $\sigma$.} In \Cref{sec:reps:cyclic} $\sigma$ is the paper's element of order $3$ and $\mathrm{diag}(1,-1)$ is $h$.
\Cref{conv:reps:sigma}.
\item \textbf{Unipotent elements.} For $u\ne1$ unipotent in $\GL_2(k)$ and $\xi\in k^\times\setminus\{1\}$: $V^u$ is a line, $V^{\xi u}=0$; if
$\operatorname{char}k\ne2$, $(\ad V)^u=(\ad V)^{\xi u}=k(u-1)$, a line; if $FuF^{-1}=u^q$ then $F$ acts on it by $q$ (and $F^{-1}$ by $q^{-1}$).
\emph{(proved here)} \Cref{prop:reps:unipotent-invariants}.
\item \textbf{Commutant lemma.} If $x\in\End_k(V)$ has $\dim V$ distinct eigenvalues in $\bar k$, or $\dim V=2$ and $x$ is not scalar, its
commutant is $k[x]\simeq k[X]/(\det(X-x))$; for $\dim V=2$ and $\operatorname{char}k\ne2$ its trace-zero part is the line $k(2x-\tr x)$; it is a
field iff $\det(X-x)$ is irreducible. \emph{(proved here)} \Cref{lem:reps:commutant}.
\item \textbf{Cyclic groups of order $2,3,4,6$.} An element $g\in\SL_2(K)$ of order $m\in\{3,4,6\}$, $\operatorname{char}K\nmid m$, has characteristic
polynomial $\Phi_m$; order $2$ means $g=-1$. Fixed spaces on $V$ and $\ad V$ (table), the invariant line of $\ad V$ spanned by $g-g^{-1}$,
and an $F$ with $FgF^{-1}=g^r$ acts on it by $\pm1$ according as $r\equiv\pm1\pmod m$. \emph{(proved here)} \Cref{prop:reps:cyclic-shapes}.
\item \textbf{Reduction of finite cyclic subgroups of $\SL_2(\Zl)$.} For $g$ of order $m\in\{2,3,4,6\}$: if $\ell\nmid m$, invariants on $T$ and
$\ad T$ commute with reduction; if $\ell=3$, $m=3$, then $\bar g$ is a nontrivial unipotent, $(T/3T)^g$ is a line while $T^g=0$, and
$(\ad T)^g=\Z_3(2g+1)$ reduces onto $(\ad T/3)^g$; for $m=6,4,2$ and $\ell=3$, $(T/3T)^g=0$; the integral commutant is $\Zl[g]\simeq\Zl[X]/(\Phi_m)$.
\emph{(proved here; the freeness of $T$ over $\Zl[g]$ stated; reference)} \Cref{prop:reps:reduction-cyclic}.
\item \textbf{Lattices} (definition): finitely generated $\mathcal{O}$-submodules spanning $V$; homothety $L'=cL$. \Cref{def:reps:lattice}.
\item \textbf{Stable lattices.} Any two lattices satisfy $\varpi^aL\subseteq L'\subseteq\varpi^{-b}L$; for finite $G$, $\sum_g\rho(g)L_0$ is stable; $\varpi^mL$,
$L^\vee$, $L\otimes\mathcal{O}(\chi)$, $\End_\mathcal{O}(L)$, $\ad L$ are stable, and $\ad L/\varpi^m=\ad(L/\varpi^mL)$. \emph{(proved here)} \Cref{prop:reps:lattice}.
\item \textbf{Reduction.} $\det(X-\rhobar_L(g))$ is the reduction of $\det(X-\rho(g))$; $\rhobar_L^{ss}$ does not depend on $L$; if $\rhobar_L$ is
irreducible for one $L$, $\rhobar_{L'}\simeq\rhobar_L$ for all $L'$. \emph{(proved here)} \Cref{prop:reps:reduction-ss}.
\item \textbf{Homothety lemma.} If $\rhobar_L$ is irreducible, every $G$-stable lattice is $\varpi^mL$ for a unique $m\in\Z$. \emph{(proved here)}
\Cref{thm:reps:homothety}.
\item \textbf{Non-homothetic lattices.} For the upper triangular group, $\Zl e_1\oplus\Zl e_2$ and $\Zl e_1\oplus\ell\Zl e_2$ are stable, not homothetic,
with non-isomorphic reductions of equal semisimplification. \emph{(proved here)} \Cref{ex:reps:nonhomothetic}.
\item \textbf{Ribet's lemma.} Let $K$ be finite over $\Ql$ with ring of integers $\mathcal{O}$, $G$ compact, and $\rho\colon G\to\GL_2(K)$ continuous and
irreducible. If $\rhobar^{ss}\simeq\eta_1\oplus\eta_2$, some stable lattice reduces to a
non-split extension with sub-character $\eta_1$ (and one with sub-character $\eta_2$). \emph{(stated; reference)} \Cref{thm:reps:ribet}.
\item \textbf{Twisted self-isomorphisms modulo $\varpi^m$.} If $T/\varpi T$ is absolutely irreducible, the module of $G$-maps
$T/\varpi^m\to(T/\varpi^m)(\chi)$ is cyclic, its elements of exact order $\varpi^m$ are isomorphisms, and the length of the invariants of $X=\End_\mathcal{O}(T)(\chi)\otimes K/\mathcal{O}$ is the
largest $h$ with $T/\varpi^h\simeq(T/\varpi^h)(\chi)$; the same for $\ad T(\chi)$ if $n\in\mathcal{O}^\times$, $\bar\chi\ne1$; $\chi^n\equiv1\pmod{\varpi^h}$; for $n=2$, $\ell$ odd,
$\chi(G)=\Zl^\times$: $h=0$ if $\ell\ge5$, $h\le1$ if $\ell=3$. \emph{(proved here)} \Cref{prop:reps:twist-length}.
\item \textbf{The mod $\ell$ dictionary.} For $\rhobar\colon G_\Q\to\GL_2(\F_\ell)$ continuous, $\ell$ odd, $\det\rhobar=\bar\chi_\ell$: irreducible equals
absolutely irreducible; for $\rhobar$ irreducible, the hypothesis of Theorem thm:main fails iff $\rhobar\simeq\rhobar\otimes\chi_{\ell^*}$, iff $\rhobar\otimes\Fbar_\ell$
is induced from $G_{\Q(\sqrt{\ell^*})}$ (for $\rhobar$ reducible it always fails); Lemma lem:h0 concerns $\rhobar\simeq\rhobar\otimes\bar\chi_\ell$, possible only
for $\ell=3$. \emph{(proved here from the cited Galois input)} \Cref{rem:reps:mod-l-galois}.
\item \textbf{Dickson.} For $\ell$ odd and $G\le\GL_2(\F_\ell)$: if $\ell\mid\#G$, $G\supseteq\SL_2(\F_\ell)$ or $G$ lies in a Borel subgroup; if $\ell\nmid\#G$, the
image in $\PGL_2(\F_\ell)$ is cyclic, dihedral, or $A_4$, $S_4$, $A_5$. \emph{(stated; reference)} \Cref{thm:reps:dickson}.
\end{enumerate}

\section*{Where the paper uses this chapter}

\begin{itemize}
\item \textbf{Proposition prop:ns} (the image $N_s(\F_3)$): the three images and their determinant-one subgroups, \Cref{thm:reps:three-images};
absolute irreducibility of $Q_8$ and $\SL_2(\F_3)$, \Cref{prop:reps:dim2-absirr}; $\psi$ of order $4$ and $\rhobar\otimes\F_9\simeq\Ind\psi$,
\Cref{ex:reps:c4}, \Cref{prop:reps:ns3-induced}, \Cref{thm:reps:frobenius-reciprocity}; $\rhobar\simeq\rhobar\otimes\chi_{-3}$, \Cref{prop:reps:projection-formula},
\Cref{thm:reps:clifford-index-two}, descended to $\F_3$ by \Cref{prop:reps:scalars}; $H^0(\Q,\ad\rhobar(1))\ne0$, \Cref{prop:reps:H0-twisted-adjoint}.
\item \textbf{Lemma lem:h0} (the term $h_\ell(E)$): $H^0$ of a twisted adjoint as trace-zero $G$-maps, \Cref{prop:reps:homG},
\Cref{prop:reps:H0-twisted-adjoint}; the trace-zero lemma, \Cref{lem:reps:trace-zero}; Schur, \Cref{thm:reps:schur}; cyclicity, isomorphisms modulo
$\ell^k$ and the determinant bound, \Cref{prop:reps:twist-length}; Clifford theory for $\Q(\sqrt{-3})$, \Cref{thm:reps:clifford-index-two}.
\item \textbf{Lemma lem:IV, Remark rem:types, Lemma lem:euler, Proposition prop:ns-structure}: unipotent and minus-unipotent elements,
\Cref{prop:reps:unipotent-invariants}; the commutant and its trace-zero line, \Cref{lem:reps:commutant}; elements of order $2,3,4,6$ and the
action of Frobenius on the invariant line, \Cref{prop:reps:cyclic-shapes}; reduction modulo $3$, \Cref{prop:reps:reduction-cyclic};
Maschke over $\Zl$, \Cref{thm:reps:maschke}.
\item \textbf{Lemma lem:old} (oldforms): Brauer-Nesbitt over a common finite field, \Cref{thm:reps:brauer-nesbitt}, \Cref{cor:reps:bn-consequences};
absolute irreducibility from complex conjugation, \Cref{prop:reps:distinct-eigenvalues}; the unique stable line of a non-split extension,
\Cref{prop:reps:nonsplit-line}.
\item \textbf{Section 4.2} (the theorem of Diamond, Flach and Guo) \textbf{and Lemma lem:fpr}: $A=\ad T$, $T^D=\Hom(A,\Zl(1))\simeq A(1)$ by the trace
form, \Cref{prop:reps:trace-form}; the adjoint of $M_{f,\lambda}\simeq T_\ell E(-1)$ is $\ad T$, \Cref{prop:reps:adjoint-invariance}; the homothety lemma
as an alternative route, \Cref{thm:reps:homothety}.
\item \textbf{Lemma lem:tamfactor} (the Tamagawa factor at a multiplicative prime): the basis $e,h,f$ and the action of $t$ and of the Frobenius
matrix, \Cref{prop:reps:adjoint}; the invariants of a unipotent, \Cref{prop:reps:unipotent-invariants}; the shape $\sm{\chi_\ell\varepsilon}{\varepsilon\kappa}0\varepsilon$,
\Cref{prop:reps:nonsplit-line}.
\item \textbf{Section 1.1 and the Introduction}: $A_f=\ad V$, the eigenvalues $\alpha/\beta,1,\beta/\alpha$, \Cref{prop:reps:adjoint}; $\ad\simeq\Sym\otimes\det^{-1}$
and $L_p(\Sym f,s)=L_p(A_f,s-1)$, \Cref{prop:reps:sym2}.
\item \textbf{Proofs of Theorems thm:main and thm:quant, Lemma lem:euler (types $\kod{I}_n^*$, $\kod{I}_0^*$)}: $\ad T$ is unchanged by twists,
\Cref{prop:reps:twist}, \Cref{prop:reps:adjoint-invariance}.
\item \textbf{The hypothesis of Theorem thm:main, Proposition prop:Sf, DDT Theorem 3.15}: absolute irreducibility on $G_{\Q(\sqrt{\ell^*})}$,
\Cref{def:reps:absirr}, \Cref{thm:reps:clifford-index-two}, \Cref{prop:reps:dihedral}, \Cref{rem:reps:mod-l-galois}.
\item \textbf{Conjecture conj:euler, Remark rem:ns, Section sec:anomaly} (the image $\Ns$ and $h_3(E)=1$): \Cref{prop:reps:ns3-induced},
\Cref{thm:reps:three-images}, \Cref{prop:reps:twist-length}.
\end{itemize}

\section{Representations over a field}\label{sec:reps:field}

This section fixes the language: representations, the linear constructions on them, and the behaviour of spaces of
$G$-maps under extension of the field of scalars. In the paper these are used silently everywhere. Three uses are explicit:
the identification of $H^0$ of a twisted adjoint with a space of $G$-maps in Lemma lem:h0 (\Cref{prop:reps:homG}), the
descent from $\F_9$ to $\F_3$ of the isomorphism $\rhobar\simeq\rhobar\otimes\chi_{-3}$ in Proposition prop:ns, and the phrase ``over a
common finite extension of $\F_\ell$'' in Lemma lem:old (\Cref{prop:reps:scalars}). The words ``irreducible'' and ``absolutely irreducible'' and
Schur's lemma, defined and proved in \Cref{sec:reps:irreducible}, are used here by forward reference; there is no circularity, since
\Cref{prop:reps:scalars}(ii) uses only \Cref{thm:reps:schur}(i), whose proof uses nothing from this section.

\begin{convention}\label{conv:reps:conventions}
The following conventions hold in this chapter.
\begin{enumerate}[label=(C\arabic*)]
\item \emph{Representations.} $G$ is a group (finite only where this is said), $k$ a field and $R$ a commutative ring. A
representation of $G$ over $k$ is a homomorphism $\rho\colon G\to\GL(V)$ with $V$ a \emph{finite-dimensional} $k$-vector space;
over $R$, $V$ is a free $R$-module of finite rank. Matrices act on column vectors on the left (as in \Cref{ch:groups}). A
\emph{$G$-map} (intertwiner) is a $k[G]$-linear map; $\Hom_G(V,W)$ is the space of $G$-maps and $V\simeq W$ means isomorphic as
representations. We often write $gv$ for $\rho(g)v$.
\item \emph{Duals.} The linear dual is $V^\vee=\Hom_k(V,k)$ with $(g\lambda)(v)=\lambda(g^{-1}v)$; in the dual basis $(e_i^\vee)$, $g$ acts by the
matrix $\rho(g^{-1})^t$. The symbol ${}^*$ is reserved for Pontryagin duals (convention B-C4 of \Cref{ch:notation}).
\item \emph{Adjoints.} $\End(V)=\End_k(V)$ is a representation by $g\cdot x=\rho(g)x\rho(g)^{-1}$; written $\Ad V$ when the action is to be
stressed. $\ad V$ is its subspace of trace zero (the paper's notation; the superscript $0$ is part of the symbol). $\Sym V$ is the
symmetric square. The matrix units are $E_{ij}=e_i\otimes e_j^\vee$, and for $\dim V=2$ the basis of $\ad V$ is the paper's
\begin{gather*}
e=E_{12}=e_1\otimes e_2^\vee,\qquad h=E_{11}-E_{22},\qquad f=E_{21}=e_2\otimes e_1^\vee,\\
e=\Mtwo0100,\qquad h=\Mtwo100{-1},\qquad f=\Mtwo0010.
\end{gather*}
\item \emph{Characters and twists.} A character of $G$ with values in $R$ is a homomorphism $\chi\colon G\to R^\times$. $R(\chi)$ is $R$ with
$G$ acting through $\chi$, and $M(\chi)=M\otimes_RR(\chi)$, also written $M\otimes\chi$ or $\rho\otimes\chi$. The trace function of $V$ is written
$\tr_V$, never $\chi_V$, so that the letter $\chi$ is free for characters such as $\chi_{-3}$ and the cyclotomic $\chi_\ell$.
\item \emph{Fields.} $\bar k$ is an algebraic closure of $k$, and $V_K=K\otimes_kV$ for an extension $K/k$. $\F_9=\F_3[i]$ with $i^2=-1$
(\Cref{ch:rings}).
\item \emph{Standard models.} $J=\sm01{-1}0$, $\sigma=\mathrm{diag}(1,-1)$, $w=\sm0110$, $u=\sm1101$, $a=\sm1{-1}11$, and the subgroups $C_s$, $N_s$,
$C_{ns}$, $N_{ns}$, $Q_8$, $C_4=\langle J\rangle$, $\mathrm B$, $\mathrm U=\langle u\rangle$ of $\GL_2(\F_3)$ are the standard models of \Cref{ch:groups}. Thus
$J=m_{-i}$ is multiplication by $-i$ on $\F_9=\F_3\oplus\F_3i$, $a=m_{1+i}$, $a^2=J$, $w=\sigma J$, $N_s(\F_3)=\langle J,\sigma\rangle$, $N_{ns}(\F_3)=\langle a,\sigma\rangle$, and
$\psi\colon C_4\to\F_9^\times$ is the character $\psi(J)=i$.
\item \emph{Lattices.} $\mathcal{O}$ is a discrete valuation ring with uniformiser $\varpi$, fraction field $K$ and residue field $\mathcal{O}/\varpi$; in the paper
$\mathcal{O}=\Zl$ and $\varpi=\ell$. A lattice in a $K$-vector space of dimension $n$ is a finitely generated $\mathcal{O}$-submodule that spans it; it is free of
rank $n$ (\Cref{ch:modules}).
\item \emph{No Galois groups.} This chapter never uses Frobenius elements. Where the paper's Frobenius appears it is an abstract element
$F$ with $FgF^{-1}=g^q$ for $g$ in an inertia image, which is the relation satisfied by the \emph{arithmetic} Frobenius (convention B-C2);
the geometric Frobenius $F^{-1}$ has the inverse eigenvalues. The letter $q$ is then the paper's bad prime.
\item \emph{Typography.} Double names take plain hyphens: Brauer-Nesbitt, Jordan-H\"older.
\end{enumerate}
\end{convention}

\begin{definition}[Representations]\label{def:reps:representation}
A \emph{representation} of $G$ over $k$ is a homomorphism $\rho\colon G\to\GL(V)$, $V$ a finite-dimensional $k$-vector space; equivalently, a
left $k[G]$-module structure on $V$, where $k[G]$ is the group algebra (finite sums $\sum_ga_gg$). Its \emph{degree} is $n=\dim V$; it is
\emph{faithful} if $\ker\rho=1$. A basis of $V$ turns $\rho$ into a homomorphism $G\to\GL_n(k)$, and two such matrix representations $\rho_1$,
$\rho_2$ are isomorphic iff there is one $P\in\GL_n(k)$ with $\rho_2(g)=P\rho_1(g)P^{-1}$ for all $g$ (the matrix of an isomorphism).
Examples:
\begin{enumerate}[label=(\alph*)]
\item the trivial representation $\mathbf 1$ ($V=k$, $g\mapsto1$);
\item one-dimensional representations, which are the characters $\chi\colon G\to k^\times$; they factor through the abelianisation $G^{\mathrm{ab}}$;
\item the permutation representation $k[X]$ of a $G$-set $X$ (basis $X$, $g$ permuting it), and the regular representation $k[G]$;
\item the \emph{standard representation} of a subgroup $H\le\GL_2(k)$ on $k^2$ (the inclusion);
\item inflation: if $\pi\colon G\to G'$ is a homomorphism and $\rho'$ a representation of $G'$, then $\rho'\circ\pi$ is one of $G$. A representation
of a finite quotient of a profinite group is thus one of the group (\Cref{ch:galois}).
\end{enumerate}
The representations of Riley-Hobson-Bence \cite{RHB06} are the case $k=\C$, $G$ finite.
\end{definition}

\begin{definition}[Constructions]\label{def:reps:constructions}
Let $V,W$ be representations of $G$ over $k$.
A \emph{subrepresentation} is a $G$-stable subspace $U\subseteq V$, and $V/U$ is the quotient representation. $G$ acts on $V\oplus W$
componentwise and on $V\otimes W$ by $g(v\otimes w)=gv\otimes gw$; on $V^\vee$ as in (C2); on $\Hom_k(V,W)$ by $(g\phi)(v)=g\,\phi(g^{-1}v)$.
The \emph{invariants} are $V^G=\{v:gv=v\text{ for all }g\}$, and $\det\rho\colon G\to k^\times$ is the character $g\mapsto\det\rho(g)$, the
representation on $\wedge^nV$. The same definitions apply to free modules of finite rank over $R$.
\end{definition}

\begin{proposition}\label{prop:reps:basic-facts}
Let $V,W$ be representations of $G$ over $k$, $n=\dim V$, $\chi$ a character.
\begin{enumerate}[label=(\roman*)]
\item $\tr(g\mid V\oplus W)=\tr(g\mid V)+\tr(g\mid W)$, $\tr(g\mid V\otimes W)=\tr(g\mid V)\tr(g\mid W)$, $\tr(g\mid V^\vee)=\tr(g^{-1}\mid V)$ and
$\det(V\otimes\chi)=\chi^n\det V$.
\item The canonical maps $V\to V^{\vee\vee}$ and $V^\vee\otimes W^\vee\to(V\otimes W)^\vee$, $\lambda\otimes\mu\mapsto(v\otimes w\mapsto\lambda(v)\mu(w))$, are
$G$-isomorphisms.
\item If $n=2$ (or $V$ is free of rank $2$ over $R$), let $\omega(x,y)=\det[x\,|\,y]$. Then $\omega(gx,gy)=\det(g)\omega(x,y)$, and
$v\mapsto\omega(v,{-})$ is a $G$-isomorphism $V\to V^\vee\otimes\det$. Hence $V^\vee\simeq V\otimes\det^{-1}$.
\end{enumerate}
\end{proposition}

\begin{proof}
(i) The first is clear in a basis adapted to $V\oplus W$; the second holds because the matrix of $g$ on $V\otimes W$ in the basis
$e_i\otimes f_j$ is the Kronecker product of the two matrices, whose trace is the product of the traces; the third because $g$ acts on
$V^\vee$ by $\rho(g^{-1})^t$; the fourth because $\det(\chi(g)\rho(g))=\chi(g)^n\det\rho(g)$.
(ii) Both maps are isomorphisms of vector spaces (compare bases). For the first, $v\mapsto(\lambda\mapsto\lambda(v))$, and
$(g\cdot\mathrm{ev}_v)(\lambda)=\mathrm{ev}_v(g^{-1}\lambda)=(g^{-1}\lambda)(v)=\lambda(gv)$, so $g\cdot\mathrm{ev}_v=\mathrm{ev}_{gv}$. For the second,
$(g\lambda\otimes g\mu)(v\otimes w)=\lambda(g^{-1}v)\mu(g^{-1}w)$, which is the value of $g$ applied to the form $\lambda\otimes\mu$ on $V\otimes W$.
(iii) $\omega(gx,gy)=\det(\rho(g)[x\,|\,y])=\det(g)\omega(x,y)$. Put $\Psi(v)=\omega(v,{-})$. As $\omega$ is a perfect pairing ($\omega(e_1,e_2)=1$),
$\Psi$ is an isomorphism of vector spaces (or of free $R$-modules). For $y\in V$,
\[
(g\cdot\Psi(v))(y)=\Psi(v)(g^{-1}y)=\omega(v,g^{-1}y)=\det(g)^{-1}\omega(gv,y),
\]
so $\Psi(gv)=\det(g)\,(g\cdot\Psi(v))$, which says that $\Psi$ is $G$-equivariant from $V$ to $V^\vee\otimes\det$.
\end{proof}

\begin{proposition}[$G$-maps as invariants]\label{prop:reps:homG}
For representations $V,W$ over $k$ (or free $R$-modules of finite rank with a $G$-action), the map
\[
\Theta\colon V^\vee\otimes W\to\Hom(V,W),\qquad \lambda\otimes w\mapsto\bigl(v\mapsto\lambda(v)w\bigr)
\]
is a $G$-isomorphism, and $\Hom_G(V,W)=\Hom(V,W)^G=(V^\vee\otimes W)^G$. In particular $\End(V)^G=\End_G(V)$, and for a character $\chi$,
$(\End(V)\otimes\chi)^G=\Hom_G(V,V\otimes\chi)$, the space of $\phi\in\End(V)$ with $\phi\rho(g)=\chi(g)\rho(g)\phi$ for all $g$.
\end{proposition}

\begin{proof}
$\Theta$ sends $e_i^\vee\otimes f_j$ to the matrix unit mapping $e_i$ to $f_j$, so it is bijective. It is equivariant:
$\Theta(g\lambda\otimes gw)(v)=\lambda(g^{-1}v)\,gw=g\bigl(\Theta(\lambda\otimes w)(g^{-1}v)\bigr)=(g\cdot\Theta(\lambda\otimes w))(v)$. A linear map $\phi$ is
fixed by $G$ iff $g\phi g^{-1}=\phi$ for all $g$, that is, iff $\phi$ is a $G$-map. For the last statement, $\Hom(V,V\otimes\chi)$ is $\End(V)$ with
$g$ acting by $\phi\mapsto\chi(g)\rho(g)\phi\rho(g)^{-1}$, that is, $\End(V)\otimes\chi$; it is fixed iff $\phi\rho(g)=\chi(g)\rho(g)\phi$.
\end{proof}

\begin{paperlink}
In the proof of Lemma lem:h0, the sentence ``the $\ell^k$-torsion of $H^0(\Q,W^D)$ is $H^0(\Q,\ad(T/\ell^k)(1))$, the group of
trace-zero $G_\Q$-maps $T/\ell^k\to(T/\ell^k)(1)$'' is \Cref{prop:reps:homG} for the free $\Z/\ell^k$-module $T/\ell^k$ and $\chi=\chi_\ell\bmod\ell^k$, together with
\Cref{prop:reps:lattice}(iii) for $W^D[\ell^k]=\ad(T/\ell^k)(1)$. In Lemma lem:tamfactor the basis $e=e_1\otimes e_2^\vee$ of $\ad T$ is the image
under $\Theta$ of a vector of $V^\vee\otimes V$, as in (C3).
\end{paperlink}

\begin{proposition}[Extension of scalars]\label{prop:reps:scalars}
Let $K/k$ be a field extension and $V,W$ representations of $G$ over $k$.
\begin{enumerate}[label=(\roman*)]
\item $\Hom_{K[G]}(V_K,W_K)=K\otimes_k\Hom_{k[G]}(V,W)$ inside $\Hom_K(V_K,W_K)=K\otimes_k\Hom_k(V,W)$. In particular the two spaces of $G$-maps
have the same dimension, and $(V_K)^G=K\otimes_kV^G$.
\item If $V$ and $W$ are irreducible (\Cref{def:reps:irreducible}) and $V_K\simeq W_K$, then $V\simeq W$.
\item If $V_K$ is irreducible (for instance if $V$ is absolutely irreducible, \Cref{prop:reps:absirr}) and $W_K\simeq V_K$, then $W\simeq V$.
\item (Noether-Deuring) If $V_K\simeq W_K$, then $V\simeq W$, for arbitrary $V,W$.
\end{enumerate}
\end{proposition}

\begin{proof}
(i) For $g\in G$ let $\delta_g\colon\Hom_k(V,W)\to\Hom_k(V,W)$ be $\phi\mapsto\rho_W(g)\phi-\phi\rho_V(g)$; then $\Hom_{k[G]}(V,W)=\bigcap_g\ker\delta_g$. As
$\Hom_k(V,W)$ is finite-dimensional, the descending chain of finite intersections stabilises, so there are $g_1,\dots,g_r$ with
$\bigcap_g\ker\delta_g=\ker\delta$ for the single linear map $\delta=(\delta_{g_1},\dots,\delta_{g_r})$. A linear map defined over $k$ has the same rank
over $K$ (rank is the size of a nonvanishing minor of its matrix), so $\ker(\delta\otimes K)=K\otimes\ker\delta$. Hence
\[
\textstyle\bigcap_g\ker(\delta_g\otimes K)\subseteq\ker(\delta\otimes K)=K\otimes\bigcap_g\ker\delta_g\subseteq\bigcap_g\ker(\delta_g\otimes K),
\]
and the left side is $\Hom_{K[G]}(V_K,W_K)$. The statement on invariants is the case $V=\mathbf1$.
(ii) By (i), $\Hom_{k[G]}(V,W)\ne0$; a nonzero $G$-map between irreducible representations is an isomorphism (\Cref{thm:reps:schur}(i)).
(iii) If $U\subseteq W$ is a subrepresentation, $U_K$ is a subrepresentation of $W_K\simeq V_K$, so $U_K\in\{0,W_K\}$ and $U\in\{0,W\}$; thus $W$
is irreducible (it is nonzero). The same argument applied to $V$ shows that $V$ is irreducible, since $V_K$ is. So (ii) applies.
(iv) We do not prove this; see \cite[Section 29]{CurtisReiner}. It is not needed in the paper: wherever an isomorphism is descended,
(ii) or (iii) applies.
\end{proof}

\begin{paperlink}
Proposition prop:ns: $\rhobar_{E,3}\otimes\F_9\simeq\Ind\psi\simeq(\rhobar_{E,3}\otimes\chi_{-3})\otimes\F_9$, and both $\rhobar_{E,3}$ and $\rhobar_{E,3}\otimes\chi_{-3}$ are irreducible
over $\F_3$; so (ii) gives $\rhobar_{E,3}\simeq\rhobar_{E,3}\otimes\chi_{-3}$ over $\F_3$ (``the isomorphism descending to $\F_3$''; the same sentence ends the
proof of Lemma lem:h0, where it is (i)). Lemma lem:old: the residual representations of $f$ and of $h$ are compared over a common finite
extension $k'$ of $\F_\ell$; by (iii), an isomorphism there between $\rhobar\otimes k'$ and $\rhobar_h$ is all that is needed.
\end{paperlink}

\section{Irreducibility, Schur's lemma, absolute irreducibility}\label{sec:reps:irreducible}

This section introduces irreducible and absolutely irreducible representations and proves Schur's lemma, the tool behind
every ``a nonzero such map is an isomorphism'' of the paper. Absolute irreducibility is the hypothesis of Theorem thm:main (on
$\Gal(\Qbar/\Q(\sqrt{\ell^*}))$) and of DDT Theorem 3.15; where the paper says ``irreducible'' and then uses Schur in the form $\End=k$
(Lemmas lem:h0 and lem:old), absolute irreducibility is what is needed, and for $\rhobar_{E,\ell}$ with $\ell$ odd it is automatic by the
distinct-eigenvalue criterion (\Cref{prop:reps:distinct-eigenvalues}) applied to complex conjugation.

\begin{definition}\label{def:reps:irreducible}
A representation $V\ne0$ is \emph{irreducible} (or simple) if its only subrepresentations are $0$ and $V$, and \emph{reducible}
otherwise; it is \emph{indecomposable} if it is not the direct sum of two nonzero subrepresentations. Irreducible implies
indecomposable, not conversely (\Cref{ex:reps:maschke-fails}). If $\dim V=2$, $V$ is irreducible iff no line of $V$ is $G$-stable, iff $\rho(G)$
lies in no Borel subgroup of $\GL(V)$ (the stabiliser of a line, \Cref{ch:groups}).
Examples: every one-dimensional representation is irreducible. The standard representations of $\GL_2(\F_3)$ and of $N_s(\F_3)$ on $\F_3^2$
are irreducible (they stabilise no line, \Cref{ch:groups}), and so is that of $C_4=\langle J\rangle$: a stable line would be spanned by an
eigenvector of $J$, but the characteristic polynomial $X^2+1$ of $J$ has no root in $\F_3$. The standard representations of
$\mathrm U=\langle u\rangle$ and of $\mathrm B(\F_3)$ are reducible: they stabilise $\F_3e_1$.
\end{definition}

\begin{theorem}[Schur's lemma]\label{thm:reps:schur}
Let $V,W$ be irreducible representations of a group $G$ over a field $k$.
\begin{enumerate}[label=(\roman*)]
\item Every $G$-map $V\to W$ is $0$ or an isomorphism. Hence $D=\End_G(V)$ is a division algebra, of finite dimension over $k$.
\item If $k$ is algebraically closed, then $\End_G(V)=k$ (the scalars), and $\dim_k\Hom_G(V,W)$ is $1$ if $V\simeq W$ and $0$ otherwise.
\item The conclusions of (ii) hold over any $k$ if $V$ and $W$ are absolutely irreducible (\Cref{def:reps:absirr}).
\item If $V$ is absolutely irreducible, every central element of $G$ acts on $V$ by a scalar; every absolutely irreducible
representation of an abelian group is one-dimensional; and if $G$ is abelian and $V$ irreducible, $\End_G(V)$ is a field, equal
to the image of $k[G]$ in $\End_k(V)$.
\end{enumerate}
\end{theorem}

\begin{proof}
(i) If $\phi\colon V\to W$ is a $G$-map, $\ker\phi$ and $\im\phi$ are subrepresentations. If $\phi\ne0$, then $\ker\phi\ne V$, so $\ker\phi=0$, and
$\im\phi\ne0$, so $\im\phi=W$. A nonzero element of $\End_G(V)$ is therefore invertible, and its inverse is again a $G$-map.
(ii) Let $\phi\in\End_G(V)$. As $V\ne0$ is finite-dimensional and $k$ algebraically closed, $\phi$ has an eigenvalue $c$; then $\phi-c$ is a
$G$-map that is not injective, so $\phi-c=0$ by (i). If $V\simeq W$ via $\alpha$, then $\Hom_G(V,W)=\alpha\End_G(V)=k\alpha$; if not, $\Hom_G(V,W)=0$
by (i).
(iii) By \Cref{prop:reps:scalars}(i), $\dim_k\End_G(V)=\dim_{\bar k}\End_G(V_{\bar k})=1$ and $\dim_k\Hom_G(V,W)=\dim_{\bar k}\Hom_G(V_{\bar k},W_{\bar k})$; apply (ii)
over $\bar k$, using \Cref{prop:reps:scalars}(ii) to compare $V\simeq W$ with $V_{\bar k}\simeq W_{\bar k}$.
(iv) A central $z$ gives $\rho(z)\in\End_G(V)=k$. If $G$ is abelian, every $\rho(g)$ is central in the image, so it is a scalar, every
line is stable, and $\dim V=1$. For the last statement let $A\subseteq\End_k(V)$ be the image of $k[G]$; it is a commutative subalgebra, and
$\End_G(V)$ is the commutant of $A$, which contains $A$. For $0\ne x\in A$, $\ker x$ is $A$-stable ($A$ is commutative), hence $G$-stable
and different from $V$, so $\ker x=0$: $A$ has no zero divisors, and a finite-dimensional commutative $k$-algebra without zero
divisors is a field. Then $V$ is an $A$-vector space whose $A$-subspaces are $G$-stable, so $\dim_AV=1$, $V\simeq A$, and the commutant
of $A$ acting on itself by multiplication is $A$.
\end{proof}

\begin{definition}\label{def:reps:absirr}
A representation $V$ of $G$ over $k$ is \emph{absolutely irreducible} if $V_{\bar k}$ is irreducible. It is \emph{absolutely reducible}
otherwise.
\end{definition}

\begin{proposition}\label{prop:reps:absirr}
Let $V$ be a representation of $G$ over $k$, $n=\dim V$.
\begin{enumerate}[label=(\roman*)]
\item $V$ is absolutely irreducible iff the map $k[G]\to\End_k(V)$ is surjective. If so, $V_K$ is irreducible for every extension $K/k$.
\item An absolutely irreducible $V$ is irreducible and $\End_G(V)=k$.
\item If $n=2$, $V$ is absolutely irreducible iff $\rho(G)$ stabilises no line of $\bar k^2$, iff $\rho(G)$ is not conjugate in $\GL_2(\bar k)$ into
the upper triangular group.
\item If $V$ is irreducible, $V$ is absolutely irreducible iff $\End_G(V)=k$. This holds over every field $k$.
\end{enumerate}
\end{proposition}

\begin{proof}
(i) Let $A\subseteq\End_k(V)$ be the image of $k[G]$. The image of $\bar k[G]$ in $\End_{\bar k}(V_{\bar k})=\bar k\otimes_k\End_k(V)$ is $\bar k\otimes_kA$, of
$\bar k$-dimension $\dim_kA$. If $V_{\bar k}$ is irreducible, the density theorem (\Cref{thm:reps:burnside}(b), whose proof does not use this
proposition) says that $\bar k\otimes_kA=\End(V_{\bar k})$, so $\dim_kA=n^2$ and $A=\End_k(V)$. Conversely, if $A=\End_k(V)$, then for every
extension $K/k$ the image of $K[G]$ is $\End_K(V_K)$, and a subspace stable under all of $\End_K(V_K)$ is $0$ or $V_K$.
(ii) If $U\subseteq V$ is a subrepresentation, $U_{\bar k}$ is one of $V_{\bar k}$; so $V$ is irreducible. $\End_G(V)=k$ is \Cref{thm:reps:schur}(iii).
(iii) A nonzero proper subrepresentation of the two-dimensional $V_{\bar k}$ is a stable line; and $\rho(G)$ stabilises the line $\bar ke_1'$ iff it is
upper triangular in a basis $e_1',e_2'$.
(iv) One direction is (ii). Conversely let $V$ be irreducible with $\End_G(V)=k$. The density theorem (\Cref{thm:reps:burnside}(a), with
$r=1$) says that the image of $k[G]$ is $\End_{\End_G(V)}(V)=\End_k(V)$, and (i) applies.
\end{proof}

\begin{proposition}[Distinct eigenvalues]\label{prop:reps:distinct-eigenvalues}
Let $V$ be a representation of any group $G$ over $k$, $n=\dim V$, and suppose that some $g\in G$ (or some element of the image of
$k[G]$) acts on $V$ with $n$ distinct eigenvalues $\lambda_1,\dots,\lambda_n\in k$. Then every $G$-stable subspace $U$ of $V_{\bar k}$ is a sum of
eigenlines of $g$, and $U=(U\cap V)_{\bar k}$ with $U\cap V$ a subrepresentation of $V$. In particular $V$ is irreducible iff it is absolutely
irreducible. This applies when $n=2$, $\operatorname{char}k\ne2$ and some $c\in G$ has $\rho(c)^2=1$, $\det\rho(c)=-1$: the eigenvalues of $\rho(c)$ are
then $1$ and $-1$, which are distinct.
\end{proposition}

\begin{proof}
Let $v_1,\dots,v_n\in V$ be eigenvectors of $g$ for $\lambda_1,\dots,\lambda_n$ (they exist over $k$ because the $\lambda_i$ lie in $k$); they form a basis of $V$
and of $V_{\bar k}$. Let $u=\sum_ic_iv_i\in U$. Applying $\prod_{j\ne i}(g-\lambda_j)$, which preserves $U$, gives $c_i\prod_{j\ne i}(\lambda_i-\lambda_j)v_i\in U$, and
the product is nonzero; so $c_iv_i\in U$ for every $i$. Hence $U$ is spanned by the $v_i$ it contains, $U=(U')_{\bar k}$ for the $k$-span $U'$
of these $v_i$, and $U'=U\cap V$ is $G$-stable. If $V$ is irreducible, $U\cap V\in\{0,V\}$, so $U\in\{0,V_{\bar k}\}$. The case $n=2$: $\rho(c)^2=1$
gives eigenvalues in $\{\pm1\}$, with product $-1$.
\end{proof}

\begin{example}[$C_4$ on $\F_3^2$]\label{ex:reps:c4}
Let $C_4=\langle J\rangle$, $J=\sm01{-1}0$, act on $\F_3^2$. It is irreducible (\Cref{def:reps:irreducible}) but not absolutely irreducible: over
$\F_9=\F_3[i]$,
\[
J\binom1i=\binom{i}{-1}=i\binom1i,\qquad J\binom1{-i}=\binom{-i}{-1}=-i\binom1{-i},
\]
so with $v=(1,i)^t$ and $\bar v=(1,-i)^t$ we get $\F_9^2=\F_9v\oplus\F_9\bar v\simeq\psi\oplus\psi^{-1}$, where $\psi(J)=i$ is a character of order $4$ with values
in $\F_9^\times$ (these are the two lines spanned by $(\mp i,1)^t$ of \Cref{ch:groups}). The vector $\bar v$ is obtained from $v$ by the Frobenius
$x\mapsto x^3$ of $\F_9$ on coordinates, and $\psi^{-1}=\psi^3$ is the Frobenius conjugate of $\psi$. The commutant is $\End_{C_4}(\F_3^2)=\F_3[J]\simeq\F_3[X]/(X^2+1)\simeq\F_9$
(\Cref{lem:reps:commutant}), a field of dimension $2$ over $\F_3$, as \Cref{prop:reps:absirr}(iv) and \Cref{thm:reps:schur}(iv) predict.
In contrast, $C_{ns}(\F_3)=\langle a\rangle\simeq\F_9^\times$ also acts irreducibly and not absolutely irreducibly, with commutant $\F_3[a]=\F_9$; and $\mathrm U$
fixes exactly one line of $\F_3^2$ and is not a direct sum of two lines (\Cref{ex:reps:maschke-fails}).
\end{example}

\begin{paperlink}
The hypothesis of Theorem thm:main is that $\rhobar_{E,\ell}$ is absolutely irreducible on $\Gal(\Qbar/\Q(\sqrt{\ell^*}))$; in DFG04 this is $\ell\notin S_f$.
\Cref{ex:reps:c4} is the reason why ``absolutely'' cannot be dropped: in the $N_s(\F_3)$ case the restriction to $\Gal(\Qbar/\Q(\sqrt{-3}))$ has
image $C_4$, which is irreducible over $\F_3$ and splits as $\psi\oplus\psi^{-1}$ over $\F_9$ (Proposition prop:ns). In Lemmas lem:h0 and lem:old and in
Proposition prop:Sf the paper deduces absolute irreducibility from irreducibility because complex conjugation has the distinct
eigenvalues $\pm1$: this is the case $n=2$ of \Cref{prop:reps:distinct-eigenvalues}, with $c$ complex conjugation and $\det\rhobar(c)=\bar\chi_\ell(c)=-1$
(\Cref{ch:galois}).
\end{paperlink}

\section{Semisimplicity, Maschke's theorem, Jordan-H\"older}\label{sec:reps:semisimple}

This section explains when a representation is a direct sum of irreducibles. Over a finite group of order prime to the
characteristic it always is (Maschke); in characteristic dividing the order it need not be, and this failure is exactly what the
paper's Lemma lem:IV records: an element of order $3$ in $\SL_2(\Z_3)$ fixes no vector of $T$ but a line of $T/3T$. The section also
proves the facts about extensions of two characters used in Lemma lem:old (a non-split extension has a unique stable line) and
gives a second proof that $Q_8$ acts absolutely irreducibly on $\F_3^2$, a step of Proposition prop:ns.

\begin{definition}\label{def:reps:semisimple}
A representation is \emph{semisimple} if it is a sum of irreducible subrepresentations. The zero representation is semisimple.
\end{definition}

\begin{proposition}\label{prop:reps:semisimple}
Let $V$ be a finite-dimensional representation of any group $G$ over $k$.
\begin{enumerate}[label=(\roman*)]
\item The following are equivalent: (a) $V$ is a sum of irreducible subrepresentations; (b) $V$ is a direct sum of irreducible
subrepresentations; (c) every subrepresentation of $V$ has a $G$-stable complement.
\item Subrepresentations and quotients of a semisimple representation are semisimple.
\item (Isotypic decomposition.) For an irreducible $S$ let $V_S$ be the sum of the subrepresentations of $V$ isomorphic to $S$. If $V$
is semisimple, $V=\bigoplus_SV_S$ over the isomorphism classes of irreducible $S$, and $V_S\simeq S^{m_S}$, where $m_S\dim_k\End_G(S)=\dim_k\Hom_G(S,V)$.
In particular $m_S=\dim\Hom_G(S,V)$ if $S$ is absolutely irreducible.
\end{enumerate}
\end{proposition}

\begin{proof}
(i) Suppose (a), $V=\sum_{i\in I}S_i$, and let $U$ be a subrepresentation. Choose $J\subseteq I$ maximal such that the sum $U+\sum_{j\in J}S_j$ is
direct. For every $i$, $S_i\cap(U\oplus\bigoplus_JS_j)$ is a subrepresentation of the irreducible $S_i$; if it were $0$, $J\cup\{i\}$ would
contradict maximality; so $S_i\subseteq U\oplus\bigoplus_JS_j$. Hence $V=U\oplus\bigoplus_JS_j$, which gives (c), and (b) for $U=0$. (b) implies (a)
trivially. Suppose (c), and let $V'$ be the sum of all irreducible subrepresentations. If $V'\ne V$, a complement $C\ne0$ exists; $C$
contains a nonzero subrepresentation of minimal dimension, which is irreducible and lies in $V'\cap C=0$, a contradiction.
(ii) By the proof of (i), $V=U\oplus\bigoplus_JS_j$, so $V/U\simeq\bigoplus_JS_j$ is semisimple. Applying this to a complement $C$ of $U$, with
$V=C\oplus U$, shows that $U\simeq V/C$ is semisimple.
(iii) Write $V=\bigoplus_iS_i$ by (b), and let $W_S$ be the sum of the $S_i$ isomorphic to $S$. If $T\subseteq V$ is irreducible and $T\simeq S$, its
projection to any $S_i\not\simeq S$ is $0$ by \Cref{thm:reps:schur}(i); so $T\subseteq W_S$. Hence $V_S=W_S$, $V=\bigoplus_SV_S$, and $V_S\simeq S^{m_S}$ with $m_S$
the number of $i$ with $S_i\simeq S$. Then $\Hom_G(S,V)=\bigoplus_i\Hom_G(S,S_i)$ has dimension $m_S\dim\End_G(S)$; the number $m_S$ does not
depend on the decomposition.
\end{proof}

\begin{theorem}[Maschke]\label{thm:reps:maschke}
\begin{enumerate}[label=(\roman*)]
\item Let $G$ be finite and $\operatorname{char}k\nmid\#G$. Every finite-dimensional representation of $G$ over $k$ is semisimple.
\item Let $G$ be finite, $R$ a commutative ring with $\#G\in R^\times$, and $M$ an $R[G]$-module. Then $e_G=\frac1{\#G}\sum_{g\in G}g$ is a central idempotent of
$R[G]$ with $e_GM=M^G$, $M=M^G\oplus(1-e_G)M$ as $R[G]$-modules, and $(M/IM)^G=M^G/IM^G$ for every ideal $I$ of $R$. In particular, for $R=\Zl$
with $\ell\nmid\#G$ and a $\Zl[G]$-lattice $T$, $(T/\ell^mT)^G=T^G/\ell^mT^G$: invariants commute with reduction.
\end{enumerate}
\end{theorem}

\begin{proof}
(i) Let $U\subseteq V$ be a subrepresentation and $\pi\colon V\to U$ any $k$-linear projection ($\pi|_U=\id$). Put $\pi'=\frac1{\#G}\sum_gg\pi g^{-1}$. Then
$\pi'(V)\subseteq U$, $\pi'|_U=\id$ (as $g\pi g^{-1}u=u$ for $u\in U$), and $h\pi'h^{-1}=\pi'$ for all $h$. So $\pi'$ is a $G$-equivariant projection
onto $U$, and $\ker\pi'$ is a stable complement; apply \Cref{prop:reps:semisimple}(i).
(ii) $he_G=e_G=e_Gh$ for $h\in G$, so $e_G$ is central and $e_G^2=\frac1{\#G}\sum_gge_G=e_G$. If $m\in M^G$, $e_Gm=m$; and $he_Gm=e_Gm$, so
$e_GM=M^G$. Since $e_G$ is a central idempotent, $M=e_GM\oplus(1-e_G)M$ as $R[G]$-modules. Finally $M/IM$ is an $R[G]$-module, so
$(M/IM)^G=e_G(M/IM)$, the image of $e_GM=M^G$; the kernel of $M^G\to M/IM$ is $M^G\cap IM$, and if $x=\sum r_im_i\in M^G$ with $r_i\in I$, then
$x=e_Gx=\sum r_ie_Gm_i\in IM^G$.
\end{proof}

\begin{paperlink}
In Lemma lem:euler and Remark rem:types, for $\ell\ge5$ the finite images of inertia (of order $2,3,4,6$) have order prime to $\ell$, so by (ii)
their invariants on $T$ and on $\ad T$ commute with reduction modulo $\ell^m$. At $\ell=3$ this fails only for the elements of order $3$ acting on $T$:
$T^g=0$ while $(T/3T)^g$ is a line (\Cref{prop:reps:reduction-cyclic}(ii)); that failure is the content of Lemma lem:IV. For an element of
order $6$, $T^g=0=(T/3T)^g$ (\Cref{prop:reps:reduction-cyclic}(iii)), and on $\ad T$ it acts as $-g$, of order $3$, whose invariant line
$\Z_3(1-2g)$ reduces onto $(\ad T/3)^g$ (\Cref{prop:reps:reduction-cyclic}(ii)): there the invariants still commute with reduction. In Proposition prop:ns the groups $Q_8$ and $C_4$ have order
prime to $3$, which is used in \Cref{prop:reps:dim2-absirr}.
\end{paperlink}

\begin{example}[Failures of semisimplicity]\label{ex:reps:maschke-fails}
\begin{enumerate}[label=(\alph*)]
\item $\mathrm U=\langle u\rangle\simeq C_3$ acting on $\F_3^2$. The line $\F_3e_1$ is stable. A stable line is spanned by an eigenvector of $u$, and
$\ker(u-1)=\F_3e_1$ is the only eigenspace; so $\F_3e_1$ has no stable complement, $\F_3^2$ is indecomposable and not semisimple, and its
composition factors (\Cref{thm:reps:jordan-holder}) are $\mathbf1,\mathbf1$. Maschke does not apply: $3\mid\#\mathrm U$.
\item The regular representation $\F_p[C_p]\simeq\F_p[x]/(x^p-1)=\F_p[x]/((x-1)^p)$ ($x$ a generator of $C_p$). Its subrepresentations are the ideals,
which are the $((x-1)^j)$, $0\le j\le p$, a chain: the representation is uniserial, with the single composition factor $\mathbf1$ of
multiplicity $p$.
\item For an infinite group Maschke fails even in characteristic $0$: $\Z$ acting on $\Ql^2$ by $m\mapsto\sm1{mn'}01$ with $n'\ne0$ stabilises
only the line $\Ql e_1$ (as in (a)), so it is not semisimple. This is the shape of tame inertia at a multiplicative prime
(Lemma lem:tamfactor).
\item $\mathrm B(\F_3)$ on $\F_3^2$: it contains $u$, so $\F_3e_1$ is its only stable line, and the representation is indecomposable and reducible.
Its commutant is $\F_3$: a matrix commuting with $\sigma=\mathrm{diag}(1,-1)$ is diagonal, and $\mathrm{diag}(x,y)$ commutes with $u$ iff $x=y$. So
``$\End_G(V)=k$'' does not imply irreducibility without semisimplicity.
\end{enumerate}
\end{example}

\begin{theorem}[Jordan-H\"older]\label{thm:reps:jordan-holder}
Let $A$ be a ring (here $k[G]$ or $\mathcal{O}[G]$) and $M$ an $A$-module of finite length, for example a finite-dimensional representation or a
finite $\mathcal{O}[G]$-module. Then $M$ has a composition series $0=M_0\subset M_1\subset\dots\subset M_r=M$ with simple quotients $M_i/M_{i-1}$, and the
length $r$ and the multiset of isomorphism classes of the quotients (the \emph{composition factors}, with multiplicity) do not depend
on the series.
\end{theorem}

\begin{proof}
This is proved in \Cref{ch:modules} for modules of finite length over any ring; we do not prove it again (decision B2). A
finite-dimensional representation has finite length because a strictly increasing chain of subspaces has at most $\dim V+1$ terms.
\end{proof}

\begin{definition}\label{def:reps:semisimplification}
The \emph{semisimplification} of a finite-dimensional representation $V$ is $V^{ss}=\bigoplus_{i=1}^rM_i/M_{i-1}$ for a composition series
$(M_i)$ of $V$.
\end{definition}

\begin{proposition}\label{prop:reps:semisimplification}
\begin{enumerate}[label=(\roman*)]
\item $V^{ss}$ is well defined up to isomorphism; $V^{ss}\simeq V$ iff $V$ is semisimple; and if $0\to U\to V\to Q\to0$ is exact, $V^{ss}\simeq U^{ss}\oplus Q^{ss}$.
\item $\det(X-g\mid V)=\det(X-g\mid V^{ss})$ and $\tr(g\mid V)=\tr(g\mid V^{ss})$ for all $g\in G$.
\item If $k$ is perfect and $K/k$ algebraic, $(V^{ss})_K\simeq(V_K)^{ss}$.
\end{enumerate}
\end{proposition}

\begin{proof}
(i) Well defined by \Cref{thm:reps:jordan-holder}. If $V=\bigoplus_iS_i$ with $S_i$ irreducible, $0\subset S_1\subset S_1\oplus S_2\subset\cdots$ is a
composition series with quotients $S_i$, so $V^{ss}\simeq V$; conversely $V^{ss}$ is semisimple by definition. For the exact sequence, a
composition series of $U$ followed by the preimages of one of $Q$ is a composition series of $V$.
(ii) In a basis adapted to a composition series the matrix of $g$ is block upper triangular with the matrices of $g$ on the
$M_i/M_{i-1}$ on the diagonal.
(iii) We do not prove this; see \cite[Section 29]{CurtisReiner} (over a perfect field an irreducible representation stays semisimple
after an algebraic extension of scalars). It is not used below: \Cref{thm:reps:brauer-nesbitt} is proved over an arbitrary field
without it.
\end{proof}

\begin{proposition}[Extensions of two characters]\label{prop:reps:nonsplit-line}
Let $V$ be a two-dimensional representation of any group $G$ over $k$ with a stable line $L$, on which $G$ acts by $\eta_1$, and let $G$
act on $V/L$ by $\eta_2$. In a basis $e_1\in L$, $e_2$,
\[
\rho(g)=\Mtwo{\eta_1(g)}{b(g)}0{\eta_2(g)},\qquad b(gh)=\eta_1(g)b(h)+b(g)\eta_2(h).
\]
\begin{enumerate}[label=(\roman*)]
\item The following are equivalent: $V$ is semisimple; $V\simeq\eta_1\oplus\eta_2$; $V$ has a stable line other than $L$; there is $c\in k$ with
$b(g)=c\,(\eta_2(g)-\eta_1(g))$ for all $g$ (then $k(ce_1+e_2)$ is stable).
\item If $V$ is not semisimple (a \emph{non-split extension}), $L$ is its only stable line. Hence if $V'$ is another non-split extension,
with sub-character $\eta_1'$ and quotient character $\eta_2'$, every isomorphism $V\simeq V'$ maps $L$ onto $L'$, and $\eta_1=\eta_1'$, $\eta_2=\eta_2'$.
\item If $\eta_1\ne\eta_2$ and $V$ is semisimple, $V$ has exactly two stable lines.
\end{enumerate}
\end{proposition}

\begin{proof}
The cocycle identity is the upper right entry of $\rho(gh)=\rho(g)\rho(h)$.
(i) If $V$ is semisimple, $L$ has a stable complement $L'$, a second stable line, and $L'\simeq V/L$, so $V=L\oplus L'\simeq\eta_1\oplus\eta_2$; if
$V\simeq\eta_1\oplus\eta_2$ it is semisimple. A line other than $L$ is $k(ce_1+e_2)$ for a unique $c$, and
$g(ce_1+e_2)=(c\eta_1(g)+b(g))e_1+\eta_2(g)e_2$ lies in $k(ce_1+e_2)$ iff $c\eta_1(g)+b(g)=c\eta_2(g)$. A second stable line is a complement of $L$.
(ii) A second stable line would make $V$ semisimple by (i). An isomorphism maps stable lines to stable lines, so $L$ to $L'$, and
induces isomorphisms $L\simeq L'$ and $V/L\simeq V'/L'$ of one-dimensional representations, which are equalities of characters.
(iii) In a basis $v_1,v_2$ of stable lines, $V=kv_1\oplus kv_2$ with characters $\eta_1\ne\eta_2$, a vector $xv_1+yv_2$ with $xy\ne0$ spans a stable line only if
$\eta_1(g)=\eta_2(g)$ for all $g$.
\end{proof}

\begin{paperlink}
Lemma lem:old, case $p\,\|\,N$ with $\rhobar$ ramified at $p$: ``a non-split extension of two characters has a unique stable line, so the
isomorphism $\rhobar|_{D_p}\simeq\rhobar_h|_{D_p}$ matches the sublines and hence the quotients'' is (ii), applied over a common finite extension
of $\F_\ell$. Lemma lem:tamfactor uses the shape $\sm{\chi_\ell\varepsilon}{\varepsilon\kappa}0\varepsilon$ of the Tate representation: the corner
$\varepsilon\kappa$ satisfies the cocycle identity above with $\eta_1=\chi_\ell\varepsilon$, $\eta_2=\varepsilon$ (convention B-C14; the cohomological reading is in
\Cref{ch:cohomology}).
\end{paperlink}

\begin{proposition}[Absolute irreducibility in dimension $2$]\label{prop:reps:dim2-absirr}
Let $\rho\colon G\to\GL_2(k)$ be a representation.
\begin{enumerate}[label=(\roman*)]
\item If $\rho(G)$ is abelian, $\rho$ is absolutely reducible.
\item If $\rho(G)$ is finite of order prime to $\operatorname{char}k$ and $\rho$ is absolutely reducible, then $\rho_{\bar k}\simeq\eta_1\oplus\eta_2$ and $\rho(G)$
is abelian. Hence a finite non-abelian subgroup of $\GL_2(k)$ of order prime to $\operatorname{char}k$ acts absolutely irreducibly on $k^2$.
\item If $H\le G$ and $\Res_H\rho$ is absolutely irreducible, so is $\rho$.
\item On $\F_3^2$: $Q_8$ acts absolutely irreducibly, and hence so do $\SL_2(\F_3)$, $N_{ns}(\F_3)$ and $\GL_2(\F_3)$, which contain it; $N_s(\F_3)\simeq D_4$
acts absolutely irreducibly; $C_4$ and $C_{ns}(\F_3)$ do not. The Borel subgroup $\mathrm B(\F_3)$ is non-abelian and acts faithfully but
reducibly: its order $12$ is divisible by $3$, and (ii) does not apply.
\end{enumerate}
\end{proposition}

\begin{proof}
(i) Commuting operators on $\bar k^2$ have a common eigenvector: if all $\rho(g)$ are scalars any line is stable; otherwise take $g_0$
with $\rho(g_0)$ not scalar and an eigenvalue $c$ of it; its eigenspace $\ker(\rho(g_0)-c)$ is a line (it is not everything), stable under
every $\rho(g)$, which commutes with $\rho(g_0)$.
(ii) $\rho_{\bar k}$ has a stable line, and by Maschke (\Cref{thm:reps:maschke}(i), applied to the finite group $\rho(G)$) a stable
complement; in the basis of the two lines $\rho(G)$ is diagonal, hence abelian.
(iii) A stable line of $\bar k^2$ for $G$ is one for $H$.
(iv) $Q_8$ is non-abelian of order $8$, prime to $3$, and acts faithfully; apply (ii) and then (iii) to the groups containing $Q_8$
(\Cref{ch:groups}). $N_s(\F_3)$ is non-abelian of order $8$. $C_4$ and $C_{ns}(\F_3)$ are cyclic; apply (i). $\mathrm B(\F_3)$ stabilises $\F_3e_1$.
\end{proof}

\begin{paperlink}
Proposition prop:ns: ``The first two act absolutely irreducibly on $\F_3^2$''. The paper's argument is the commutator one of
\Cref{ch:groups}: a subgroup of $\GL_2(\Fbar_3)$ fixing a line lies in a Borel subgroup, whose commutator subgroup is unipotent, so abelian
of exponent $3$, while the commutator subgroups $\{\pm1\}$ of $Q_8$ and $Q_8$ of $\SL_2(\F_3)$ are not. \Cref{prop:reps:dim2-absirr}(ii), (iii) give
a second proof: $Q_8$ is non-abelian of order prime to $3$, and $\SL_2(\F_3)$ contains $Q_8$. The phrase ``non-abelian and faithful'' of
the original version of the paper is not enough, by the example of $\mathrm B(\F_3)$, which is why the revised version argues with commutators. ``The third is abelian,
hence absolutely reducible'' is (i).
\end{paperlink}

\section{Characters in characteristic zero}\label{sec:reps:characters}

This short section connects the chapter with the character theory of Riley-Hobson-Bence: over $\C$ a representation of a finite
group is determined by its trace function, and the trace functions of the irreducible representations are orthonormal. The paper
uses none of this directly; it is here for two reasons. It shows the reader where the familiar theory stops: in characteristic $\ell$
traces no longer determine semisimple representations (\Cref{ex:reps:traces-fail}), which is why Brauer-Nesbitt in Lemma lem:old is
applied with characteristic polynomials. And the two groups of order $8$ of Proposition prop:ns, $N_s(\F_3)\simeq D_4$ and $Q_8$, are a classical pair
with the same character table.

\begin{definition}\label{def:reps:trace-character}
The \emph{trace function} of a representation $V$ is $\tr_V\colon G\to k$, $\tr_V(g)=\tr\rho(g)$. It is a class function ($\tr_V(xgx^{-1})=\tr_V(g)$),
$\tr_V(1)=\dim V$ (in $k$), and by \Cref{prop:reps:basic-facts} $\tr_{V\oplus W}=\tr_V+\tr_W$, $\tr_{V\otimes W}=\tr_V\tr_W$, $\tr_{V^\vee}(g)=\tr_V(g^{-1})$. For
$k=\C$ and $G$ finite, $\rho(g)$ has finite order, its eigenvalues are roots of unity, and $\tr_V(g^{-1})=\overline{\tr_V(g)}$. In
Riley-Hobson-Bence, and in most books on finite groups, $\tr_V$ is called the \emph{character} of $V$; we avoid the word for trace
functions because the paper's characters $\chi_{-3}$, $\chi_\ell$, $\psi$ are homomorphisms to $k^\times$ (one-dimensional representations, whose trace
function is the character itself).
\end{definition}

\begin{theorem}[Orthogonality]\label{thm:reps:orthogonality}
Let $G$ be finite, $k=\C$, and $\langle\phi,\psi\rangle=\frac1{\#G}\sum_{g\in G}\phi(g)\overline{\psi(g)}$ on class functions.
\begin{enumerate}[label=(\roman*)]
\item For irreducible $V,W$: $\langle\tr_V,\tr_W\rangle=1$ if $V\simeq W$ and $0$ otherwise.
\item For all $V,W$: $\dim\Hom_G(V,W)=\langle\tr_W,\tr_V\rangle$. A representation $V$ is irreducible iff $\langle\tr_V,\tr_V\rangle=1$, and two
representations with the same trace function are isomorphic.
\item The trace functions $\tr_{V_1},\dots,\tr_{V_r}$ of the irreducible representations (up to isomorphism) form an orthonormal basis of
the space of class functions; so $r$ is the number of conjugacy classes; and $\sum_i(\dim V_i)^2=\#G$.
\item Each $\dim V_i$ divides $\#G$.
\end{enumerate}
\end{theorem}

\begin{proof}
(ii) By \Cref{prop:reps:homG}, $\Hom_G(V,W)=(V^\vee\otimes W)^G$, the image of the idempotent $e_G$ of \Cref{thm:reps:maschke}(ii). The trace of an
idempotent is its rank, so
\begin{align*}
\dim\Hom_G(V,W)&=\tr(e_G\mid V^\vee\otimes W)=\frac1{\#G}\sum_g\tr_{V^\vee}(g)\tr_W(g)\\
&=\frac1{\#G}\sum_g\tr_W(g)\overline{\tr_V(g)}=\langle\tr_W,\tr_V\rangle.
\end{align*}
(i) follows from (ii) and \Cref{thm:reps:schur}(ii); the value is real, so the order of the arguments does not matter. By Maschke,
$V\simeq\bigoplus_iV_i^{m_i}$, and by (i) $\langle\tr_V,\tr_V\rangle=\sum_im_i^2$, which is $1$ iff $V$ is irreducible; and $m_i=\langle\tr_V,\tr_{V_i}\rangle$ is determined
by $\tr_V$.
(iii) Orthonormality is (i). The regular representation $\C[G]$ has $\tr(1)=\#G$ and $\tr(g)=0$ for $g\ne1$ (a nontrivial $g$ moves every
basis vector), so $\langle\tr_{\C[G]},\tr_{V_i}\rangle=\dim V_i$, $\C[G]\simeq\bigoplus_iV_i^{\dim V_i}$, and $\#G=\sum_i(\dim V_i)^2$. Let $\phi$ be a class function
orthogonal to every $\tr_{V_i}$. For each irreducible $V$, the operator $\Phi_V=\sum_g\overline{\phi(g)}\rho_V(g)$ commutes with $G$ ($\phi$ is a class
function), so it is a scalar $c$ by Schur, and $c\dim V=\tr\Phi_V=\#G\,\langle\tr_V,\phi\rangle=0$. So $\Phi$ vanishes on every irreducible, hence on $\C[G]$,
and applied to $1\in\C[G]$ it gives $\sum_g\overline{\phi(g)}g=0$, so $\phi=0$. Hence the $\tr_{V_i}$ span the class functions, whose dimension is the
number of conjugacy classes.
(iv) We do not prove this; see \cite[Section 6.5]{SerreLinearReps}.
\end{proof}

\begin{corollary}\label{cor:reps:traces-char0}
If $\operatorname{char}k=0$, two semisimple finite-dimensional representations of any group $G$ with the same trace function are
isomorphic.
\end{corollary}

\begin{proof}
This is the last statement of \Cref{thm:reps:brauer-nesbitt} (traces suffice in characteristic $0$), together with
\Cref{prop:reps:semisimplification}(i); the proof of that theorem does not use this corollary.
\end{proof}

\begin{example}[The character tables of $D_4$ and $Q_8$]\label{ex:reps:tables}
Take $D_4=N_s(\F_3)=\langle J,\sigma\rangle$ (rotation $J$, reflection $\sigma$; \Cref{ch:groups}) and $Q_8=\{\pm1,\pm J,\pm j,\pm k\}$ with $j=\sm1{-1}{-1}{-1}$, $k=Jj$.
Both have five conjugacy classes, and their irreducible trace functions over $\C$ are:
\begin{center}
\begin{tabular}{lccccc}
\toprule
$D_4$ & $1$ & $-1$ & $\pm J$ & $\pm\sigma$ & $\pm w$\\
$Q_8$ & $1$ & $-1$ & $\pm J$ & $\pm j$ & $\pm k$\\
\midrule
$\lambda_0$ & $1$ & $1$ & $1$ & $1$ & $1$\\
$\lambda_1$ & $1$ & $1$ & $1$ & $-1$ & $-1$\\
$\lambda_2$ & $1$ & $1$ & $-1$ & $1$ & $-1$\\
$\lambda_3$ & $1$ & $1$ & $-1$ & $-1$ & $1$\\
$\tr_2$ & $2$ & $-2$ & $0$ & $0$ & $0$\\
\bottomrule
\end{tabular}
\end{center}
The four characters factor through the quotient by $\{\pm1\}$, which is $V_4$ for both groups; for $D_4$, $\lambda_1=\det$ (kernel $C_4$), $\lambda_2$ has kernel
$C_s(\F_3)=\{\pm1,\pm\sigma\}$ and $\lambda_3$ kernel $\{\pm1,\pm w\}$. The orthogonality relations and $4\cdot1^2+2^2=8$ show the table is complete. Yet
$D_4\not\simeq Q_8$ ($D_4$ has five involutions, $Q_8$ one). The two-dimensional representation of $D_4$ is given by the integral matrices $J$, $\sigma$,
so it is defined over $\Q$, and its reduction modulo $3$ is the standard representation of $N_s(\F_3)$. That of $Q_8$ is not defined over
$\R$: its Frobenius-Schur indicator $\frac18\sum_g\tr_2(g^2)$ is $\frac18(2\cdot2+6\cdot(-2))=-1$, which for an irreducible representation means that it
is not realisable over $\R$ (we do not prove this; see \cite[Section 13.2]{SerreLinearReps}). Over $\F_3$, however, it is realised by the
standard representation of $Q_8\subset\SL_2(\F_3)$, whose trace function is the reduction of $\tr_2$.
\end{example}

\begin{example}[Traces do not suffice in characteristic $p$]\label{ex:reps:traces-fail}
Let $G=C_2=\{1,s\}$ and $k=\F_3$. The semisimple representations $\mathbf1^{\oplus3}$ and $\sgn^{\oplus3}$ ($s\mapsto-1$) have trace functions $(3,3)$ and
$(3,-3)$, both $(0,0)$ in $\F_3$, but they are not isomorphic: $\det(X-s)$ is $(X-1)^3$, respectively $(X+1)^3$. In general, if $n<p$ (or
$p=0$), the traces of $g,g^2,\dots,g^n$ determine the characteristic polynomial of $g$ on an $n$-dimensional space by Newton's identities,
$je_j=\sum_{i=1}^{j}(-1)^{i-1}e_{j-i}p_i$ for the elementary symmetric functions $e_j$ and power sums $p_i=\tr(g^i)$ of the eigenvalues, which
can be solved for the $e_j$ because $1,\dots,n$ are invertible. For $n=2$ the characteristic polynomial $X^2-\tr(g)X+\det(g)$ is determined
by the trace and the determinant in every characteristic.
\end{example}

\begin{warning}
In characteristic $\ell$ the right invariant of a representation is the family of characteristic polynomials, not the trace function.
The paper's representations have dimension $2$ and determinant $\bar\chi_\ell$, so for them traces do suffice
(\Cref{cor:reps:bn-consequences}(i)): ``$a_p(h)\equiv a_p(f)\pmod\lambda$ for all $p\nmid\ell N$'' is a statement about traces of Frobenius
elements, and it determines the semisimplification.
\end{warning}

\section{The theorem of Brauer and Nesbitt}\label{sec:reps:brauer-nesbitt}

This section proves that a semisimple representation over any field is determined by the characteristic polynomials of the
elements of the group (Brauer and Nesbitt). With the Chebotarev density theorem (\Cref{ch:galois}) it is the reason why a congruence
$a_p(h)\equiv a_p(f)\pmod\lambda$ for almost all $p$ gives $\rhobar_h\simeq\rhobar$, which is the first step of the paper's Lemma lem:old. The proof needs the
density theorem of Jacobson and Burnside, proved first.

\begin{theorem}[Density theorem of Jacobson and Burnside]\label{thm:reps:burnside}
Let $G$ be any group, $S_1,\dots,S_r$ pairwise non-isomorphic irreducible finite-dimensional representations over $k$, and
$D_i=\End_G(S_i)$.
\begin{enumerate}[label=(\alph*)]
\item The image of $k[G]\to\prod_i\End_k(S_i)$ is $\prod_i\End_{D_i}(S_i)$.
\item If $k$ is algebraically closed, the image is $\prod_i\End_k(S_i)$. In particular, for each $i$ there is $x\in k[G]$ acting on $S_i$ as a
projector of rank one and on every $S_j$, $j\ne i$, as $0$.
\end{enumerate}
\end{theorem}

\begin{proof}
(a) Let $A$ be the image. Each element of $k[G]$ commutes with $D_i$ on $S_i$, so $A\subseteq\prod_i\End_{D_i}(S_i)$. By \Cref{thm:reps:schur}(i), $D_i$ is a
division algebra and $S_i$ is a finite-dimensional (left) $D_i$-vector space; let $x_{i1},\dots,x_{in_i}$ be a $D_i$-basis. Put
$M=\bigoplus_iS_i^{n_i}$, a semisimple representation, and $x=(x_{ij})_{i,j}\in M$. The submodule $k[G]x=Ax$ has a $G$-stable complement
(\Cref{prop:reps:semisimple}), so there is a $G$-equivariant projection $\pi\in\End_G(M)$ with image $Ax$. Since $\Hom_G(S_i,S_j)=0$ for $i\ne j$
(Schur), $\End_G(M)=\prod_i\mathrm M_{n_i}(D_i)$, the matrices with entries in $D_i$ acting on $S_i^{n_i}$. Now let $\phi=(\phi_i)\in\prod_i\End_{D_i}(S_i)$ and let $\Phi$ act on
$M$ by $\phi_i$ on each copy of $S_i$. As $\phi_i$ commutes with $D_i$, $\Phi$ commutes with every matrix in $\prod_i\mathrm M_{n_i}(D_i)$, in particular with
$\pi$. Hence $\Phi(x)=\Phi(\pi x)=\pi\Phi(x)\in Ax$: there is $\alpha\in A$ with $\alpha x_{ij}=\phi_i(x_{ij})$ for all $i,j$. On $S_i$, both $\alpha$ and $\phi_i$ are
$D_i$-linear and agree on a $D_i$-basis, so they are equal. Thus $\phi=\alpha\in A$.
(b) If $k=\bar k$ then $D_i=k$ by \Cref{thm:reps:schur}(ii), so $\End_{D_i}(S_i)=\End_k(S_i)$; take the element of $A$ that is a rank-one projector
on $S_i$ and $0$ on the other factors.
\end{proof}

\begin{theorem}[Brauer-Nesbitt]\label{thm:reps:brauer-nesbitt}
Let $G$ be any group, $k$ any field, and $V_1,V_2$ finite-dimensional representations of $G$ over $k$ with
\[
\det(X-\rho_1(g))=\det(X-\rho_2(g))\quad\text{in }k[X]\text{ for every }g\in G.
\]
Then $V_1$ and $V_2$ have the same composition factors with multiplicities, that is, $V_1^{ss}\simeq V_2^{ss}$. If $\operatorname{char}k=0$, or if
$\dim V_1<\operatorname{char}k$ and $\dim V_2<\operatorname{char}k$, the equality $\tr\rho_1(g)=\tr\rho_2(g)$ for all $g$ suffices.
\end{theorem}

\begin{proof}
We prove both statements together; ``the hypothesis'' means equal characteristic polynomials, or, in the two special cases, equal
traces.

\emph{Step 0.} By \Cref{prop:reps:semisimplification}(ii) we may replace $V_1,V_2$ by $V_1^{ss},V_2^{ss}$.

\emph{Step 1: reduction to $k=\bar k$.} Let $S,S'$ be non-isomorphic irreducible representations over $k$. By \Cref{thm:reps:burnside}(a) with
$r=2$ there is $x\in k[G]$ acting as $1$ on $S$ and as $0$ on $S'$. Then $x$ acts as $1$ on every composition factor of $S_{\bar k}$ and as $0$
on every composition factor of $S'_{\bar k}$, so $S_{\bar k}$ and $S'_{\bar k}$ have no common composition factor. If $V$ has composition factors $S$ with
multiplicities $m_S(V)$, a composition series of $V$ gives a filtration of $V_{\bar k}$ with quotients $S_{\bar k}$, so the multiset of composition
factors of $V_{\bar k}$ is $\sum_Sm_S(V)\,c(S_{\bar k})$, where the multisets $c(S_{\bar k})$ of composition factors of the $S_{\bar k}$ are nonempty with pairwise
disjoint supports. Hence the composition factors of $V_{1,\bar k}$ and $V_{2,\bar k}$ agree iff $m_S(V_1)=m_S(V_2)$ for all $S$. As the hypothesis is
unchanged by extending scalars, we may assume $k=\bar k$, and (Step 0) that $V_1,V_2$ are semisimple.

\emph{Step 2: cancellation.} Write $V_1\simeq C\oplus B_1$ and $V_2\simeq C\oplus B_2$ where $C=\bigoplus_SS^{\min(m_S(V_1),m_S(V_2))}$; then $B_1$ and $B_2$ are
semisimple without a common irreducible constituent. We show $B_1=B_2=0$. The hypothesis passes to $B_1,B_2$: traces by subtraction,
and characteristic polynomials because $\det(X-g\mid V_i)=\det(X-g\mid C)\det(X-g\mid B_i)$ and $k[X]$ is a domain. We argue by induction on
$\dim B_1+\dim B_2$; if it is $0$ there is nothing to prove.

\emph{Step 3: multiplicities modulo $p$.} The traces of $B_1$ and $B_2$ agree on $G$ (they are coefficients of the characteristic
polynomials), hence on all of $k[G]$ by linearity. Let $S$ be an irreducible constituent of $B_1$, of multiplicity $m\ge1$. By
\Cref{thm:reps:burnside}(b) applied to the distinct irreducible constituents of $B_1\oplus B_2$, there is $x\in k[G]$ acting on $S$ as a projector of
rank one and as $0$ on all the other constituents; $S$ does not occur in $B_2$, so $m\cdot1=\tr(x\mid B_1)=\tr(x\mid B_2)=0$ in $k$. If
$\operatorname{char}k=0$ this is impossible, so $B_1=0$, and symmetrically $B_2=0$. If $p=\operatorname{char}k>0$, then $p\mid m$; if moreover
$\dim V_1<p$, then $m\le\dim V_1<p$ is impossible, and again $B_1=B_2=0$. This proves the statements about traces.

\emph{Step 4: dividing by $p$.} In the remaining case ($p>0$, characteristic polynomials) every multiplicity in $B_1$ and in $B_2$ is divisible
by $p$, so $B_i\simeq B_i'^{\oplus p}$ with $B_i'$ semisimple, and $\det(X-g\mid B_i)=\det(X-g\mid B_i')^p$. The map $f\mapsto f^p$ is an injective ring
endomorphism of $k[X]$ (it is additive in characteristic $p$, and $k[X]$ is a domain), so $\det(X-g\mid B_1')=\det(X-g\mid B_2')$ for all $g$. The
representations $B_1',B_2'$ have no common constituent and $\dim B_1'+\dim B_2'<\dim B_1+\dim B_2$ unless both are $0$; by induction $B_1'=B_2'=0$,
so $B_1=B_2=0$.
\end{proof}

\begin{corollary}\label{cor:reps:bn-consequences}
\begin{enumerate}[label=(\roman*)]
\item If $\dim V_1=\dim V_2=2$, $\det\rho_1=\det\rho_2$ and $\tr\rho_1=\tr\rho_2$ on $G$, then $V_1^{ss}\simeq V_2^{ss}$ (in every characteristic).
\item Let $k\subseteq k'$ be fields, $V_1$ a representation over $k$ and $V_2$ one over $k'$, with $\det(X-\rho_1(g))=\det(X-\rho_2(g))$ for all $g$. If
$V_1\otimes_kk'$ is irreducible (for instance if $V_1$ is absolutely irreducible), then $V_2\simeq V_1\otimes_kk'$; no semisimplicity of $V_2$ is
assumed.
\item The irreducibility of $V_1\otimes_kk'$ cannot be weakened to that of $V_1$. First, the irreducibility of $V_1$ does not make $V_2$
irreducible: for $C_4$ acting on $V_1=\F_3^2$ and $V_2=\psi\oplus\psi^{-1}$ over $\F_9$ (\Cref{ex:reps:c4}), $V_1$ is irreducible and $V_2$ is not. This
example does not contradict the conclusion of (ii), since here $V_2\simeq V_1\otimes\F_9$. A genuine counterexample to the conclusion is the
following. For the group $G=\bigl\{\sm x y0{x^3}:x\in\F_9^\times,y\in\F_9\bigr\}$, its standard
representation $V_2$ on $\F_9^2$ and the representation $V_1$ on $\F_9=\F_3^2$ by multiplication by $x$ have the same characteristic
polynomials $(X-x)(X-x^3)$ and $V_1$ is irreducible over $\F_3$, but $V_2$ is not semisimple while $V_1\otimes\F_9$ is, so $V_2\not\simeq V_1\otimes\F_9$.
\end{enumerate}
\end{corollary}

\begin{proof}
(i) The characteristic polynomial of a $2\times2$ matrix is $X^2-\tr(g)X+\det(g)$; apply \Cref{thm:reps:brauer-nesbitt}.
(ii) Apply \Cref{thm:reps:brauer-nesbitt} over $k'$ to $V_1\otimes_kk'$ and $V_2$: $V_2^{ss}\simeq V_1\otimes_kk'$ is irreducible, so $V_2$ has a single composition
factor, is irreducible, and equals $V_2^{ss}$.
(iii) In the second example, an element $x\in\F_9^\times$ of order $8$ acts on $\F_9$ by multiplication with eigenvalues $x,x^3$ over $\F_9$; the
$\F_3$-subspaces of $\F_9$ stable under multiplication by $\F_9^\times$ are $0$ and $\F_9$, so $V_1$ is irreducible, and $V_1\otimes\F_9$ is the sum of
the two characters $g\mapsto x$, $g\mapsto x^3$. In $V_2$ the line $\F_9e_1$ is stable and is the only stable line, because $G$ contains $\sm1101$;
so $V_2$ is not semisimple (\Cref{prop:reps:nonsplit-line}).
\end{proof}

\begin{paperlink}
Lemma lem:old: ``$\rhobar_h\simeq\rhobar$ over a common finite extension of $\F_\ell$, by Chebotarev and Brauer-Nesbitt, $\rhobar$ being irreducible, hence
absolutely irreducible''. The residual representation $\rhobar_h$ of the newform $h$ lives over the residue field $k'$ of the coefficient
ring of $h$, and $\rhobar$ over $\F_\ell$. For $p\nmid\ell N$ both have characteristic polynomial $X^2-a_pX+p$ at $\Frob_p$, with
$a_p(h)\equiv a_p(f)$; by Chebotarev these elements are dense, so the characteristic polynomials agree on all of $G_\Q$ (\Cref{ch:galois}).
Then \Cref{cor:reps:bn-consequences}(ii), with $\rhobar$ absolutely irreducible (\Cref{prop:reps:distinct-eigenvalues}), gives
$\rhobar_h\simeq\rhobar\otimes k'$; with only $\rhobar$ irreducible over $\F_\ell$ one would get $\rhobar_h^{ss}\simeq(\rhobar\otimes k')^{ss}$ and nothing more, by (iii).
\end{paperlink}

\section[Restriction, induction, Mackey and Clifford]{Restriction, induction, Frobenius reciprocity, Mackey and Clifford}\label{sec:reps:induction}

This section develops induced representations and their behaviour under restriction to a subgroup of index $2$. In the paper,
the image $N_s(\F_3)$ of Proposition prop:ns gives a representation induced, over $\F_9$, from a character of order $4$ of
$\Gal(\Qbar/\Q(\sqrt{-3}))$, and for $\rhobar$ irreducible the condition ``absolutely irreducible on $\Gal(\Qbar/\Q(\sqrt{\ell^*}))$'' of Theorem thm:main
is, by Clifford theory (\Cref{thm:reps:clifford-index-two}), the condition ``not induced from $\Q(\sqrt{\ell^*})$'', equivalently $\rhobar\not\simeq\rhobar\otimes\chi_{\ell^*}$
(Lemma lem:h0 for $\ell=3$, Proposition prop:Sf).

\begin{definition}\label{def:reps:induction}
Let $H\le G$ and let $W$ be a representation of $H$ over $k$ (or an $R[H]$-module). The \emph{induced representation} is
$\Ind_H^GW=k[G]\otimes_{k[H]}W$, with $G$ acting on the left factor. For a representation $V$ of $G$, $\Res^G_HV$ (or $\Res_HV$) is $V$ viewed as a
representation of $H$. For $s\in G$ and a representation $W$ of a subgroup $H$, $W^s$ is the representation of $sHs^{-1}$ on the space $W$
given by $x\mapsto\rho_W(s^{-1}xs)$. If $H$ is normal, $W^s$ is again a representation of $H$; for a character $\psi$ of $H$, $\psi^s(x)=\psi(s^{-1}xs)$.
\end{definition}

\begin{proposition}\label{prop:reps:induction}
Let $H\le G$ have finite index $m$ and choose $G=\bigsqcup_{i=1}^ms_iH$ with $s_1=1$.
\begin{enumerate}[label=(\roman*)]
\item $\Ind_H^GW=\bigoplus_is_i\otimes W$, $\dim\Ind_H^GW=m\dim W$, and if $gs_i=s_jx$ with $x\in H$, then $g(s_i\otimes w)=s_j\otimes xw$. The subspace
$s_i\otimes W$ is $s_iHs_i^{-1}$-stable and isomorphic to $W^{s_i}$.
\item $\tr(g\mid\Ind_H^GW)=\sum_{i:\,s_i^{-1}gs_i\in H}\tr(s_i^{-1}gs_i\mid W)$.
\item Let $[G:H]=2$, $s\notin H$. Then $\tr(x\mid\Ind W)=\tr(x\mid W)+\tr(s^{-1}xs\mid W)$ for $x\in H$ and $\tr(g\mid\Ind W)=0$ for $g\notin H$. For a
character $\psi$ of $H$, in the basis $1\otimes1$, $s\otimes1$ of $\Ind_H^G\psi$,
\[
x\mapsto\Mtwo{\psi(x)}00{\psi^s(x)}\ (x\in H),\qquad s\mapsto\Mtwo0{\psi(s^2)}10.
\]
\end{enumerate}
\end{proposition}

\begin{proof}
(i) $k[G]$ is a free right $k[H]$-module with basis $s_1,\dots,s_m$, so $k[G]\otimes_{k[H]}W=\bigoplus_is_i\otimes W$, and
$g(s_i\otimes w)=s_jx\otimes w=s_j\otimes xw$. For $y=s_ixs_i^{-1}$, $y(s_i\otimes w)=s_ix\otimes w=s_i\otimes xw$, so $w\mapsto s_i\otimes w$ is an isomorphism
$W^{s_i}\to s_i\otimes W$.
(ii) By (i), $g$ maps the block $s_i\otimes W$ into $s_j\otimes W$; the diagonal blocks are those with $j=i$, that is, $s_i^{-1}gs_i=x\in H$, where $g$ acts
as $x$ on $W$.
(iii) With $s_1=1$, $s_2=s$: for $x\in H$ both blocks are diagonal ($s^{-1}xs\in H$ as $H$ is normal), and for $g\notin H$ neither is. For $\psi$:
$x(1\otimes1)=1\otimes\psi(x)$, $x(s\otimes1)=s(s^{-1}xs)\otimes1=\psi^s(x)\,s\otimes1$, $s(1\otimes1)=s\otimes1$ and $s(s\otimes1)=s^2\otimes1=\psi(s^2)\,1\otimes1$.
\end{proof}

\begin{theorem}[Frobenius reciprocity]\label{thm:reps:frobenius-reciprocity}
Let $H\le G$, $W$ a representation of $H$ and $V$ one of $G$, over a field or a commutative ring.
\begin{enumerate}[label=(\roman*)]
\item The map $\phi\mapsto\phi(1\otimes{-})$ is an isomorphism $\Hom_G(\Ind_H^GW,V)\simeq\Hom_H(W,\Res_HV)$.
\item If $[G:H]$ is finite, with $\pi\colon\Ind_H^GW\to W$ the projection onto $1\otimes W$ along $\bigoplus_{i\ge2}s_i\otimes W$, the map $\phi\mapsto\pi\circ\phi$ is an
isomorphism $\Hom_G(V,\Ind_H^GW)\simeq\Hom_H(\Res_HV,W)$.
\end{enumerate}
In characteristic $0$ with $G$ finite, (i) and \Cref{thm:reps:orthogonality}(ii) give $\langle\tr_{\Ind W},\tr_V\rangle_G=\langle\tr_W,\tr_{\Res V}\rangle_H$.
\end{theorem}

\begin{proof}
(i) $w\mapsto\phi(1\otimes w)$ is $H$-linear because $\phi(1\otimes xw)=\phi(x\otimes w)=x\phi(1\otimes w)$. The inverse sends $\psi\in\Hom_H(W,V)$ to
$y\otimes w\mapsto y\psi(w)$, which is well defined on $k[G]\otimes_{k[H]}W$ ($yx\otimes w$ and $y\otimes xw$ have the same image, $\psi$ being $H$-linear) and
$G$-linear.
(ii) $\pi$ is $H$-linear: $H$ preserves $1\otimes W$, and permutes the blocks $s_i\otimes W$, $i\ge2$ (if $xs_i\in s_1H=H$ then $s_i\in H$). Given $\psi\in
\Hom_H(V,W)$ put $\Psi(v)=\sum_is_i\otimes\psi(s_i^{-1}v)$. This does not depend on the representatives, since $s_ix\otimes\psi(x^{-1}s_i^{-1}v)=s_i\otimes\psi(s_i^{-1}v)$.
It is $G$-linear: $g\Psi(v)=\sum_igs_i\otimes\psi(s_i^{-1}v)$, and $t_i=gs_i$ is another system of representatives with $t_i^{-1}gv=s_i^{-1}v$, so
$g\Psi(v)=\sum_it_i\otimes\psi(t_i^{-1}gv)=\Psi(gv)$. Then $\pi\Psi=\psi$ (the term $i=1$). Conversely, for a $G$-map $\phi$ write $\phi(v)=\sum_is_i\otimes w_i$;
then $\pi(\phi(s_i^{-1}v))=\pi(s_i^{-1}\phi(v))=w_i$, because $s_i^{-1}s_j\in H$ only for $j=i$; so the map $\Psi$ built from $\psi=\pi\phi$ is $\phi$.
\end{proof}

\begin{proposition}[Projection formula]\label{prop:reps:projection-formula}
Let $H\le G$ have finite index.
\begin{enumerate}[label=(\roman*)]
\item For a representation $W$ of $H$ and $U$ of $G$, $y\otimes(w\otimes u)\mapsto(y\otimes w)\otimes yu$ is a $G$-isomorphism
$\Ind_H^G(W\otimes\Res_HU)\simeq\Ind_H^GW\otimes U$. In particular, if $\chi$ is a character of $G$ trivial on $H$, then $\Ind_H^GW\otimes\chi\simeq\Ind_H^GW$.
\item $\Ind_H^G\mathbf1=k[G/H]$, the permutation representation on the cosets. If $[G:H]=2$ and $\operatorname{char}k\ne2$, then $\Ind_H^G\mathbf1\simeq\mathbf1\oplus\chi_H$,
where $\chi_H\colon G\to\{\pm1\}$ is the character with kernel $H$.
\item $\Ind_K^G\Ind_H^K\simeq\Ind_H^G$ for $H\le K\le G$; $\Ind(W\oplus W')\simeq\Ind W\oplus\Ind W'$; and $(\Ind_H^GW)_K\simeq\Ind_H^G(W_K)$ for a field extension
$K/k$.
\end{enumerate}
\end{proposition}

\begin{proof}
(i) The formula is compatible with $yx\otimes(w\otimes u)=y\otimes(xw\otimes xu)$ ($x\in H$), since both go to $(yx\otimes w)\otimes yxu$; it is $G$-linear.
On the block $s_i\otimes(W\otimes U)$ it is $w\otimes u\mapsto(s_i\otimes w)\otimes s_iu$, a bijection onto $(s_i\otimes W)\otimes U$; so it is bijective. If
$\chi|_H=1$, then $\Res_H\chi=\mathbf1$ and $W\otimes\mathbf1=W$.
(ii) $k[G]\otimes_{k[H]}k$ has basis $s_i\otimes1$, permuted by $G$ as the cosets $s_iH$. For index $2$ the vectors $1\otimes1\pm s\otimes1$ are fixed,
respectively multiplied by $\chi_H$, and they form a basis when $2\ne0$.
(iii) Associativity of the tensor product, $k[G]\otimes_{k[K]}(k[K]\otimes_{k[H]}W)=k[G]\otimes_{k[H]}W$; additivity of $\otimes$; and $K\otimes_k(k[G]\otimes_{k[H]}W)=
K[G]\otimes_{K[H]}W_K$.
\end{proof}

\begin{theorem}[Mackey]\label{thm:reps:mackey}
Let $H\le G$ have finite index and $W$ a representation of $H$.
\begin{enumerate}[label=(\alph*)]
\item For $K\le G$, $\Res_K\Ind_H^GW\simeq\bigoplus_{s\in K\backslash G/H}\Ind^K_{K\cap sHs^{-1}}\Res_{K\cap sHs^{-1}}W^s$, the sum over a set of representatives of the
double cosets.
\item (Mackey's irreducibility criterion.) Let $G$ be finite and $k=\C$. Then $\Ind_H^GW$ is irreducible iff $W$ is
irreducible and, for every $s\in G\setminus H$, the representations $W^s$ and $W$ of $H_s=sHs^{-1}\cap H$ have no common irreducible
constituent.
\item (Index two, any characteristic.) Let $[G:H]=2$, $s\in G\setminus H$, and $W$ absolutely irreducible. Then $\Res_H\Ind_H^GW\simeq W\oplus W^s$, and
$\Ind_H^GW$ is absolutely irreducible iff $W_{\bar k}\not\simeq W^s_{\bar k}$.
\end{enumerate}
\end{theorem}

\begin{proof}
(a) We do not prove this; see \cite[Section 7.3]{SerreLinearReps}. (The summand for the double coset $KsH$ is the sum of the blocks
$t\otimes W$ with $t$ running over representatives of the cosets $tH\subseteq KsH$.)
(c) The restriction is \Cref{prop:reps:induction}(i) with $s_1=1$, $s_2=s$ ($H$ is normal). Work over $\bar k$, which is allowed by
\Cref{prop:reps:projection-formula}(iii); $W$ and $W^s$ are irreducible. Suppose $W\not\simeq W^s$. The $H$-representation $W\oplus W^s$ is semisimple
with non-isomorphic summands, so by \Cref{prop:reps:semisimple}(iii) its subrepresentations are $0$, $W=1\otimes W$, $W^s=s\otimes W$ and the whole
space. A nonzero $G$-stable subspace is $H$-stable and $s$-stable, and $s$ exchanges $1\otimes W$ and $s\otimes W$; so it is everything, and $\Ind W$
is irreducible over $\bar k$. Suppose $W\simeq W^s$. By \Cref{thm:reps:frobenius-reciprocity}(i) and Schur over $\bar k$,
$\End_G(\Ind W)=\Hom_H(W,W\oplus W^s)$ has dimension $2$; an irreducible representation over $\bar k$ has endomorphism ring $\bar k$
(\Cref{thm:reps:schur}(ii)), so $\Ind W$ is reducible.
\end{proof}

\begin{proof}[Sketch of proof of (b)]
By (i) of \Cref{thm:reps:frobenius-reciprocity} and (a) with $K=H$,
\[
\End_G(\Ind W)\simeq\Hom_H(W,\Res_H\Ind W)=\bigoplus_{s\in H\backslash G/H}\Hom_H(W,\Ind^H_{H_s}W^s),
\]
and by Frobenius reciprocity again the summand for $s$ is $\Hom_{H_s}(\Res_{H_s}W,W^s)$. The summand for $s=1$ is $\End_H(W)$. In characteristic $0$ all
representations are semisimple, so $\Ind W$ is irreducible iff $\dim\End_G(\Ind W)=1$, iff $\End_H(W)=\C$ (that is, $W$ irreducible) and the
other summands vanish. See \cite[Section 7.4]{SerreLinearReps}. In characteristic $p$ the criterion fails as stated, because
$\End_G(V)=k$ does not force irreducibility without semisimplicity (\Cref{ex:reps:maschke-fails}(d)); only (c) is needed here.
\end{proof}

\begin{theorem}[Clifford]\label{thm:reps:clifford}
Let $N\trianglelefteq G$ have finite index and $V$ an irreducible representation of $G$ over any field. Then $\Res_NV$ is semisimple, its
irreducible constituents are the conjugates $W^g$ ($g\in G$) of any one of them, and they all occur with the same multiplicity.
\end{theorem}

\begin{proof}
Let $W\subseteq V$ be an irreducible $N$-subrepresentation (one of minimal positive dimension). For $g\in G$, $gW$ is $N$-stable, since
$n\,gW=g(g^{-1}ng)W=gW$, and $w\mapsto gw$ is an isomorphism $W^g\to gW$ of representations of $N$; so $gW$ is irreducible. The sum $\sum_ggW$ is
$G$-stable and nonzero, hence equal to $V$. So $\Res_NV$ is a sum of irreducibles of the form $W^g$, semisimple by \Cref{prop:reps:semisimple}, and
every irreducible constituent is isomorphic to some $W^g$ (\Cref{thm:reps:jordan-holder}). For an irreducible $S$ of $N$, the element $g$ maps the
isotypic component $V_S$ bijectively onto $V_{S^g}$, so $\dim V_S=\dim V_{S^g}$, and the multiplicities agree.
\end{proof}

\begin{theorem}[Clifford theory for a subgroup of index $2$]\label{thm:reps:clifford-index-two}
Let $G$ be any group, $H\le G$ of index $2$, $s\in G\setminus H$, $\chi_H\colon G\to\{\pm1\}$ the character with kernel $H$, $k$ a field with
$\operatorname{char}k\ne2$, and $V$ an absolutely irreducible representation of $G$ over $k$ of dimension $n$. Exactly one of the following holds.
\begin{enumerate}[label=(\Alph*)]
\item $\Res_HV$ is absolutely irreducible. Then $V\not\simeq V\otimes\chi_H$, and $\Ind_H^G\Res_HV\simeq V\oplus(V\otimes\chi_H)$.
\item $\Res_HV$ is not absolutely irreducible. Then $n=2d$, there is an absolutely irreducible representation $W$ of $H$ over $\bar k$ of
dimension $d$ with $W\not\simeq W^s$, $\Res_HV_{\bar k}\simeq W\oplus W^s$ and $V_{\bar k}\simeq\Ind_H^GW$, and $V\simeq V\otimes\chi_H$ over $k$.
\end{enumerate}
Hence $V\simeq V\otimes\chi_H$ iff (B) holds. In that case $\Hom_G(V,V\otimes\chi_H)$ is a line over $k$, spanned by an isomorphism $\phi$ with
$\phi^2\in k^\times$ and $\tr\phi=0$, whose two eigenspaces over $\bar k$ are $W$ and $W^s$. For $n=2$, (B) holds iff $V_{\bar k}\simeq\Ind_H^G\psi$ for a
character $\psi$ of $H$ with $\psi^s\ne\psi$.
\end{theorem}

\begin{proof}
(A) An isomorphism $\phi\colon V\to V\otimes\chi_H$ commutes with $\rho(H)$, so it is a scalar $c$ (\Cref{thm:reps:schur}(iii) for $\Res_HV$). Then
$c\rho(s)=\chi_H(s)\rho(s)c=-c\rho(s)$, so $2c=0$ and $c=0$, a contradiction. The induced statement is
$\Ind_H^G\Res_HV\simeq V\otimes\Ind_H^G\mathbf1\simeq V\otimes(\mathbf1\oplus\chi_H)$ (\Cref{prop:reps:projection-formula}(i), (ii)).

(B) $V_{\bar k}$ is irreducible and $\Res_HV_{\bar k}$ is reducible. By \Cref{thm:reps:clifford}, $\Res_HV_{\bar k}$ is semisimple and its constituents are
$W$ and $W^s$ for an irreducible $W$ (absolutely irreducible, as $\bar k$ is algebraically closed), each of dimension $d$; being reducible, it
has at least two constituents, so $n\ge2d$. By \Cref{thm:reps:frobenius-reciprocity}(i) the inclusion $W\subseteq\Res_HV_{\bar k}$ gives a nonzero $G$-map
$\Ind_H^GW\to V_{\bar k}$, surjective as $V_{\bar k}$ is irreducible; so $n\le2d$. Hence $n=2d$ and the map is an isomorphism $\Ind_H^GW\simeq V_{\bar k}$. As
$\Ind W$ is irreducible, $W\not\simeq W^s$ by \Cref{thm:reps:mackey}(c), which also gives $\Res_HV_{\bar k}\simeq W\oplus W^s$. By
\Cref{prop:reps:projection-formula}(i), $V_{\bar k}\otimes\chi_H\simeq\Ind_H^G(W\otimes\Res_H\chi_H)=\Ind_H^GW\simeq V_{\bar k}$, and since $V_{\bar k}$ is irreducible,
\Cref{prop:reps:scalars}(iii) gives $V\simeq V\otimes\chi_H$ over $k$.

The two cases are exclusive and exhaustive, so $V\simeq V\otimes\chi_H$ iff (B). In case (B), $V\otimes\chi_H$ is absolutely irreducible
(\Cref{prop:reps:twist}(ii)), so $\Hom_G(V,V\otimes\chi_H)$ is a line (\Cref{thm:reps:schur}(iii)) spanned by an isomorphism $\phi$. The composite
$\phi\circ\phi$ is a $G$-isomorphism $V\to V\otimes\chi_H^2=V$, hence a nonzero scalar. By \Cref{lem:reps:trace-zero}, $\tr\phi=0$ ($\chi_H(s)=-1\ne1$). Over
$\bar k$, $\phi$ commutes with $H$, so it preserves the isotypic components $W$ and $W^s$ of $\Res_HV_{\bar k}$ (\Cref{prop:reps:semisimple}(iii)) and acts on
each by a scalar (Schur); the two scalars differ, since otherwise $\phi$ would be a scalar and (A) would hold. For $n=2$, $d=1$ and $W=\psi$ is a
character; conversely if $V_{\bar k}\simeq\Ind\psi$ then $\Res_HV_{\bar k}\simeq\psi\oplus\psi^s$ is reducible.
\end{proof}

\begin{paperlink}
Proposition prop:ns: ``by Frobenius reciprocity $\rhobar\otimes\F_9\simeq\Ind\psi$, so $\rhobar\simeq\rhobar\otimes\chi_{-3}$''. In the language above: $G=G_\Q$, $H=G_{\Q(\sqrt{-3})}$,
$\chi_H=\chi_{-3}$; the eigenline of $\rhobar(H)=C_4$ with character $\psi$ (\Cref{ex:reps:c4}) is a nonzero $H$-map $\psi\to\Res_H(\rhobar\otimes\F_9)$, which
by \Cref{thm:reps:frobenius-reciprocity}(i) gives a nonzero $G$-map $\Ind_H^G\psi\to\rhobar\otimes\F_9$, surjective because $\rhobar\otimes\F_9$ is irreducible, hence
an isomorphism (both have dimension $2$). Then $\Ind\psi\otimes\chi_{-3}\simeq\Ind\psi$ (\Cref{prop:reps:projection-formula}(i)), and the isomorphism
descends to $\F_3$ (\Cref{prop:reps:scalars}(ii)). The last equivalence of Lemma lem:h0 ($\rhobar\simeq\rhobar\otimes\chi_{-3}$ iff $\rhobar$ is not absolutely
irreducible on $\Gal(\Qbar/\Q(\sqrt{-3}))$) and the first sentence of the proof of Proposition prop:Sf (the same with $\chi_{\ell^*}$) are
\Cref{thm:reps:clifford-index-two}; it needs $\rhobar$ absolutely irreducible, which complex conjugation provides
(\Cref{prop:reps:distinct-eigenvalues}). The representation $\rhobar$ is induced from $H$ only after extending scalars: over $\F_3$, $\Res_H\rhobar$ is
irreducible (\Cref{prop:reps:ns3-induced}).
\end{paperlink}

\begin{proposition}[Dihedral representations]\label{prop:reps:dihedral}
Let $\operatorname{char}k\ne2$ and $\rho\colon G\to\GL_2(k)$ absolutely irreducible, with finite projective image $P$, the image of $G$ in $\PGL_2(\bar k)$.
\begin{enumerate}[label=(\roman*)]
\item $P$ is not cyclic.
\item The following are equivalent: (a) $P$ is dihedral of order $2m$ for some $m\ge2$ (for $m=2$, $P\simeq V_4$); (b) $\rho\simeq\rho\otimes\chi$ for a
character $\chi$ of order $2$; (c) $\rho_{\bar k}\simeq\Ind_H^G\psi$ for a subgroup $H$ of index $2$ and a character $\psi$ of $H$.
\item The characters $\chi\ne1$ with $\rho\simeq\rho\otimes\chi$ correspond bijectively, by $\chi\mapsto$ the image of $\ker\chi$ in $P$, to the cyclic
subgroups of index $2$ of $P$: there is exactly one if $P$ is dihedral of order at least $6$, and there are three if $P\simeq V_4$.
\end{enumerate}
\end{proposition}

\begin{proof}
Let $Z=\rho^{-1}(\bar k^\times)$ be the preimage of the scalars, so $P=G/Z$. (i) If $P$ is generated by the image of $g_0$, then $\rho(G)\subseteq\bar k^\times\langle\rho(g_0)\rangle$
is abelian, and $\rho$ is absolutely reducible by \Cref{prop:reps:dim2-absirr}(i).
First a claim: \emph{for $H\le G$ of index $2$, $\Res_H\rho$ is absolutely reducible iff the image of $H$ in $P$ is cyclic.} If the image is cyclic,
$\rho(H)$ is abelian as in (i) and \Cref{prop:reps:dim2-absirr}(i) applies. If $\Res_H\rho_{\bar k}$ has a stable line, then by
\Cref{thm:reps:clifford-index-two}(B) it is $\psi\oplus\psi^s$, so $\rho(H)$ is diagonal in a suitable basis, and its image in $\PGL_2(\bar k)$ is a finite
subgroup of the group of classes of $\mathrm{diag}(x,1)$, which is isomorphic to $\bar k^\times$; finite subgroups of $\bar k^\times$ are cyclic.
(ii) (b)$\Leftrightarrow$(c): a character of order $2$ is $\chi_H$ for $H=\ker\chi$ of index $2$, and \Cref{thm:reps:clifford-index-two} says that
$\rho\simeq\rho\otimes\chi_H$ iff (B) holds for $H$, iff $\rho_{\bar k}$ is induced from a character of $H$. (c)$\Rightarrow$(a): in the basis of
\Cref{prop:reps:induction}(iii), $\rho(H)$ is diagonal and $\rho(G\setminus H)$ antidiagonal. The image $C$ of $H$ in $P$ is cyclic by the claim, of index
$1$ or $2$ in $P$; by (i), $P\ne C$. An antidiagonal matrix $\sm0ab0$ squares to the scalar $ab$, and conjugating $\mathrm{diag}(x,y)$ by it gives
$\mathrm{diag}(y,x)$, which is the inverse of $\mathrm{diag}(x,y)$ modulo scalars. So $P=C\rtimes C_2$ with the involution inverting $C$, and $\#C=m\ge2$ by
(i): $P$ is dihedral of order $2m$ (for $m=2$, $V_4$). (a)$\Rightarrow$(c): a dihedral group of order $2m$ has a cyclic subgroup of index $2$
(for $V_4$, any subgroup of order $2$); its preimage $H$ in $G$ has index $2$ and cyclic image, so $\Res_H\rho$ is absolutely reducible by the
claim, and \Cref{thm:reps:clifford-index-two}(B) gives (c).
(iii) If $\rho\simeq\rho\otimes\chi$ with $\chi\ne1$, then $\chi^2=1$ (\Cref{prop:reps:twist}(iv)), so $\chi=\chi_H$ with $H=\ker\chi$ of index $2$. For $z\in Z$,
$\tr\rho(z)=\chi(z)\tr\rho(z)$ and $\tr\rho(z)=2c\ne0$ for the scalar $c=\rho(z)$; so $\chi(z)=1$ and $Z\subseteq H$. Hence such $H$ correspond to the
subgroups of index $2$ of $P$, and by the claim and \Cref{thm:reps:clifford-index-two} exactly those with cyclic image occur. In a dihedral
group of order $2m$, $m\ge3$, the only cyclic subgroup of index $2$ is the rotation subgroup (the others, which exist for $m$ even, are
dihedral of order $m\ge4$, hence not cyclic); $V_4$ has three subgroups of index $2$, all cyclic.
\end{proof}

\begin{paperlink}
Among irreducible $\rhobar$ (a reducible one fails it trivially), the hypothesis of Theorem thm:main excludes exactly the representations
induced from $\Gal(\Qbar/\Q(\sqrt{\ell^*}))$, not all dihedral ones (\Cref{rem:reps:mod-l-galois}(ii)). At
$\ell=3$, the image $N_s(\F_3)$ has projective image $V_4$ and is induced in three ways, one of them from the determinant-one subgroup
(\Cref{prop:reps:ns3-induced}); the image $N_{ns}(\F_3)$ has projective image $D_4$ and is induced only from $C_{ns}(\F_3)$, whose fixed field is not
$\Q(\sqrt{-3})$ (\Cref{prop:reps:nns3}), so it satisfies the hypothesis.
\end{paperlink}

\section[Twisting, the adjoint and the trace form]{Twisting, the adjoint, the symmetric square and the trace form}\label{sec:reps:adjoint}

This section treats the operations the paper performs on $\rho$: twisting by a character, passing to the adjoint $\End(V)$ and to its
trace-zero part $\ad V$, and the identification of $\ad V$ with a twist of the symmetric square and with its own dual. In the paper,
$A_f=\ad V$ (Section 1.1), $A=\ad T$ and $T^D=\Hom(A,\Zl(1))\simeq A(1)$ (Section 4.2), $H^0(\Q,W^D)$ is a space of trace-zero $G$-maps
$\rhobar\to\rhobar\otimes\bar\chi_\ell$ (Lemma lem:h0), and the proof of Lemma lem:tamfactor computes the adjoint action of two triangular
matrices in the basis $e,h,f$.

\begin{proposition}[Twisting]\label{prop:reps:twist}
Let $V$ be a representation of $G$ over $k$ (or a free $R$-module) of rank $n$ and $\chi\colon G\to k^\times$ (or $R^\times$) a character.
\begin{enumerate}[label=(\roman*)]
\item $\tr(\rho\otimes\chi)(g)=\chi(g)\tr\rho(g)$, $\det(X-(\rho\otimes\chi)(g))=\det(X-\chi(g)\rho(g))$ and $\det(\rho\otimes\chi)=\chi^n\det\rho$.
\item $V$ is irreducible (absolutely irreducible) iff $V\otimes\chi$ is; $\Hom_G(V\otimes\chi,W\otimes\chi)=\Hom_G(V,W)$; $(V\otimes\chi)^\vee=V^\vee\otimes\chi^{-1}$; and
$\Res_{\ker\chi}(V\otimes\chi)=\Res_{\ker\chi}V$.
\item $\End(V\otimes\chi)=\End(V)$ as representations (the same space with the same action), hence $\ad(V\otimes\chi)=\ad V$.
\item If $V\simeq V\otimes\chi$, then $\chi^n=1$ on $G$; for $n=2$, $\chi^2=1$.
\end{enumerate}
\end{proposition}

\begin{proof}
On the space $V$ the twist acts by $\chi(g)\rho(g)$; this gives (i) and the last statement of (ii). The $G$-stable subspaces of $V$ and of
$V\otimes\chi$ coincide (a scalar preserves every subspace), also after extending scalars. A linear map $\phi$ satisfies
$\phi\,\chi(g)\rho_V(g)=\chi(g)\rho_W(g)\phi$ iff $\phi\rho_V(g)=\rho_W(g)\phi$. The dual of $\chi(g)\rho(g)$ is $\chi(g)^{-1}\rho(g^{-1})^t$. (iii)
$(\chi(g)\rho(g))x(\chi(g)\rho(g))^{-1}=\rho(g)x\rho(g)^{-1}$. (iv) Take determinants in $\rho(g)\simeq\chi(g)\rho(g)$: $\det\rho(g)=\chi(g)^n\det\rho(g)$.
\end{proof}

\begin{lemma}[The trace-zero lemma]\label{lem:reps:trace-zero}
Let $\phi\colon V\to V\otimes\chi$ be a $G$-map, that is, $\phi\rho(g)=\chi(g)\rho(g)\phi$ for all $g\in G$. Then $(1-\chi(g))\tr\phi=0$ for every $g\in G$. Hence
$\tr\phi=0$ if $\chi\ne1$ (over a field), or, over a ring $R$, if $1-\chi(g)\in R^\times$ for some $g$. If moreover $V$ is absolutely irreducible,
$\Hom_G(V,V\otimes\chi)$ has dimension at most $1$, and its nonzero elements are isomorphisms.
\end{lemma}

\begin{proof}
$\tr\phi=\tr(\rho(g)^{-1}\phi\rho(g))=\tr(\rho(g)^{-1}\chi(g)\rho(g)\phi)=\chi(g)\tr\phi$. The last statement is \Cref{thm:reps:schur}(iii) for the absolutely
irreducible representations $V$ and $V\otimes\chi$ (\Cref{prop:reps:twist}(ii)).
\end{proof}

\begin{paperlink}
Lemma lem:h0: ``every such map has trace $0$ automatically, since $\tr\phi=\chi_\ell(g)\tr\phi$ for every $g$ and $1-\chi_\ell(c)=2$ is a unit for
complex conjugation $c$'' is \Cref{lem:reps:trace-zero} over $R=\Z/\ell^k$, $\ell$ odd. ``Taking determinants'' (the step $\chi_\ell^2\equiv1$) is
\Cref{prop:reps:twist}(iv) modulo $\ell^k$ (see \Cref{prop:reps:twist-length}(v)). ``$\ad T$ is unchanged by the twist'' (proofs of Theorems thm:main
and thm:quant, Lemma lem:tamfactor at an additive potentially multiplicative prime, Lemma lem:euler for $\kod{I}_n^*$, $\kod{I}_0^*$) is (iii).
\end{paperlink}

\begin{definition}\label{def:reps:adjoint}
The \emph{adjoint representation} of $V$ is $\End(V)=V\otimes V^\vee$ with $g\cdot x=\rho(g)x\rho(g)^{-1}$, also written $\Ad V$ or $\Ad\rho$; the \emph{trace-zero
adjoint} is $\ad V=\{x\in\End(V):\tr x=0\}$, a subrepresentation since the trace is conjugation invariant. The paper writes $A_f=\ad V$ for
$V=V_\ell E$ and $A=\ad T$ for $T=T_\ell E$.
\end{definition}

\begin{proposition}\label{prop:reps:adjoint}
Let $V$ be a representation of $G$ over $k$ (or a free $R$-module) of rank $n$.
\begin{enumerate}[label=(\roman*)]
\item If $n$ is invertible in $k$ (or $R$), $\End(V)=k\cdot1\oplus\ad V$ as representations, $x=\frac{\tr x}n\cdot1+\bigl(x-\frac{\tr x}n\cdot1\bigr)$. For $n=2$ this needs
$\operatorname{char}k\ne2$, that is, $\ell$ odd.
\item If $\rho(g)$ has eigenvalues $\lambda_1,\dots,\lambda_n$ in $\bar k$ (with multiplicity), then $\Ad\rho(g)$ has the eigenvalues $\lambda_i/\lambda_j$ on $\End(V)$. For
$n=2$ it has the eigenvalues $\lambda_1/\lambda_2$, $1$, $\lambda_2/\lambda_1$ on $\ad V$, in every characteristic, and
$\det(X-\Ad\rho(g)\mid\ad V)=(X-1)\bigl(X^2-(\tr(g)^2/\det(g)-2)X+1\bigr)$.
\item In the basis $e,h,f$ of $\ad V$ (convention (C3)):
\begin{align*}
&\Ad\mathrm{diag}(x,y):\quad e\mapsto\tfrac xy\,e,\quad h\mapsto h,\quad f\mapsto\tfrac yx\,f;\\
&t=\Mtwo1{n'}01:\quad (t-1)e=0,\quad (t-1)h=-2n'e,\quad (t-1)f=n'h-n'^2e;\\
&F=\Mtwo{\varepsilon q}{c'}0\varepsilon:\quad FeF^{-1}=qe,\quad FhF^{-1}=h-2c'\varepsilon^{-1}e,\\
&\phantom{F=\Mtwo{\varepsilon q}{c'}0\varepsilon:\quad}FfF^{-1}=q^{-1}f+c'\varepsilon^{-1}q^{-1}h-c'^2\varepsilon^{-2}q^{-1}e,
\end{align*}
for $x,y,\varepsilon\in R^\times$, $q\in R^\times$, $n',c'\in R$. In the paper $\varepsilon=\pm1$ is the value at Frobenius of the unramified character of the Tate data
(convention B-C14), so $\varepsilon^{-1}=\varepsilon$.
\end{enumerate}
\end{proposition}

\begin{proof}
(i) Scalars are fixed by conjugation, and the decomposition is preserved by conjugation since it is defined by the trace.
(ii) Over $\bar k$, conjugate $\rho(g)$ to an upper triangular matrix with diagonal $\lambda_1,\dots,\lambda_n$. On $V^\vee$, $g$ acts by $\rho(g^{-1})^t$, with eigenvalues
$\lambda_j^{-1}$, and the eigenvalues on $V\otimes V^\vee$ are the products (triangularise both factors). For $n=2$ and $g=\sm{\lambda_1}b0{\lambda_2}$, the matrix of
$\Ad g$ in the ordered basis $e,h,f$ is upper triangular with diagonal $\lambda_1/\lambda_2,1,\lambda_2/\lambda_1$: $geg^{-1}=\frac{\lambda_1}{\lambda_2}e$, $ghg^{-1}\in h+ke$ and
$gfg^{-1}\in\frac{\lambda_2}{\lambda_1}f+kh+ke$, as the formulas of (iii) for $F$ show (with $\varepsilon q=\lambda_1$, $\varepsilon=\lambda_2$). The factor $X^2-(\tr(g)^2/\det(g)-2)X+1$
has roots $\lambda_1/\lambda_2$, $\lambda_2/\lambda_1$, since their sum is $(\lambda_1^2+\lambda_2^2)/\lambda_1\lambda_2$.
(iii) Direct multiplication. For instance $th=\sm1{-n'}0{-1}$, so $tht^{-1}=\sm1{-n'}0{-1}\sm1{-n'}01=\sm1{-2n'}0{-1}=h-2n'e$; and
$tf=\sm{n'}010$, $tft^{-1}=\sm{n'}{-n'^2}1{-n'}=f+n'h-n'^2e$. With $F^{-1}=\frac1{\varepsilon^2q}\sm\varepsilon{-c'}0{\varepsilon q}$: $Fe=\sm0{\varepsilon q}00$ and $FeF^{-1}=qe$;
$Fh=\sm{\varepsilon q}{-c'}0{-\varepsilon}$ and $FhF^{-1}=\frac1{\varepsilon^2q}\sm{\varepsilon^2q}{-2\varepsilon qc'}0{-\varepsilon^2q}=h-2c'\varepsilon^{-1}e$; $Ff=\sm{c'}0\varepsilon0$ and
$FfF^{-1}=\frac1{\varepsilon^2q}\sm{\varepsilon c'}{-c'^2}{\varepsilon^2}{-\varepsilon c'}$, which is the stated combination. These identities were also checked symbolically.
\end{proof}

\begin{paperlink}
Section 1.1: $A_f=\ad V$, and at a good prime the Frobenius has the eigenvalues $\alpha/\beta,1,\beta/\alpha$ on $\ad V$ by (ii), for either Frobenius
(the set is stable under inversion; convention B-C3), which gives the good factor of $L(A_f,s)$ (\Cref{ch:lfunctions}). Lemma lem:tamfactor:
``$(t-1)e=0$, $(t-1)h=-2n'e$, $(t-1)f=n'h-n'^2e$, and as $2$ is a unit, $(t-1)M=n'(\Zl e\oplus\Zl h)$'', and ``the adjoint action sends $e$ to $qe$
and $h$ to $h-2c'\varepsilon e$'', are (iii) with $\varepsilon^{-1}=\varepsilon$; $F=\rho(\Frob)$ there is the arithmetic Frobenius (convention B-C2(c)). The
cohomology built on these formulas is in \Cref{ch:cohomology}.
\end{paperlink}

\begin{proposition}[Invariants of twisted adjoints]\label{prop:reps:H0-twisted-adjoint}
Let $V$ be a representation over $k$, or a free module of rank $n$ over a local ring $R$, and $\chi$ a character.
\begin{enumerate}[label=(\roman*)]
\item $(\End(V)\otimes\chi)^G=\Hom_G(V,V\otimes\chi)$.
\item If $n$ is invertible, $(\End(V)\otimes\chi)^G=(\ad V\otimes\chi)^G\oplus R(\chi)^G$, and $(\ad V\otimes\chi)^G=\{\phi\in\Hom_G(V,V\otimes\chi):\tr\phi=0\}$.
\item If $\chi\ne1$ (over a local ring $R$: if the reduction $\bar\chi$ of $\chi$ modulo the maximal ideal is nontrivial), then $R(\chi)^G=0$, every
$G$-map $V\to V\otimes\chi$ has trace $0$, and $(\ad V\otimes\chi)^G=\Hom_G(V,V\otimes\chi)$.
\item If $V$ is absolutely irreducible over $k$ and $n\in k^\times$, then $(\ad V)^G=0$; and for $\chi\ne1$, $(\ad V\otimes\chi)^G\ne0$ iff $V\simeq V\otimes\chi$, in
which case it is a line spanned by an isomorphism.
\end{enumerate}
\end{proposition}

\begin{proof}
(i) is \Cref{prop:reps:homG}. (ii) Twist the decomposition of \Cref{prop:reps:adjoint}(i) by $\chi$ and take invariants. (iii) If $\bar\chi(g)\ne1$ then
$\chi(g)-1$ is a unit, and $\chi(g)x=x$ forces $x=0$; the rest is \Cref{lem:reps:trace-zero} and (ii). (iv) $\End_G(V)=k$, and a scalar $c$ has
trace $nc$, which vanishes only for $c=0$. For $\chi\ne1$, by (iii) $(\ad V\otimes\chi)^G=\Hom_G(V,V\otimes\chi)$, which is described by
\Cref{lem:reps:trace-zero}.
\end{proof}

\begin{paperlink}
Lemma lem:h0, case $k=1$ of the sentence quoted after \Cref{prop:reps:homG}: the $\ell$-torsion of $H^0(\Q,W^D)$ is $H^0(\Q,\ad\rhobar\otimes\bar\chi_\ell)$, the
space of trace-zero $G_\Q$-maps $\rhobar\to\rhobar\otimes\bar\chi_\ell$; by (iii) the trace condition is automatic because $\bar\chi_\ell\ne1$ for $\ell$ odd. Proposition prop:ns: ``$H^0(\Q,\ad\rhobar_{E,3}(1))\ne0$'' is (iv) with $\chi=\bar\chi_3=\chi_{-3}$
and the isomorphism $\rhobar\simeq\rhobar\otimes\chi_{-3}$.
\end{paperlink}

\begin{proposition}[Adjoint and symmetric square]\label{prop:reps:sym2}
Let $V$ be free of rank $2$ over $R$ with basis $e_1,e_2$, $\omega(x,y)=\det[x\,|\,y]$, and let $\Sym V$ be the quotient of $V\otimes V$ by the span of the
$x\otimes y-y\otimes x$, with basis $e_1^2,e_1e_2,e_2^2$. If $\rho(g)$ has eigenvalues $\lambda,\mu$, then $\Sym\rho(g)$ has eigenvalues $\lambda^2,\lambda\mu,\mu^2$.
\begin{enumerate}[label=(\roman*)]
\item The map $\Xi\colon\Sym V\to\End(V)$, $\Xi(xy)(v)=\omega(v,x)y+\omega(v,y)x$, is well defined, takes values in $\ad V$, and satisfies
$\Xi(g\cdot m)=\det(g)\,\rho(g)\Xi(m)\rho(g)^{-1}$. So $\Xi$ is a $G$-map $\Sym V\otimes\det^{-1}\to\ad V$. In coordinates
\[
\Xi(e_1^2)=-2e,\qquad\Xi(e_1e_2)=h,\qquad\Xi(e_2^2)=2f,
\]
so $\Xi$ is an isomorphism iff $2\in R^\times$; over $\Z$ its image has index $4$ in $\ad\Z^2$, and over $\Zl$ it is an isomorphism of lattices iff
$\ell$ is odd.
\item Hence, if $2\in R^\times$, $\ad\rho\simeq\Sym\rho\otimes\det^{-1}$ and $\Sym\rho\simeq\ad\rho\otimes\det$; moreover $V\otimes V\simeq\Sym V\oplus\det$ and
$V\otimes V^\vee\simeq\ad V\oplus\mathbf1$.
\item (Normalisations.) Let $V=V_\ell E$, with $\det=\chi_\ell$ (\Cref{ch:ellcurves}). Then $\Sym V\simeq\ad V(1)$ and $\Sym(V^\vee)\simeq\ad V(-1)$. The
classical symmetric square $L$-function is that of $\Sym(V^\vee)=\Sym H^1$, so with the convention $L_p(M(n),s)=L_p(M,s+n)$ (B-C2(d)),
$L_p(\Sym f,s)=L_p(A_f,s-1)$.
\end{enumerate}
\end{proposition}

\begin{proof}
The eigenvalue statement follows by triangularising $\rho(g)$: in the basis $e_1^2,e_1e_2,e_2^2$ the matrix of $\Sym\rho(g)$ is then triangular with
diagonal $\lambda^2,\lambda\mu,\mu^2$.
(i) The expression is bilinear and symmetric in $(x,y)$, so it factors through $\Sym V$. The trace of $v\mapsto\lambda(v)y$ is $\lambda(y)$, so
$\tr\Xi(xy)=\omega(y,x)+\omega(x,y)=0$. Using $\omega(v,gx)=\det(g)\omega(g^{-1}v,x)$ (\Cref{prop:reps:basic-facts}(iii)),
\[
\Xi(gx\,gy)(v)=\det(g)\bigl(\omega(g^{-1}v,x)gy+\omega(g^{-1}v,y)gx\bigr)=\det(g)\,\rho(g)\Xi(xy)\rho(g)^{-1}(v).
\]
In coordinates, $\omega(v,e_1)=-v_2$ and $\omega(v,e_2)=v_1$ for $v=v_1e_1+v_2e_2$, so $\Xi(e_1^2)(v)=-2v_2e_1$, $\Xi(e_1e_2)(v)=v_1e_1-v_2e_2$ and
$\Xi(e_2^2)(v)=2v_1e_2$, which are $-2e$, $h$, $2f$. The matrix of $\Xi$ in these bases is $\mathrm{diag}(-2,1,2)$, of determinant $-4$.
(ii) The first two statements are (i). For the others, $x\otimes y\mapsto(xy,\,\omega(x,y))$ is a $G$-map $V\otimes V\to\Sym V\oplus\det$, bijective when
$2\in R^\times$ (inverse $(xy,c)\mapsto\frac12(x\otimes y+y\otimes x)+\frac c2(e_1\otimes e_2-e_2\otimes e_1)$); and $V\otimes V^\vee=\End(V)=\ad V\oplus\mathbf1$ by
\Cref{prop:reps:adjoint}(i).
(iii) $\Sym V\simeq\ad V\otimes\det=\ad V(1)$ by (ii). By \Cref{prop:reps:basic-facts}(iii), $V^\vee\simeq V\otimes\chi_\ell^{-1}=V(-1)$, so
$\Sym(V^\vee)\simeq\Sym V(-2)\simeq\ad V(-1)$. Then $L_p(\Sym(V^\vee),s)=L_p(\ad V(-1),s)=L_p(\ad V,s-1)$; the definitions of the local factors are in
\Cref{ch:lfunctions} (convention B-C7).
\end{proof}

\begin{paperlink}
Section 1.1: ``In the normalisation of the symmetric square, $L_p(\Sym f,s)=L_p(A_f,s-1)$'' is (iii); the Introduction's comparison of
Flach's theorem on $\Sym f$ with the adjoint uses (ii). Over $\Z_2$ the map $\Xi$ is not an isomorphism, one of the places where $\ell$
odd is needed for lattice-level statements.
\end{paperlink}

\begin{proposition}[The trace form]\label{prop:reps:trace-form}
Let $T$ be free of rank $n$ over a commutative ring $R$, with an action of $G$.
\begin{enumerate}[label=(\roman*)]
\item $B(x,y)=\tr(xy)$ on $\End_R(T)$ is $G$-invariant, $B(gxg^{-1},gyg^{-1})=B(x,y)$, and perfect: $B(E_{ij},E_{kl})=\delta_{jk}\delta_{il}$, a permutation Gram
matrix, of determinant $\pm1$.
\item For $n=2$ the Gram matrix of $B$ on $\ad T$ in the basis $e,h,f$ is $\begin{psmallmatrix}0&0&1\\0&2&0\\1&0&0\end{psmallmatrix}$, of determinant $-2$, so $B$ is
perfect on $\ad T$ iff $2\in R^\times$.
\item Let $R=\Zl$ with $\ell$ odd and $A=\ad T$. Then $x\mapsto B(x,{-})$ is a $G$-isomorphism $A\simeq A^\vee=\Hom_{\Zl}(A,\Zl)$, and for every character
$\chi\colon G\to\Zl^\times$, $\Hom_{\Zl}(A,\Zl(\chi))\simeq A(\chi)$; the same holds with $\Ql$ or $\Ql/\Zl$ in place of $\Zl$. For $\ell=2$ the image of
$A$ in $A^\vee$ has index $2$.
\end{enumerate}
\end{proposition}

\begin{proof}
(i) $\tr(gxg^{-1}gyg^{-1})=\tr(gxyg^{-1})=\tr(xy)$; and $E_{ij}E_{kl}=\delta_{jk}E_{il}$ has trace $\delta_{jk}\delta_{il}$. (ii) $e^2=f^2=0$, $eh=-e$, $he=e$,
$hf=-f$, so $\tr(e^2)=\tr(f^2)=\tr(eh)=\tr(hf)=0$, while $\tr(ef)=\tr(E_{11})=1$ and $\tr(h^2)=2$. A bilinear form on a free module of finite rank is
perfect iff its Gram determinant is a unit (\Cref{ch:modules}). (iii) By (ii) the map is an isomorphism; it is equivariant because
$B(gx,y)=B(x,g^{-1}y)$ for the conjugation action, which says $B(gx,{-})=g\cdot B(x,{-})$ in $A^\vee$. Twisting by $\chi$ gives
$\Hom(A,\Zl(\chi))=A^\vee(\chi)\simeq A(\chi)$; tensor with $\Ql$ or $\Ql/\Zl$ for the other coefficients. For $\ell=2$ the cokernel of the Gram matrix
has order $\abs{\det}_2^{-1}=2$.
\end{proof}

\begin{paperlink}
Section 4.2: ``$T^D=\Hom(A,\Zl(1))\simeq A(1)$ (via the trace form)'' is (iii) with $\chi=\chi_\ell$; it is one of the places where $\ell$ odd is used.
Then $V^D=A(1)\otimes\Ql$ and $W^D=\ad T(1)\otimes\Ql/\Zl$, the module of Lemma lem:h0 and Lemma lem:dfg21 (\Cref{ch:blochkato}).
\end{paperlink}

\begin{proposition}[Invariance of the adjoint]\label{prop:reps:adjoint-invariance}
Let $V$ be a representation of $G$ over a field and $\chi$ a character of $G$.
\begin{enumerate}[label=(\roman*)]
\item Twisting does not change the adjoint: $\End(V\otimes\chi)$ is $\End(V)$, and $\ad(V\otimes\chi)$ is $\ad V$, by
\Cref{prop:reps:twist}(iii).
\item The transpose $x\mapsto x^t$, where $x^t(\lambda)=\lambda\circ x$, is a $G$-isomorphism $\End(V)\to\End(V^\vee)$ preserving traces, so it restricts to a
$G$-isomorphism $\ad V\simeq\ad(V^\vee)$; the map $x\mapsto-x^t$ has the same properties and also respects commutators.
\item Let $L\subseteq V$ be an $\mathcal{O}$-lattice in a $K$-vector space, and put $\ad L=\End_\mathcal{O}(L)\cap\ad V$. For $c\in K^\times$, $\End_\mathcal{O}(cL)=\End_\mathcal{O}(L)$ and
$\ad(cL)=\ad L$ inside $\End_K(V)$. For $L^\vee=\Hom_\mathcal{O}(L,\mathcal{O})\subseteq V^\vee$, the transpose maps $\End_\mathcal{O}(L)$ onto $\End_\mathcal{O}(L^\vee)$ and $\ad L$ onto
$\ad(L^\vee)$.
\item Consequently, if $L'$ is a stable lattice homothetic to $L$, to $L^\vee$, or to a twist $L\otimes\mathcal{O}(\chi)$, then $\ad L'$ is canonically isomorphic to
$\ad L$ as a representation.
\end{enumerate}
\end{proposition}

\begin{proof}
(ii) On $V^\vee$, $g$ acts by $\rho^\vee(g)=\rho(g^{-1})^t$. Then
\[
(\rho(g)x\rho(g)^{-1})^t=\rho(g)^{-t}x^t\rho(g)^t=\rho^\vee(g)x^t\rho^\vee(g)^{-1},
\]
and $\tr(x^t)=\tr x$.
Also $[x,y]^t=-[x^t,y^t]$, so $-x^t$ respects commutators. (iii) $x(cL)\subseteq cL$ iff $xL\subseteq L$. A $K$-linear $x$ preserves $L$ iff $x^t$ preserves
$L^\vee$, because $L=(L^\vee)^\vee$ (both are free and dual bases correspond). (iv) Combine (i) to (iii); a twist does not change the
underlying module.
\end{proof}

\begin{paperlink}
Section 4.2: the lattice $M_{f,\lambda}$ of DFG04 is identified by $\varphi_E^*$ with $H^1(E,\Zl)=\Hom(T_\ell E,\Zl)\simeq T_\ell E(-1)$ (Proposition prop:pullback;
convention B-C5(b)), ``so its adjoint is $\ad T$'': this is (ii) and (i), with no irreducibility needed. The original version of the paper
said instead that $M_{f,\lambda}$ and $T_\ell E$ are homothetic when $\rhobar$ is irreducible; the correct form of that argument is that
$M_{f,\lambda}$ is homothetic to the dual $T_\ell E^\vee$ (\Cref{thm:reps:homothety}), and then (iv) identifies the adjoints. Theorem thm:quant,
Step 3 (``$A_g=A_f$'' for a twist) is (i).
\end{paperlink}

\section{\texorpdfstring{$\GL_2(\F_3)$ and its subgroups on $\F_3^2$ and $\F_9^2$}{GL2(F3) and its subgroups over F3 and F9}}\label{sec:reps:gl23}

This section applies the theory to the standard representations of the subgroups of $\GL_2(\F_3)$ that occur as mod $3$ images. It is the
representation-theoretic core of the paper's Proposition prop:ns: for the three possible images $\GL_2(\F_3)$, $N_{ns}(\F_3)$, $N_s(\F_3)$ of
an irreducible $\rhobar_{E,3}$, the restriction to the determinant-one subgroup (the image of $\Gal(\Qbar/\Q(\sqrt{-3}))$) is absolutely
irreducible except for $N_s(\F_3)$, where it is $\psi\oplus\psi^{-1}$ over $\F_9$ and the representation is isomorphic to its twist by $\det$
(\Cref{thm:reps:three-images}). The group theory (subgroups, orders, orbits on lines, the three images) is \Cref{ch:groups}'s; only
the representation-theoretic reading is added here. Throughout, $\rho$ is the standard representation (the inclusion $H\subseteq\GL_2(\F_3)$),
and since $\det$ takes values in $\{\pm1\}$, $\rho\otimes\det$ is the twist by the character $\det|_H$.

\begin{proposition}[The standard representations of the subgroups of $\GL_2(\F_3)$]\label{prop:reps:gl23-table}
For the standard models, the standard representation $\rho$ of $H$ on $\F_3^2$ has the following properties. In the first table,
``irr.'' means irreducible over $\F_3$, ``abs.'' absolutely irreducible, and the last column is the image of $H$ in $\PGL_2(\F_3)\simeq S_4$.
\begin{center}
\begin{tabular}{lcccccc}
\toprule
$H$ & order & $\det$ onto $\F_3^\times$ & irr. & abs. & $\#$ of $H/(H\cap\{\pm1\})$ & projective image\\
\midrule
$\GL_2(\F_3)$ & $48$ & yes & yes & yes & $24$ & $S_4$\\
$\SL_2(\F_3)$ & $24$ & no & yes & yes & $12$ & $A_4$\\
$N_{ns}(\F_3)$ & $16$ & yes & yes & yes & $8$ & $D_4$\\
$Q_8$ & $8$ & no & yes & yes & $4$ & $V_4$ (normal in $S_4$)\\
$N_s(\F_3)$ & $8$ & yes & yes & yes & $4$ & $V_4$ (not normal)\\
$C_{ns}(\F_3)$ & $8$ & yes & yes & no & $4$ & $C_4$\\
$C_4=\langle J\rangle$ & $4$ & no & yes & no & $2$ & $C_2$\\
$C_s(\F_3)$ & $4$ & yes & no & no & $2$ & $C_2$\\
$\mathrm B(\F_3)$ & $12$ & yes & no & no & $6$ & $S_3$\\
$\mathrm U$ & $3$ & no & no & no & $3$ & $C_3$\\
\bottomrule
\end{tabular}
\end{center}
In the second table, $\mathrm{Hom}$ is $\Hom_H(\rho,\rho\otimes\det)$, with the dimension of its trace-zero part in parentheses, and the last
column says whether $\rho\simeq\rho\otimes\det$ (``triv.'' when $\det|_H=1$, so that the question is empty).
\begin{center}
\begin{tabular}{lccccl}
\toprule
$H$ & $\dim\End_H(\rho)$ & $\dim(\ad\rho)^H$ & $\dim\mathrm{Hom}$ & $\rho\simeq\rho\otimes\det$ & trace-zero part\\
\midrule
$\GL_2(\F_3)$ & $1$ & $0$ & $0\ (0)$ & no & \\
$\SL_2(\F_3)$ & $1$ & $0$ & $1\ (0)$ & triv. & \\
$N_{ns}(\F_3)$ & $1$ & $0$ & $0\ (0)$ & no & \\
$Q_8$ & $1$ & $0$ & $1\ (0)$ & triv. & \\
$N_s(\F_3)$ & $1$ & $0$ & $1\ (1)$ & \textbf{yes} & $J$\\
$C_{ns}(\F_3)$ & $2$ & $1$ & $0\ (0)$ & no & \\
$C_4$ & $2$ & $1$ & $2\ (1)$ & triv. & $J$\\
$C_s(\F_3)$ & $2$ & $1$ & $2\ (2)$ & yes & $w$, $J=\sigma w$\\
$\mathrm B(\F_3)$ & $1$ & $0$ & $1\ (1)$ & no & $e$ (not invertible)\\
$\mathrm U$ & $2$ & $1$ & $2\ (1)$ & triv. & $e=u-1$\\
\bottomrule
\end{tabular}
\end{center}
\end{proposition}

\begin{proof}
Orders, surjectivity of $\det$ and the stable lines over $\F_3$ (hence the column ``irr.'') are in \Cref{ch:groups}. Absolute irreducibility is
\Cref{prop:reps:dim2-absirr}(iv). The order of $H/(H\cap\{\pm1\})$ is $\#H/2$ when $-1\in H$, which is the case for all rows but $\mathrm U$; the
identifications $\PGL_2(\F_3)\simeq S_4$, $\PSL_2(\F_3)\simeq A_4$ and of the smaller images are in \Cref{ch:groups} (a non-abelian group of order $6$ is $S_3$; $N_{ns}(\F_3)/\{\pm1\}$ is
non-abelian of order $8$ with a cyclic subgroup $C_{ns}/\{\pm1\}$ of order $4$ and the involution $\sigma$ outside it, so it is $D_4$).
For the second table: in the absolutely irreducible rows $\End_H=\F_3$ (\Cref{thm:reps:schur}(iii)); for the cyclic groups $C_{ns}=\langle a\rangle$,
$C_4=\langle J\rangle$, $\mathrm U=\langle u\rangle$ the commutant of the non-scalar generator is $2$-dimensional (\Cref{lem:reps:commutant}); for $C_s$ it is the
diagonal matrices (the commutant of $\sigma$); for $\mathrm B$ it is $\F_3$ (\Cref{ex:reps:maschke-fails}(d)). As $2\in\F_3^\times$,
$\dim(\ad\rho)^H=\dim\End_H-1$ (\Cref{prop:reps:adjoint}(i)). When $\det|_H=1$, $\mathrm{Hom}=\End_H$ and its trace-zero part is $(\ad\rho)^H$. When $\det|_H\ne1$,
every element of $\mathrm{Hom}$ has trace $0$ (\Cref{lem:reps:trace-zero}). For $\GL_2(\F_3)$ and $N_{ns}(\F_3)$, $\mathrm{Hom}=0$ by \Cref{prop:reps:gl23-sl23}
and \Cref{prop:reps:nns3}; for $N_s(\F_3)$ it is $\F_3J$ by \Cref{prop:reps:ns3-induced}. For $C_{ns}$: if $\phi a=-a\phi$ with $\phi\ne0$, then $\ker\phi$ is
$a$-stable, hence $0$ ($C_{ns}$ is irreducible), so $a$ and $-a$ are conjugate; but their eigenvalues are $\{1+i,1-i\}$ and $\{-1-i,-1+i\}$.
For $C_s$ and $\mathrm B$: $\phi$ must anticommute with $\sigma=\mathrm{diag}(1,-1)$ (of determinant $-1$), so $\phi=\sm0xy0$, and the other generator
$\mathrm{diag}(-1,1)$ of $C_s$ imposes nothing new; in $\mathrm B$, commuting with $u$ forces $y=0$. Finally $\rho\simeq\rho\otimes\det$ iff $\mathrm{Hom}$ contains an
invertible element. Every entry was also checked by enumerating $\GL_2(\F_3)$, for instance with the following PARI/GP script, which prints,
for each group, its order, $\dim\End_H$, $\dim(\ad\rho)^H$, $\dim\mathrm{Hom}$ and the dimension of its trace-zero part.
\end{proof}

\begin{lstlisting}[language=GP]
red(M) = lift(Mod(M, 3));
gen(L) = my(S = Set([matid(2)]), T);\
  while(1, T = S; foreach(T, x, foreach(L, g, S = setunion(S, [red(x*g)])));\
    if(#S == #T, return(S)));
mats = concat(vector(81, k, my(d = digits(k-1+81, 3)); [[d[2],d[3];d[4],d[5]]]));
sols(H, tw, t0) = #select(X -> (!t0 || trace(X)%3 == 0) &&\
  #select(g -> red(X*g - matdet(g)^tw*g*X) != 0, H) == 0, mats);
row(H) = [#H, logint(sols(H,0,0),3), logint(sols(H,0,1),3),\
  logint(sols(H,1,0),3), logint(sols(H,1,1),3)];
J = [0,1;-1,0]; s = [1,0;0,-1]; a = [1,-1;1,1]; u = [1,1;0,1];
{
  G = [["Ns",[J,s]], ["Nns",[a,s]], ["Q8",[J,[1,-1;-1,-1]]], ["C4",[J]],
       ["Cns",[a]], ["U",[u]], ["B",[u,s,-matid(2)]]];
  foreach(G, x, print(x[1], ": ", row(gen(x[2]))));
}
\\ Ns: [8, 1, 0, 1, 1]    Nns: [16, 1, 0, 0, 0]   Q8: [8, 1, 0, 1, 0]
\\ C4: [4, 2, 1, 2, 1]    Cns: [8, 2, 1, 0, 0]    U: [3, 2, 1, 2, 1]
\\ B: [12, 1, 0, 1, 1]
\end{lstlisting}

\begin{proposition}[$N_s(\F_3)$: induced over $\F_9$, not over $\F_3$]\label{prop:reps:ns3-induced}
Let $\rho$ be the standard representation of $N_s(\F_3)=\langle J,\sigma\rangle$ and $C_4=\langle J\rangle=N_s(\F_3)\cap\SL_2(\F_3)$, the kernel of $\det$.
\begin{enumerate}[label=(\roman*)]
\item $\rho$ is absolutely irreducible, but $\Res_{C_4}\rho\otimes\F_9\simeq\psi\oplus\psi^{-1}$ with $\psi(J)=i$, the eigenvectors being $v=(1,i)^t$ and $\sigma v=(1,-i)^t$. Since
$\sigma J\sigma^{-1}=J^{-1}$, $\psi^\sigma=\psi^{-1}\ne\psi$.
\item In the basis $(v,\sigma v)$ of $\F_9^2$, $\rho(J)=\mathrm{diag}(i,-i)$ and $\rho(\sigma)=w$. This is the matrix form of $\Ind_{C_4}^{N_s}\psi$ of
\Cref{prop:reps:induction}(iii) (with $s=\sigma$, $\sigma^2=1$), so $\rho\otimes\F_9\simeq\Ind_{C_4}^{N_s}\psi$. Over $\F_3$, $\rho$ is not induced from a character of
$C_4$: $\Res_{C_4}\rho$ is irreducible over $\F_3$.
\item $\rho\simeq\rho\otimes\det$ over $\F_3$, with intertwiner $J$: $J\rho(g)=\det(g)\rho(g)J$ for all $g$. Moreover $\tr J=0$, $J^2=-1$, and the eigenlines of $J$
over $\F_9$ are the two summands of (i), as \Cref{thm:reps:clifford-index-two}(B) predicts.
\item The projective image is $V_4$, and $\rho\simeq\rho\otimes\chi$ for each of the three nontrivial characters $\chi$ of $N_s(\F_3)/\{\pm1\}$: their kernels are $C_4$
($\chi=\det$, intertwiner $J$), $C_s(\F_3)=\{\pm1,\pm\sigma\}$ (intertwiner $\sigma$) and $\{\pm1,\pm w\}$ (intertwiner $w$). So $\rho\otimes\F_9$ is induced from
each of these three subgroups.
\item $(\ad\rho\otimes\det)^{N_s}=\F_3J$, a line, and $(\ad\rho)^{N_s}=0$.
\end{enumerate}
\end{proposition}

\begin{proof}
(i) Absolute irreducibility is \Cref{prop:reps:dim2-absirr}(iv), and the decomposition over $\F_9$ is \Cref{ex:reps:c4}; $\sigma v=(1,-i)^t$ is
the $(-i)$-eigenvector. $\sigma J\sigma^{-1}=\sm0{-1}10=J^{-1}$, so $\psi^\sigma(J)=\psi(J^{-1})=-i$.
(ii) $Jv=iv$ and $J(\sigma v)=-i\,\sigma v$; $\sigma(v)=\sigma v$ and $\sigma(\sigma v)=v$. The matrix of \Cref{prop:reps:induction}(iii) for $\psi$ and $s=\sigma$ is
$J\mapsto\mathrm{diag}(\psi(J),\psi^\sigma(J))=\mathrm{diag}(i,-i)$ and $\sigma\mapsto\sm0{\psi(1)}10=w$. Over $\F_3$, $\Res_{C_4}\rho$ is irreducible (\Cref{ex:reps:c4}), whereas
the restriction of a representation induced from a character of $C_4$ would contain that character as a line.
(iii) It suffices to check the generators: $\det J=1$ and $J$ commutes with $J$; $\det\sigma=-1$ and $J\sigma=\sm0{-1}{-1}0=-\sigma J$. And $\tr J=0$, $J^2=-1$, with
eigenlines $\F_9v$, $\F_9\sigma v$.
(iv) $N_s(\F_3)/\{\pm1\}$ has order $4$ and its three involutions $\bar J,\bar\sigma,\bar w$ commute, so it is $V_4$. The relations $wJ=-Jw$, $\sigma J=-J\sigma$,
$w\sigma=-\sigma w$ show that $x\in\{J,\sigma,w\}$ commutes with $\pm1$ and $\pm x$ and anticommutes with the other four elements, so $x\rho(g)=\chi(g)\rho(g)x$ for the
character $\chi$ with kernel $\{\pm1,\pm x\}$. Induction from each kernel is \Cref{thm:reps:clifford-index-two} (or \Cref{prop:reps:dihedral}(iii)).
(v) is \Cref{prop:reps:H0-twisted-adjoint}(iii), (iv) with $\chi=\det\ne1$ and (iii).
\end{proof}

\begin{proposition}[$N_{ns}(\F_3)$]\label{prop:reps:nns3}
Let $\rho$ be the standard representation of $N_{ns}(\F_3)=\langle a,\sigma\rangle$, with $Q_8=N_{ns}(\F_3)\cap\SL_2(\F_3)$ the kernel of $\det$.
\begin{enumerate}[label=(\roman*)]
\item $\rho$ and $\Res_{Q_8}\rho$ are absolutely irreducible.
\item $\rho\not\simeq\rho\otimes\det$, and $(\ad\rho\otimes\det)^{N_{ns}}=0$.
\item Over $\F_9$, $C_{ns}(\F_3)=\langle a\rangle$ acts on the lines spanned by $(1,-i)^t$ and $(1,i)^t$ by the characters $\theta$ and $\theta^3$, where $\theta(a)=1+i$ has order
$8$; $\theta^\sigma=\theta^3\ne\theta$, and $\rho\otimes\F_9\simeq\Ind_{C_{ns}}^{N_{ns}}\theta$. The projective image is $D_4$, and $\rho\simeq\rho\otimes\chi$ for exactly one
nontrivial character $\chi$, the character $\chi_{ns}$ with kernel $C_{ns}(\F_3)$, with intertwiner $J$.
\end{enumerate}
\end{proposition}

\begin{proof}
(i) is \Cref{prop:reps:dim2-absirr}(iv). (ii) If $\phi$ were an isomorphism, $\phi\rho(g)\phi^{-1}=\det(g)\rho(g)$ would give $\tr\rho(g)=\det(g)\tr\rho(g)$, so
$\tr\rho(g)=0$ whenever $\det g=-1$ (as $2\ne0$). But $\det a=2=-1$ and $\tr a=2\ne0$. The invariants vanish by \Cref{prop:reps:H0-twisted-adjoint}(iv).
(iii) $a\binom1{-i}=\binom{1+i}{1-i}=(1+i)\binom1{-i}$, since $(1+i)(-i)=1-i$; and $a\binom1i=(1-i)\binom1i$, with $1-i=(1+i)^3$. The element $1+i$ has order $8$ in
$\F_9^\times$ ($(1+i)^2=2i=-i$). From $\sigma a\sigma^{-1}=a^3$ (\Cref{ch:groups}), $\theta^\sigma(a)=\theta(a^3)=\theta(a)^3$, and $\theta^3\ne\theta$. So $\Res_{C_{ns}}\rho\otimes\F_9\simeq\theta\oplus\theta^3$ is
reducible, case (B) of \Cref{thm:reps:clifford-index-two} for $H=C_{ns}$, and $\rho\otimes\F_9\simeq\Ind\theta$, $\rho\simeq\rho\otimes\chi_{ns}$. Directly, $J=a^2$ commutes with
$a$ and anticommutes with $\sigma$, so it is an intertwiner $\rho\to\rho\otimes\chi_{ns}$. The projective image $N_{ns}(\F_3)/\{\pm1\}$ is $D_4$ (\Cref{prop:reps:gl23-table}),
of order $8$, so by \Cref{prop:reps:dihedral}(iii) there is exactly one such character.
\end{proof}

\begin{proposition}[$\SL_2(\F_3)$ and $\GL_2(\F_3)$]\label{prop:reps:gl23-sl23}
The standard representations of $\SL_2(\F_3)$ and $\GL_2(\F_3)$ are absolutely irreducible. For $\GL_2(\F_3)$, $\rho\not\simeq\rho\otimes\det$ and
$(\ad\rho\otimes\det)^{\GL_2(\F_3)}=0$. The projective images are $A_4$ and $S_4$, exceptional in Dickson's list (\Cref{thm:reps:dickson}), so neither
representation is induced from a subgroup of index $2$.
\end{proposition}

\begin{proof}
Absolute irreducibility: both contain $Q_8$ (\Cref{prop:reps:dim2-absirr}(iii), (iv)). For $\GL_2(\F_3)$ argue as in \Cref{prop:reps:nns3}(ii) with
$a\in\GL_2(\F_3)$ (or with $\sm0111$, of determinant $-1$ and trace $1$). Neither $A_4$ nor $S_4$ is dihedral, so \Cref{prop:reps:dihedral}(ii)
applies.
\end{proof}

\begin{theorem}[The three images]\label{thm:reps:three-images}
Let $H\le\GL_2(\F_3)$ be conjugate to $\GL_2(\F_3)$, $N_{ns}(\F_3)$ or $N_s(\F_3)$ (by \Cref{ch:groups} these are exactly the subgroups with
surjective determinant, no stable line in $\F_3^2$, and an element of determinant $-1$ and trace $0$). Let $H_1=H\cap\SL_2(\F_3)$, the kernel of
$\det|_H$, of index $2$, and $\rho$ the standard representation. The following are equivalent:
\begin{enumerate}[label=(\alph*)]
\item $H$ is conjugate to $N_s(\F_3)$;
\item $\Res_{H_1}\rho$ is not absolutely irreducible;
\item $\rho\simeq\rho\otimes\det$;
\item $(\ad\rho\otimes\det)^H\ne0$;
\item $\rho\otimes\F_9\simeq\Ind_{H_1}^H\psi$ for a character $\psi$ of $H_1$ of order $4$;
\item $H_1$ is abelian.
\end{enumerate}
When they hold, $(\ad\rho\otimes\det)^H$ is the line spanned by a trace-zero isomorphism $\rho\to\rho\otimes\det$; when they fail, $(\ad\rho\otimes\det)^H=0$.
\end{theorem}

\begin{proof}
All statements are invariant under conjugation, so let $H$ be one of the three models; then $H_1$ is $\SL_2(\F_3)$, $Q_8$ or $C_4$ (\Cref{ch:groups}).
(a)$\Leftrightarrow$(f): of these three only $C_4$ is abelian. (a)$\Leftrightarrow$(b): $\SL_2(\F_3)$ and $Q_8$ act absolutely irreducibly and $C_4$ does
not (\Cref{prop:reps:dim2-absirr}(iv)). (b)$\Leftrightarrow$(c): \Cref{thm:reps:clifford-index-two} with $\chi_H=\det|_H$ (its kernel is $H_1$, of index $2$
because $\det$ is onto $\{\pm1\}$), $\operatorname{char}\F_3\ne2$ and $\rho$ absolutely irreducible. (c)$\Leftrightarrow$(d) and the description of the
invariants: \Cref{prop:reps:H0-twisted-adjoint}(iii), (iv). (a)$\Rightarrow$(e): \Cref{prop:reps:ns3-induced}(ii), $\psi(J)=i$ of order $4$.
(e)$\Rightarrow$(b): $\Res_{H_1}\Ind_{H_1}^H\psi\simeq\psi\oplus\psi^s$ is reducible over $\F_9$.
\end{proof}

\begin{paperlink}
Proposition prop:ns and its proof. As $\det\rhobar_{E,3}=\bar\chi_3=\chi_{-3}$ cuts out $\Q(\sqrt{-3})$ (\Cref{ch:galois}), the image of $\Gal(\Qbar/\Q(\sqrt{-3}))$ is $H_1$,
and \Cref{thm:reps:three-images} gives at once: $\rhobar_{E,3}$ is absolutely irreducible on $\Gal(\Qbar/\Q(\sqrt{-3}))$ unless the image is $N_s(\F_3)$; in that
case $\rhobar\otimes\F_9$ is induced from a character of order $4$, $\rhobar\simeq\rhobar\otimes\chi_{-3}$ and $H^0(\Q,\ad\rhobar(1))\ne0$. The value $h_3(E)=1$ then comes from
Lemma lem:h0 (\Cref{prop:reps:twist-length}). This is also the reading of the hypothesis of Theorem thm:main at $\ell=3$ (Proposition prop:Sf),
and of Remark rem:ns and Section sec:anomaly, which concern exactly the image $\Ns$. That $N_s(\F_3)$ is the only one of the three images
whose determinant-one subgroup is abelian is the orbit dichotomy of \Cref{ch:groups} in another form.
\end{paperlink}

\section{Unipotent elements and cyclic groups of order 2, 3, 4, 6}\label{sec:reps:cyclic}

This section computes the fixed vectors of the groups that occur as images of inertia at a bad prime $q\ne\ell$: unipotent groups and
``minus unipotent'' ones (potentially multiplicative reduction), and cyclic groups of order $2$, $3$, $4$, $6$ in $\SL_2$ (potentially good
reduction), on $V$, on $\ad V$, and modulo $3$. In the paper these computations are Lemma lem:IV, Remark rem:types, Lemma lem:euler, the
table of Section 1 and Proposition prop:ns-structure. Which shape occurs for which Kodaira type is a theorem about elliptic curves
(Serre-Tate, Kraus) proved in \Cref{ch:reduction} and \Cref{ch:galrep}; here it is only recorded in the table of \Cref{prop:reps:cyclic-shapes}.

\begin{convention}\label{conv:reps:sigma}
In this section $\sigma$ denotes an element of order $3$ of $\SL_2$, as in the paper's Lemmas lem:IV and lem:euler, and $\tau$ an element of
order $4$; the matrix $\mathrm{diag}(1,-1)$, called $\sigma$ in \Cref{ch:groups} and in \Cref{sec:reps:gl23}, is written $h$ here (it is the basis vector $h$ of $\ad V$
viewed as a matrix). A Frobenius element is an abstract $F$ satisfying $FgF^{-1}=g^q$ on the inertia image (convention (C8)).
\end{convention}

\begin{proposition}[Unipotent elements]\label{prop:reps:unipotent-invariants}
Let $k$ be a field, $V=k^2$, $u\in\GL_2(k)$ unipotent (characteristic polynomial $(X-1)^2$) with $u\ne1$, $\nu=u-1$, and $\xi\in k^\times$, $\xi\ne1$.
\begin{enumerate}[label=(\roman*)]
\item $\nu\ne0$ and $\nu^2=0$; in a basis $e_1\in\ker\nu$, $e_2\notin\ker\nu$ one has $u=\sm1{n'}01$ with $n'\ne0$. $V^u=\ker\nu=\im\nu$ is a line, and $V^{\xi u}=0$. In
particular ``minus a unipotent'' ($\xi=-1$, $\operatorname{char}k\ne2$) fixes no nonzero vector.
\item If $\operatorname{char}k\ne2$, then $(\ad V)^u=(\ad V)^{\xi u}=k\nu=k\,e$, a line. If $\operatorname{char}k=2$, $(\ad V)^u=ke+kh$. The same holds for any group $H$ acting
through matrices $\xi'u^m$ ($\xi'$ scalar, $m\in\Z$) one of which has $u^m\ne1$.
\item If $F\in\GL_2(k)$ satisfies $FuF^{-1}=u^q$ with $q\in\Z$, then $F$ preserves $(\ad V)^u$ and acts on it by $q$: $F\nu F^{-1}=u^q-1=q\nu$. Hence $F^{-1}$ acts by
$q^{-1}$ when $q\in k^\times$. The same holds for $q\in\Zl$ when $k\supseteq\Ql$ and $u^q$ is defined by continuity, $u^q=1+q\nu$.
\end{enumerate}
\end{proposition}

\begin{proof}
(i) By Cayley-Hamilton $\nu^2=0$, and $\nu\ne0$ as $u\ne1$. Then $\im\nu\subseteq\ker\nu$, both are lines, so they are equal. For $e_2\notin\ker\nu$,
$\nu e_2\in\ker\nu\setminus\{0\}$, say $\nu e_2=n'e_1$. A vector fixed by $\xi u$ would be an eigenvector of $u$ with eigenvalue $\xi^{-1}\ne1$, but $1$ is the only
eigenvalue of $u$.
(ii) Scalars act trivially on $\ad V$, so $(\ad V)^{\xi u}=(\ad V)^u$. By \Cref{prop:reps:adjoint}(iii) with $t=u$, for $x=\alpha e+\beta h+\gamma f$,
$(u-1)x=(-2n'\beta-n'^2\gamma)e+n'\gamma h$, which vanishes iff $\gamma=0$ and $2\beta=0$. For $H$: $u^m=1+m\nu$, and if $u^m\ne1$ it is again a nontrivial
unipotent with $(\ad V)^{u^m}=ke$; every element of $H$ fixes $e$.
(iii) $F\nu F^{-1}=FuF^{-1}-1=u^q-1$, and $u^q=(1+\nu)^q=1+q\nu$ for every $q\in\Z$ (binomial theorem with $\nu^2=0$, and $(1+\nu)^{-1}=1-\nu$). By continuity the
same holds for $q\in\Zl$.
\end{proof}

\begin{lemma}[The commutant of one matrix]\label{lem:reps:commutant}
Let $x\in\End_k(V)$, $n=\dim V$.
\begin{enumerate}[label=(\roman*)]
\item If the characteristic polynomial of $x$ has $n$ distinct roots in $\bar k$, the commutant $\{y:xy=yx\}$ is $k[x]\simeq k[X]/(\det(X-x))$, of
dimension $n$.
\item If $n=2$ and $x$ is not a scalar, the commutant is $k\cdot1\oplus k\cdot x=k[x]$; if moreover $\operatorname{char}k\ne2$, its trace-zero part is the line
spanned by $2x-\tr x$.
\item In (i) and (ii) the commutant is a field iff $\det(X-x)$ is irreducible over $k$.
\end{enumerate}
Examples: if $\sigma^2+\sigma+1=0$ ($\sigma$ of order $3$ in $\SL_2$), the trace-zero line is spanned by $2\sigma+1=\sigma-\sigma^{-1}$, and $(2\sigma+1)^2=-3$; if $\tau^2=-1$, it is $k\tau$.
For $k=\Ql$ and $\sigma$ of order $3$ in $\SL_2(\Ql)$ the commutant $\Ql[\sigma]\simeq\Ql[X]/(X^2+X+1)$ is a field iff $\ell\not\equiv1\pmod3$, and for $\tau$ of order $4$,
$\Ql[\tau]$ is a field iff $\ell\not\equiv1\pmod4$ (the roots of unity of $\Ql$, \Cref{ch:local}).
\end{lemma}

\begin{proof}
(i) The minimal polynomial of $x$ vanishes at all $n$ eigenvalues, so it has degree $n$ and equals $\det(X-x)$; hence $k[x]\simeq k[X]/(\det(X-x))$ has
dimension $n$, and $k[x]$ lies in the commutant. Over $\bar k$, $x$ is diagonal with distinct entries in a suitable basis, and the matrices
commuting with it are the diagonal ones, of dimension $n$. The commutant is the solution space of a linear system with coefficients
in $k$, so its dimension over $k$ is also $n$ (as in \Cref{prop:reps:scalars}(i)); it therefore equals $k[x]$.
(ii) Choose $v$ that is not an eigenvector of $x$ (the eigenvectors lie in at most two lines, or in one line if $x$ has a single eigenvalue
and is not scalar). Then $v,xv$ is a basis. If $y$ commutes with $x$, write $yv=\alpha v+\beta xv$; then $y$ and $\alpha+\beta x$ agree on $v$ and on $xv$
($y\,xv=x\,yv=(\alpha+\beta x)xv$), so $y=\alpha+\beta x$. The trace of $\alpha+\beta x$ is $2\alpha+\beta\tr x$, which vanishes iff $\alpha+\beta x=\frac\beta2(2x-\tr x)$.
(iii) $k[X]/(P)$ is a field iff $P$ is irreducible. For the examples: $\tr\sigma=-1$, so $2\sigma-\tr\sigma=2\sigma+1$, $\sigma^{-1}=\sigma^2=-1-\sigma$, and
$(2\sigma+1)^2=4(\sigma^2+\sigma)+1=-3$; $\tr\tau=0$.
\end{proof}

\begin{proposition}[Cyclic groups of order $2,3,4,6$]\label{prop:reps:cyclic-shapes}
Let $K$ be a field, $V=K^2$, and $g\in\SL_2(K)$ of finite order $m\in\{2,3,4,6\}$ with $\operatorname{char}K\nmid m$.
\begin{enumerate}[label=(\roman*)]
\item If $m=2$ then $g=-1$. If $m\ge3$, $g$ has the distinct eigenvalues $\zeta,\zeta^{-1}$ with $\zeta$ a primitive $m$-th root of unity, and characteristic
polynomial $\Phi_m$: $X^2+X+1$, $X^2+1$, $X^2-X+1$ for $m=3,4,6$. For $m=6$, $\sigma=-g$ has order $3$.
\item $V^g=0$. If $m=2$, $(\ad V)^g=\ad V$. If $m\ge3$, $(\ad V)^g$ is the line spanned by $g-g^{-1}$, which is $2g+1$ for $m=3$, $2g$ for $m=4$, and
$2g-1=-(2\sigma+1)$ for $m=6$.
\item Let $F\in\GL_2(K)$ with $FgF^{-1}=g^r$, $r\in\Z$ prime to $m$. If $m\ge3$, then $F$ preserves $(\ad V)^g$ and acts on it by $+1$ if $r\equiv1\pmod m$ and by
$-1$ if $r\equiv-1\pmod m$, and every $r$ prime to $m$ is $\equiv\pm1\pmod m$. For $r=q$ a prime not dividing $m$ the sign is $\chi_{-3}(q)$ if
$m\in\{3,6\}$ and $\chi_{-4}(q)$ if $m=4$; it is the same for $F^{-1}$.
\end{enumerate}
\end{proposition}

\begin{proof}
(i) $X^m-1$ is separable since $\operatorname{char}K\nmid m$, so $g$ is diagonalisable over $\bar K$ with eigenvalues $m$-th roots of unity of product
$\det g=1$. If the two eigenvalues are equal, to $\zeta$ say, then $\zeta^2=1$ and $g=\zeta=\pm1$, of order at most $2$; this is the case $m=2$, $g=-1$. Otherwise
they are $\zeta\ne\zeta^{-1}$, the order of $g$ is that of $\zeta$, so $\zeta$ is a primitive $m$-th root, and $(X-\zeta)(X-\zeta^{-1})=\Phi_m$ because $\Phi_m$ has degree $2$
for these $m$. For $m=6$, $g^3$ has the eigenvalues $\zeta^{\pm3}=-1$, so $g^3=-1$ and $(-g)^3=1$, with $-g\ne1$.
(ii) $1$ is not an eigenvalue. Conjugation by $-1$ is trivial. For $m\ge3$, \Cref{lem:reps:commutant}(ii) gives the line $K(2g-\tr g)$; and $g-g^{-1}$
has trace $0$ and lies in $K[g]$, and is nonzero since $g^2\ne1$. The three expressions follow from $g^{-1}=-1-g$, $-g$, $1-g$ respectively.
(iii) $F$ normalises $\langle g\rangle$, so it preserves the invariants, and $F(g-g^{-1})F^{-1}=g^r-g^{-r}$. The units modulo $3$, $4$, $6$ are $\pm1$, and
$g^{\pm1}-g^{\mp1}=\pm(g-g^{-1})$. Finally $q\equiv\chi_{-3}(q)\pmod3$ for a prime $q\ne3$ (including $q=2$, where both sides are $-1$ modulo $3$) and $q\equiv\chi_{-4}(q)\pmod4$ for an odd
prime $q$ (\Cref{ch:rings}); and for $m=6$ a prime $q\ge5$ is $\equiv\pm1\pmod6$ with the sign of $q\bmod3$. $F^{-1}$ acts by the inverse sign, which is the same.
\end{proof}

The following table collects the fixed spaces for the shapes of the image of inertia $I=I_q$ at a prime $q\ne\ell$, $q\ge5$, on $V=V_\ell E$
(dimension $2$ over $\Ql$). The first column is \Cref{ch:reduction} and \Cref{ch:galrep}'s (Serre-Tate, Kraus), the other columns are
\Cref{prop:reps:unipotent-invariants}, \Cref{prop:reps:cyclic-shapes} and \Cref{prop:reps:reduction-cyclic}. In the unipotent rows the image of $I$ is
infinite, topologically generated, up to the finite part, by $t=\sm1{n'}01$; $\xi$ is a ramified quadratic character. The column ``$F$''
gives the action of an arithmetic Frobenius $F$ ($FgF^{-1}=g^q$) on $(\ad V)^I$; the geometric Frobenius $F^{-1}$ gives the inverse, which matters
only in the unipotent rows ($q$ against $q^{-1}$; convention B-C3(d)). The last column is $\dim(T/3T)^I$ for $T=T_3E$.
\begin{center}
\small
\begin{tabular}{llccccc}
\toprule
image of $I$ & types & $\dim V^I$ & $\dim(\ad V)^I$ & generator & $F$ on $(\ad V)^I$ & $\dim(T/3T)^I$\\
\midrule
trivial & $\kod{I}_0$ & $2$ & $3$ & $\ad V$ & $\alpha/\beta,1,\beta/\alpha$ & $2$\\
$\{\pm1\}$ & $\kod{I}_0^*$ & $0$ & $3$ & $\ad V$ & as for the twist & $0$\\
unipotent, $t\ne1$ & $\kod{I}_n$ & $1$ & $1$ & $e$ & $q$ & $1$ or $2$\\
$\xi\cdot$unipotent & $\kod{I}_n^*$ & $0$ & $1$ & $e$ & $q$ & $0$\\
order $3$ & $\kod{IV}$, $\kod{IV}^*$ & $0$ & $1$ & $2\sigma+1$ & $\chi_{-3}(q)$ & $1$\\
order $4$ & $\kod{III}$, $\kod{III}^*$ & $0$ & $1$ & $\tau$ & $\chi_{-4}(q)$ & $0$\\
order $6$, $-\sigma$ & $\kod{II}$, $\kod{II}^*$ & $0$ & $1$ & $2\sigma+1$ & $\chi_{-3}(q)$ & $0$\\
\bottomrule
\end{tabular}
\end{center}
In the row $\kod{I}_n$, $\dim(T/3T)^I$ is $1$ if $3\nmid n'$ and $2$ if $3\mid n'$; for the Tate curve $v_3(n')=v_3(n)$ (\Cref{ch:reduction}). In the row $\kod{I}_0^*$ the
eigenvalues of $F$ on $\ad V$ are those of the unramified twist $V'$ (\Cref{prop:reps:twist}(iii)). The degree of $L_q(A_f,s)^{-1}$ is $\dim(\ad V)^I$, that is,
$3,3,1,1,1,1,1$, and $\delta_q=\dim V^I$ is $2,0,1,0,0,0,0$ (Definition def:terms). The row $\kod{I}_n$ holds at every $q\ne\ell$, the rows $\kod{I}_0^*$, $\kod{I}_n^*$ at every
odd $q\ne\ell$; for the finite images the tame relation $FgF^{-1}=g^q$ needs the image to be tame, which holds for $q\ge5$ (\Cref{ch:local}).

\begin{proposition}[Reduction of finite cyclic subgroups of $\SL_2(\Zl)$]\label{prop:reps:reduction-cyclic}
Let $\ell$ be a prime, $T=\Zl^2$, $g\in\SL_2(\Zl)$ of finite order $m\in\{2,3,4,6\}$, and $\bar g$ its reduction modulo $\ell$.
\begin{enumerate}[label=(\roman*)]
\item If $\ell\nmid m$, then $\bar g$ has order $m$, and the invariants of $\langle g\rangle$ on $T$ and on $\ad T$ commute with reduction:
$(T/\ell^jT)^g=T^g/\ell^jT^g$ and $(\ad T/\ell^j)^g=(\ad T)^g/\ell^j$. In particular $(T/\ell T)^g=0$.
\item Let $\ell=3$, $m=3$. Then $\bar g$ is a nontrivial unipotent element of $\SL_2(\F_3)$, so $(T/3T)^g$ is a line while $T^g=0$: invariants do not
commute with reduction. On the adjoint they do: $(\ad T)^g=\Z_3(2g+1)$ and $2g+1\not\equiv0\pmod 3$ spans $(\ad T/3)^g$.
\item Let $\ell=3$, $m=6$. Then $\bar g=-\bar u$ with $\bar u$ a nontrivial unipotent, and $(T/3T)^g=0$.
\item Let $\ell=3$, $m\in\{2,4\}$. Then $(T/3T)^g=0$; for $m=4$, $\bar g$ has the irreducible characteristic polynomial $X^2+1$ over $\F_3$ and eigenvalues
$\pm i\in\F_9$.
\item Let $\ell$ be odd and $m\ge3$. Then $\bar g$ is not a scalar, and the commutant of $g$ in $\End_{\Zl}(T)$ is $\Zl[g]=\Zl\oplus\Zl g\simeq\Zl[X]/(\Phi_m)$. Moreover $T$
is free of rank one over $\Zl[g]$; for $\ell=3$, $m=3$, this gives $T/3T\simeq\Z_3[\zeta_3]/(\pi^2)$, $\pi=1-\zeta_3$, with $(T/3T)^g=\pi T/3T$.
\end{enumerate}
\end{proposition}

\begin{proof}
First, the kernel $1+\ell\mathrm M_2(\Zl)$ of reduction contains no element $\gamma\ne1$ of finite order prime to $\ell$, and for $\ell$ odd no element of order $\ell$.
Write $\gamma=1+\ell^ay$ with $a\ge1$ and $y\not\equiv0\pmod\ell$. For $r$ prime to $\ell$, $\gamma^r\equiv1+r\ell^ay\pmod{\ell^{a+1}}$, so $\gamma^r\ne1$. For $\ell$ odd,
$\gamma^\ell=1+\ell^{a+1}y+\binom\ell2\ell^{2a}y^2+\dots\equiv1+\ell^{a+1}y\pmod{\ell^{a+2}}$, because $\ell\mid\binom\ell j$ for $0<j<\ell$, $2a+1\ge a+2$ and $\ell a\ge a+2$ (using
$\ell\ge3$); so $\gamma^\ell\ne1$ (this is also in \Cref{ch:groups}). A nontrivial finite group contains an element of prime order, so the finite
subgroups of the kernel are trivial when $\ell$ is odd, and the finite subgroups of order prime to $\ell$ are trivial for every $\ell$. Hence
reduction modulo $\ell$ is injective on $\langle g\rangle$ if $\ell$ is odd or $\ell\nmid m$.
(i) $\bar g$ has order $m$ by the above, and $\#\langle g\rangle=m\in\Zl^\times$; apply \Cref{thm:reps:maschke}(ii) to the lattices $T$ and $\ad T$. $T^g=0$ by
\Cref{prop:reps:cyclic-shapes}(ii).
(ii) $\bar g\ne1$ by injectivity, and its characteristic polynomial is $X^2+X+1\equiv(X-1)^2\pmod3$; so $\bar g$ is a nontrivial unipotent and
$(T/3T)^g$ is a line (\Cref{prop:reps:unipotent-invariants}(i)), while $T^g=0$. On the adjoint, $(\ad T)^g=\ad T\cap\Q_3(2g+1)$ by
\Cref{prop:reps:cyclic-shapes}(ii). If $2g+1\equiv0\pmod3$ then $g\equiv-2^{-1}\cdot1=1\pmod 3$, which is false; so $2g+1$ is not divisible by $3$ in $\End(T)$, and
$\ad T\cap\Q_3(2g+1)=\Z_3(2g+1)$. Its reduction $2\bar g+1=-(\bar g-1)$ spans $(\ad T/3)^g=\F_3(\bar g-1)$ (\Cref{prop:reps:unipotent-invariants}(ii)).
(iii) $g=-\sigma$ with $\sigma=-g$ of order $3$; by (ii) $\bar\sigma$ is a nontrivial unipotent, and $-\bar\sigma$ fixes no nonzero vector
(\Cref{prop:reps:unipotent-invariants}(i)).
(iv) is (i), as $3\nmid m$; for $m=4$ the characteristic polynomial is $\Phi_4=X^2+1$, without roots in $\F_3$.
(v) $\bar g$ is not a scalar: if $\ell\nmid m$, $\bar g$ has order $m\ge3$ while a scalar of determinant $1$ is $\pm1$; if $\ell=3$, $\bar g$ is $\pm$ a nontrivial
unipotent. So $1$ and $g$ are linearly independent modulo $\ell$, $\Zl\oplus\Zl g$ is a direct summand of $\End_{\Zl}(T)$, and the commutant
$\End_{\Zl}(T)\cap(\Ql\oplus\Ql g)$ (\Cref{lem:reps:commutant}(ii) over $\Ql$) is $\Zl\oplus\Zl g$. The freeness of $T$ over $\Zl[g]$ and the description of
$T/3T$ are proved in \Cref{ch:ant}; we do not prove them here. (The fixed space $\pi T/3T$ is the kernel of $g-1$, which acts as multiplication by $\zeta_3-1=-\pi$.)
\end{proof}

\begin{paperlink}
Lemma lem:IV (types $\kod{IV}$, $\kod{IV}^*$ at $q\ne3$): ``its image in $\Aut(E[3])$ is a nontrivial unipotent subgroup. Consequently $E[3]^{I_q}$ is a
line'' is \Cref{prop:reps:reduction-cyclic}(ii), and ``reduction modulo $3$ is injective on it, the kernel of $\GL_2(\Z_3)\to\GL_2(\F_3)$ being
torsion-free'' is the first paragraph of its proof; the route through $T_3E$ free of rank one over $\Z_3[X]/(X^2+X+1)$, in the unnumbered paragraph after the lemma in the paper, is (v). ``The trace-zero
endomorphisms commuting with $\sigma$ form the line spanned by $\sigma-\sigma^{-1}=2\sigma+1$'' is \Cref{lem:reps:commutant}, and ``$\Frob_q$, and therefore also
$\Frob_q^{-1}$, acts on the line by $\chi_{-3}(q)$'' is \Cref{prop:reps:cyclic-shapes}(iii). In the exclusion of potentially multiplicative reduction,
``$\rhobar_{E,3}(I_q)$ would contain $-1$, which fixes no line'' is trivial: the scalar $-1$ fixes no nonzero vector of $E[3]$, so $E[3]^{I_q}$ would be $0$. Remark rem:types (``$E[3]^{I_q}=0$''
for $\kod{I}_0^*$, $\kod{III}$, $\kod{III}^*$, $\kod{II}$, $\kod{II}^*$, $\kod{I}_n^*$) is (iii), (iv) and \Cref{prop:reps:unipotent-invariants}(i). Lemma lem:euler is the
table above: types $\kod{II}$, $\kod{IV}$ and their duals use $2\sigma+1$ inside $\End_{I_q}(T)=\Zl[\sigma]$, which is not a field for $\ell\equiv1\pmod3$ ((v) and
\Cref{lem:reps:commutant}); types $\kod{III}$, $\kod{III}^*$ use the line $\Ql\tau$, with no need to diagonalise $\tau$ over $\Ql$; type $\kod{I}_n$ uses
\Cref{prop:reps:unipotent-invariants}(iii) with the geometric Frobenius, eigenvalue $q^{-1}$. Proposition prop:ns-structure (ii): ``the trace-zero
endomorphisms commuting with $\tau$ form the line spanned by $\tau$'' and ``the sign can be read modulo $3$'' are \Cref{lem:reps:commutant} and the
injectivity of reduction on $\langle\tau\rangle$ (order $4$, prime to $3$).
\end{paperlink}

\section{Representations over \texorpdfstring{$\Zl$}{Zl}: lattices, reduction, homothety}\label{sec:reps:lattices}

This section passes from representations over $\Ql$ to representations over $\Zl$ and modulo $\ell^m$. In the paper $T=T_\ell E$ is a
$G_\Q$-stable lattice in $V=V_\ell E$, $\rhobar$ is its reduction, $\rho\bmod\ell^k$ its reduction modulo $\ell^k$, $A=\ad T$ is a lattice in $A_f$, and the
term $h_\ell(E)$ of Lemma lem:h0 is ``the largest $k$ such that $\rho\bmod\ell^k\simeq(\rho\otimes\chi_\ell)\bmod\ell^k$''. Everything is proved for an abstract
group and a discrete valuation ring $\mathcal{O}$ (convention (C7)).

\begin{definition}\label{def:reps:lattice}
Let $V$ be a representation of $G$ over $K=\Frac\mathcal{O}$. A \emph{lattice} in $V$ is a finitely generated $\mathcal{O}$-submodule $L$ spanning $V$; it is
\emph{$G$-stable} if $\rho(g)L=L$ for all $g$. Then $\rho_L\colon G\to\GL(L)\simeq\GL_n(\mathcal{O})$, conjugate in $\GL_n(K)$ to $\rho$, and $\rhobar_L$ is the representation of
$G$ on $L/\varpi L$ over the residue field; $\rho_L\bmod\varpi^m$ is the representation on $L/\varpi^mL$. Two lattices are \emph{homothetic} if $L'=cL$ for some
$c\in K^\times$, equivalently $L'=\varpi^mL$ for some $m\in\Z$.
\end{definition}

\begin{proposition}[Stable lattices]\label{prop:reps:lattice}
\begin{enumerate}[label=(\roman*)]
\item For any two lattices $L,L'$ there are $a,b\ge0$ with $\varpi^aL\subseteq L'\subseteq\varpi^{-b}L$.
\item If $G$ is finite, or more generally if the stabiliser of a lattice $L_0$ has finite index in $G$, then $L=\sum_g\rho(g)L_0$ (a finite sum over coset
representatives) is a $G$-stable lattice.
\item If $L$ is stable, so are $\varpi^mL$, $L^\vee=\Hom_\mathcal{O}(L,\mathcal{O})\subseteq V^\vee$, $L\otimes\mathcal{O}(\chi)$ for a character $\chi\colon G\to\mathcal{O}^\times$, $\End_\mathcal{O}(L)\subseteq\End_K(V)$,
and $\ad L=\End_\mathcal{O}(L)\cap\ad V$. Moreover $\End_\mathcal{O}(L)=\ad L\oplus\mathcal{O}E_{11}$ as $\mathcal{O}$-modules, so $\ad L/\varpi^m\ad L=\ad(L/\varpi^mL)$, the trace-zero endomorphisms
of $L/\varpi^mL$.
\end{enumerate}
\end{proposition}

\begin{proof}
(i) Each of the finitely many generators of $L'$ is a $K$-combination of a basis of $L$, so lies in $\varpi^{-b}L$ for $b$ large; exchange the roles for
$a$. (ii) A finite sum of lattices is finitely generated and spans $V$, and it is permuted by $G$. (iii) $(g\lambda)(L)=\lambda(g^{-1}L)=\lambda(L)\subseteq\mathcal{O}$;
$\chi(g)$ is a unit; $\rho(g)x\rho(g)^{-1}$ preserves $L$ if $x$ does; the trace is invariant. The trace $\End_\mathcal{O}(L)\to\mathcal{O}$ is onto ($\tr E_{11}=1$) and split by
$c\mapsto cE_{11}$, so its kernel $\ad L$ is a direct summand, and the split exact sequence $0\to\ad L\to\End_\mathcal{O}(L)\to\mathcal{O}\to0$ stays exact modulo $\varpi^m$, with
$\End_\mathcal{O}(L)/\varpi^m=\End(L/\varpi^mL)$ since $L$ is free.
\end{proof}

\begin{proposition}[Reduction and semisimplification]\label{prop:reps:reduction-ss}
Let $L$ be a $G$-stable lattice in $V$.
\begin{enumerate}[label=(\roman*)]
\item $\det(X-\rhobar_L(g))$ is the reduction of $\det(X-\rho(g))\in\mathcal{O}[X]$; it does not depend on $L$.
\item Hence (Brauer-Nesbitt, \Cref{thm:reps:brauer-nesbitt}) $\rhobar_L^{ss}$ does not depend on $L$; it is written $\rhobar^{ss}$.
\item If $\rhobar_L$ is irreducible for one $L$, then $\rhobar_{L'}\simeq\rhobar_L$ for every stable lattice $L'$.
\end{enumerate}
\end{proposition}

\begin{proof}
(i) The characteristic polynomial of $\rho_L(g)\in\GL_n(\mathcal{O})$ is that of the conjugate $\rho(g)$, with coefficients in $\mathcal{O}$, and reduction commutes with
taking characteristic polynomials. (ii) is immediate. (iii) $\rhobar_{L'}^{ss}\simeq\rhobar_L$ is irreducible of full dimension, so $\rhobar_{L'}$ has one composition factor
and $\rhobar_{L'}\simeq\rhobar_L$ (\Cref{cor:reps:bn-consequences}(ii) with $k'=k$).
\end{proof}

\begin{theorem}[The homothety lemma]\label{thm:reps:homothety}
If $\rhobar_L$ is irreducible, every $G$-stable lattice $L'\subseteq V$ is homothetic to $L$: $L'=\varpi^mL$ for a unique $m\in\Z$.
\end{theorem}

\begin{proof}
By \Cref{prop:reps:lattice}(i) the set of $m\in\Z$ with $L'\subseteq\varpi^mL$ is nonempty and bounded above (as $L'\ne0$ and $\bigcap_m\varpi^mL=0$); let $m$ be its
maximum and replace $L'$ by $\varpi^{-m}L'$. Then $L'\subseteq L$ and $L'\not\subseteq\varpi L$, so the image of $L'$ in $L/\varpi L$ is a nonzero $G$-stable subspace,
hence all of $L/\varpi L$ since $\rhobar_L$ is irreducible. Thus $L'+\varpi L=L$, and $L'=L$ by Nakayama's lemma (\Cref{ch:modules}). Uniqueness: $\varpi^mL=L$
with $m\ne0$ would give $L=\varpi L$, hence $L=0$ by Nakayama.
\end{proof}

\begin{paperlink}
Section 4.2. The original version of the paper argued that ``the lattice $M_{f,\lambda}$ of DFG04 and $T_\ell E$ are homothetic when $\rhobar$ is
irreducible, so their adjoints coincide''. As $M_{f,\lambda}\subseteq H^1$ is a lattice in the dual $V_\ell E^\vee$, the correct application of
\Cref{thm:reps:homothety} is to $M_{f,\lambda}$ and $T_\ell E^\vee$ inside $V_\ell E^\vee$ (both stable, and $T_\ell E^\vee/\ell\simeq\rhobar^\vee$ is irreducible), after which
\Cref{prop:reps:adjoint-invariance}(iv) identifies $\ad M_{f,\lambda}$ with $\ad T$. The revised paper avoids the homothety lemma altogether: Proposition
prop:pullback identifies $M_{f,\lambda}$ with $H^1(E,\Zl)$ itself. The homothety lemma is still the reason why, when $E[\ell]$ is irreducible, every curve
in the isogeny class has the same $T_\ell$ up to homothety, hence the same $\ad T$ (Section 1, ``at such an $\ell$ the choice of the optimal curve does
not matter''; \Cref{ch:ellcurves}).
\end{paperlink}

\begin{example}[Non-homothetic lattices]\label{ex:reps:nonhomothetic}
Let $G\subseteq\GL_2(\Zl)$ consist of upper triangular matrices and contain $t=\sm1101$. Then $L=\Zl e_1\oplus\Zl e_2$ and $L'=\Zl e_1\oplus\ell\Zl e_2$ are $G$-stable and
not homothetic ($L'$ has index $\ell$ in $L$, while $\varpi^mL$ has index $\ell^{2m}$). In the basis $e_1,\ell e_2$ of $L'$ the matrix of $g=\sm ab0d$ is $\sm a{\ell b}0d$, so
$\rhobar_{L'}$ is the direct sum of the two characters $g\mapsto\bar a$, $g\mapsto\bar d$, while $\rhobar_L$ is the non-split extension of the same two characters
(\Cref{prop:reps:nonsplit-line}: $\bar t$ is a nontrivial unipotent, so $\F_\ell e_1$ is the only stable line). Same semisimplification
(\Cref{prop:reps:reduction-ss}), non-isomorphic reductions. For elliptic curves, the Tate modules of two $\ell$-isogenous curves give such a pair
(\Cref{ch:ellcurves}, \Cref{ch:galrep}); this is why the Eisenstein case of the paper (Section sec:eis, Conjecture conj:eis) behaves
differently.
\end{example}

\begin{theorem}[Ribet's lemma]\label{thm:reps:ribet}
Let $K$ be a finite extension of $\Ql$ with ring of integers $\mathcal{O}$, $G$ a compact group and $\rho\colon G\to\GL_2(K)$ a continuous irreducible
representation such that $\rhobar^{ss}\simeq\eta_1\oplus\eta_2$ for characters $\eta_1,\eta_2$. Then there is a $G$-stable lattice $L$ whose reduction $\rhobar_L$ is a non-split
extension with sub-character $\eta_1$ (and one with sub-character $\eta_2$).
\end{theorem}

\begin{proof}
We do not prove this; see \cite[Section 2]{Ribet76}. (The idea: among the stable lattices, which up to homothety form a bounded convex subset (a finite subtree) of
the tree of $\PGL_2(K)$ because $\rho$ is irreducible, an extremal one has a reduction with the required non-split shape.) In the literature
``Ribet's lemma'' usually means this statement; the name is sometimes also attached to \Cref{thm:reps:homothety}, which concerns the opposite
situation of an irreducible reduction.
\end{proof}

\begin{proposition}[Twisted self-isomorphisms modulo $\varpi^m$]\label{prop:reps:twist-length}
Let $G$ be any group, $T$ free of rank $n$ over $\mathcal{O}$ with a $G$-action such that $\bar T=T/\varpi T$ is absolutely irreducible, and $\chi\colon G\to\mathcal{O}^\times$ a
character, with reduction $\bar\chi$. For $m\ge1$ put $M_m=\Hom_G\bigl(T/\varpi^m,(T/\varpi^m)(\chi)\bigr)$.
\begin{enumerate}[label=(\roman*)]
\item For $\phi\in\End(T/\varpi^m)$, $\varpi^{m-1}\phi\ne0$ iff $\phi(T/\varpi^m)\not\subseteq\varpi(T/\varpi^m)$. The $\varpi$-torsion of $M_m$ is isomorphic to $\Hom_G(\bar T,\bar T(\bar\chi))$, of
dimension at most $1$, so $M_m\simeq\mathcal{O}/\varpi^j$ for some $j\le m$.
\item An element $\phi\in M_m$ with $\varpi^{m-1}\phi\ne0$ is an isomorphism $T/\varpi^m\simeq(T/\varpi^m)(\chi)$. Conversely an isomorphism has $\varpi^{m-1}\phi\ne0$, so it
generates $M_m\simeq\mathcal{O}/\varpi^m$.
\item Let $h=\sup\{m\ge0:T/\varpi^m\simeq(T/\varpi^m)(\chi)\}\in\{0,1,\dots,\infty\}$ and $X=\End_\mathcal{O}(T)(\chi)\otimes_\mathcal{O}K/\mathcal{O}$. Then $X[\varpi^m]=\End(T/\varpi^m)(\chi)$,
$X^G[\varpi^m]=M_m\simeq\mathcal{O}/\varpi^{\min(m,h)}$, and $X^G$ has length $h$; if $h<\infty$, $X^G\simeq\mathcal{O}/\varpi^h$. If $h=\infty$, then $T\simeq T(\chi)$ (no
completeness of $\mathcal{O}$ is needed).
\item If $n\in\mathcal{O}^\times$ and $\bar\chi\ne1$, the same holds for $\ad T(\chi)\otimes K/\mathcal{O}$: $(\ad T(\chi)\otimes K/\mathcal{O})^G=X^G$.
\item If $T/\varpi^m\simeq(T/\varpi^m)(\chi)$ then $\chi^n\equiv1\pmod{\varpi^m}$ on $G$. In particular, if $n=2$, $\mathcal{O}=\Zl$ with $\ell$ odd and $\chi(G)=\Zl^\times$, then $h=0$
for $\ell\ge5$ and $h\le1$ for $\ell=3$.
\end{enumerate}
\end{proposition}

\begin{proof}
(i) $T/\varpi^m$ is free over $\mathcal{O}/\varpi^m$, so its $\varpi^{m-1}$-torsion is $\varpi(T/\varpi^m)$; this gives the first statement. If $\varpi\phi=0$, then $\phi$ kills
$\varpi(T/\varpi^m)$ and takes values in the $\varpi$-torsion $\varpi^{m-1}T/\varpi^mT$, which multiplication by $\varpi^{m-1}$ identifies with $\bar T$ (with the twist by $\bar\chi$
for the target); so $M_m[\varpi]\simeq\Hom_G(\bar T,\bar T(\bar\chi))$, of dimension at most one by \Cref{lem:reps:trace-zero} ($\bar T(\bar\chi)$ is absolutely irreducible).
$M_m$ is a finitely generated $\mathcal{O}/\varpi^m$-module whose $\varpi$-torsion has length at most one, hence cyclic (\Cref{ch:modules}).
(ii) By (i), $\phi\bmod\varpi$ is a nonzero $G$-map $\bar T\to\bar T(\bar\chi)$, an isomorphism by Schur. So $\phi$ is surjective modulo $\varpi$, hence surjective by
Nakayama, hence bijective, the two modules having the same finite length. Conversely, an isomorphism is nonzero modulo $\varpi$.
(iii) $X[\varpi^m]=\End(T)(\chi)\otimes\varpi^{-m}\mathcal{O}/\mathcal{O}\simeq\End(T/\varpi^m)(\chi)$, and taking invariants gives $M_m$. Let $M_m\simeq\mathcal{O}/\varpi^j$ with generator $\phi$. If
$j=m$, $\phi$ is an isomorphism by (ii), so $m\le h$. In general $\varpi^j\phi=0$ means that $\phi$ takes values in $\varpi^{m-j}(T/\varpi^m)$, so $\phi=\varpi^{m-j}\phi'$ for a
$G$-map $\phi'\colon T/\varpi^j\to(T/\varpi^j)(\chi)$ (defined by dividing the values by $\varpi^{m-j}$, well defined modulo $\varpi^j$), with $\varpi^{j-1}\phi'\ne0$ because
$\varpi^{j-1}\phi\ne0$. By (ii) $\phi'$ is an isomorphism, so $j\le h$. Conversely, for $m\le h$ the reduction of an isomorphism at a level $\ge m$ is an
isomorphism at level $m$, which generates $M_m\simeq\mathcal{O}/\varpi^m$ by (ii); and for $m>h$ with $h$ finite, if $\psi$ is an isomorphism at level $h$ and $\tilde\psi\in\End_\mathcal{O}(T)$
lifts it, then $x\mapsto\varpi^{m-h}\tilde\psi(x)\bmod\varpi^m$ is an element of $M_m$ (it is $G$-equivariant because $\tilde\psi g-\chi(g)g\tilde\psi$ has values in
$\varpi^hT$) of exact order $\varpi^h$. Hence $j=\min(m,h)$, and $X^G=\bigcup_mX^G[\varpi^m]$ has length $h$, and equals $X^G[\varpi^h]\simeq\mathcal{O}/\varpi^h$ when $h$ is finite.

Now let $h=\infty$. \emph{Compatible isomorphisms.} The reduction map $M_{m+1}\to M_m$ is surjective: $M_{m+1}$ contains an isomorphism, whose
reduction is an isomorphism at level $m$ and so generates $M_m$ by (ii). Choose an isomorphism $\phi_1\in M_1$ and, inductively, $\phi_{m+1}\in M_{m+1}$
reducing to $\phi_m$. Each $\phi_m$ reduces modulo $\varpi$ to $\phi_1\ne0$, so $\varpi^{m-1}\phi_m\ne0$ by (i), and $\phi_m$ is an isomorphism by (ii).
\emph{Passage to the completion.} Let $\widehat{\mathcal{O}}=\varprojlim\mathcal{O}/\varpi^m$ be the $\varpi$-adic completion, a discrete valuation ring with
uniformiser $\varpi$ and $\widehat{\mathcal{O}}/\varpi^m=\mathcal{O}/\varpi^m$, containing $\mathcal{O}$, with fraction field $\widehat K\supseteq K$ (for $\mathcal{O}=\Zl$ nothing
changes). In a basis of $T$, the matrices of the $\phi_m$ are compatible, so they define a matrix $\hat\phi\in\mathrm M_n(\widehat{\mathcal{O}})$ with
$\hat\phi\equiv\phi_m\pmod{\varpi^m}$ for all $m$. For $g\in G$, every entry of $\hat\phi\rho(g)-\chi(g)\rho(g)\hat\phi$ lies in $\varpi^m\widehat{\mathcal{O}}$ for all $m$, hence is $0$;
and $\hat\phi\ne0$ since $\hat\phi\equiv\phi_1\pmod\varpi$. So $\Hom_{\widehat K[G]}(V_{\widehat K},V_{\widehat K}(\chi))\ne0$, where $V=T\otimes_{\mathcal{O}}K$.
\emph{Descent.} By \Cref{prop:reps:scalars}(i) for the field extension $\widehat K/K$, $\Hom_{K[G]}(V,V(\chi))\ne0$. Let $x$ be a nonzero element. Multiplying
$x$ by a power of $\varpi$, we may assume that $x(T)\subseteq T$ and $x(T)\not\subseteq\varpi T$ (the entries of its matrix lie in $\mathcal{O}$, not all in
$\varpi\mathcal{O}$). Then $\bar x=x\bmod\varpi$ is a nonzero $G$-map $\bar T\to\bar T(\bar\chi)$ between irreducible representations, an isomorphism by
\Cref{thm:reps:schur}(i). So $\det x\in\mathcal{O}$ has nonzero reduction, $\det x\in\mathcal{O}^\times$, and $x$ is an isomorphism $T\simeq T(\chi)$ of
$\mathcal{O}[G]$-modules.
(iv) $\End(T)=\ad T\oplus\mathcal{O}$ (\Cref{prop:reps:adjoint}(i)), and $(K/\mathcal{O}(\chi))^G=0$ because $\chi(g)-1$ is a unit for some $g$; so the $\mathcal{O}$-summand contributes
nothing.
(v) Determinants (\Cref{prop:reps:twist}(iv) modulo $\varpi^m$). For the last statement, $\chi^2\equiv1\pmod{\ell^m}$ on $G$ means $x^2\equiv1\pmod{\ell^m}$ for every
$x\in\Zl^\times$; with $x=2$, $\ell^m\mid3$, so $m=0$ for $\ell\ge5$ and $m\le1$ for $\ell=3$ (\Cref{ch:rings}).
\end{proof}

\begin{paperlink}
Lemma lem:h0 with $G=G_\Q$, $\mathcal{O}=\Zl$, $T=T_\ell E$, $\chi=\chi_\ell$: $W^D=\ad T(1)\otimes\Ql/\Zl$, and by (iv) and (iii), $H^0(\Q,W^D)$ has length
$h_\ell(E)$, ``the largest $k$ such that $\rho\bmod\ell^k\simeq(\rho\otimes\chi_\ell)\bmod\ell^k$'', and is cyclic when this is finite. Absolute irreducibility of $\rhobar$ comes from complex conjugation
(\Cref{prop:reps:distinct-eigenvalues}), $\bar\chi_\ell\ne1$ for $\ell$ odd, and $\chi_\ell(G_\Q)=\Zl^\times$ (\Cref{ch:galois}); so (v) gives $h_\ell(E)=0$ for $\ell\ge5$ and
$h_3(E)\le1$ (in particular $H^0(\Q,W^D)$ is finite and cyclic), with equality iff $\rhobar\simeq\rhobar\otimes\bar\chi_3=\rhobar\otimes\chi_{-3}$, for every $E$ with $\rhobar$ irreducible, with or without complex multiplication.
In Lemma lem:h0 the phrase ``a submodule of $(\Ql/\Zl)^3$ with $\ell$-torsion of rank at most one, is cyclic'' should be read as in (iii): such a
module is finite cyclic or isomorphic to $\Ql/\Zl$, and the determinant bound (v) excludes the second case. Remark rem:ns and Section
sec:anomaly use $h_3(E)=1$ for the image $\Ns$ (\Cref{thm:reps:three-images}).
\end{paperlink}

\section{The mod \texorpdfstring{$\ell$}{l} world: a foretaste, and Dickson's list}\label{sec:reps:modl}

This last section reads the chapter's results in the setting of the paper: a continuous representation
$\rhobar\colon G_\Q\to\GL_2(\F_\ell)$ has finite image, so everything above applies to it, and the Galois input needed is collected in
\Cref{ch:galois} and \Cref{ch:galrep}. It also records Dickson's classification, which organises the possible images.

\begin{remark}[The mod $\ell$ dictionary]\label{rem:reps:mod-l-galois}
Let $\rhobar\colon G_\Q\to\GL_2(\F_\ell)$ be continuous, $\ell$ odd, with $\det\rhobar=\bar\chi_\ell$ (as for $\rhobar_{E,\ell}$, \Cref{ch:ellcurves}).
\begin{enumerate}[label=(\roman*)]
\item Complex conjugation $c$ has $\rhobar(c)^2=1$ and $\det\rhobar(c)=\bar\chi_\ell(c)=-1$, so $\rhobar$ is irreducible iff it is absolutely irreducible
(\Cref{prop:reps:distinct-eigenvalues}).
\item The hypothesis of Theorem thm:main concerns $\Res_H\rhobar$ for the subgroup $H=G_{\Q(\sqrt{\ell^*})}$ of index $2$, with $\chi_H=\chi_{\ell^*}=\bar\chi_\ell^{(\ell-1)/2}$. If
$\rhobar$ is reducible, so is $\Res_H\rhobar$, and the hypothesis fails. If $\rhobar$ is irreducible, it is absolutely irreducible by (i), and
\Cref{thm:reps:clifford-index-two} applies: the hypothesis fails iff $\rhobar\simeq\rhobar\otimes\chi_{\ell^*}$, iff $\rhobar\otimes\Fbar_\ell$ is induced from a character of
$H$. Irreducibility cannot be dropped here: for $\ell\ge5$, $\rhobar=\mathbf1\oplus\bar\chi_\ell$ has reducible restriction to $H$, but
$\rhobar\otimes\chi_{\ell^*}=\chi_{\ell^*}\oplus\bar\chi_\ell^{(\ell+1)/2}\not\simeq\rhobar$ (its constituents are not $\mathbf1$, as $\chi_{\ell^*}\ne1$ and $\bar\chi_\ell^{(\ell+1)/2}\ne1$ for
$\ell\ge5$), and $\rhobar\otimes\Fbar_\ell$ is not induced from $H$ (\Cref{prop:reps:projection-formula}: an induced representation is
isomorphic to its twist by $\chi_H$).
\item Lemma lem:h0 concerns $\rhobar\simeq\rhobar\otimes\bar\chi_\ell$, which forces $\bar\chi_\ell^2=1$ (\Cref{prop:reps:twist}(iv)), that is $\ell=3$; there $\bar\chi_3=\chi_{-3}=\det\rhobar$
and \Cref{thm:reps:three-images} applies. For $\ell\ge5$ the two conditions differ: for $\rhobar=\rhobar_{E,\ell}$, $\rhobar\simeq\rhobar\otimes\chi_{\ell^*}$ can occur (only when $\ell\mid N$, by
Proposition prop:Sf), while $\rhobar\simeq\rhobar\otimes\bar\chi_\ell$ never does.
\item The objects of the paper and where they are built:
\end{enumerate}
\begin{center}
\small
\begin{tabular}{p{0.3\textwidth}p{0.33\textwidth}p{0.29\textwidth}}
\toprule
object & meaning & here\\
\midrule
$V=V_\ell E$, $\rho$ & two-dimensional $\Ql$-representation of $G_\Q$ & \Cref{def:reps:representation}\\
$T=T_\ell E$ & a $G_\Q$-stable lattice in $V$ & \Cref{def:reps:lattice}\\
$\rhobar$, $\rho\bmod\ell^k$ & $T/\ell T$, $T/\ell^kT$ & \Cref{prop:reps:reduction-ss}, \Cref{prop:reps:twist-length}\\
$A_f=\ad V$ & trace-zero adjoint & \Cref{def:reps:adjoint}, \Cref{prop:reps:adjoint}\\
$A=\ad T$ & adjoint lattice, a direct summand of $\End(T)$ & \Cref{prop:reps:lattice}(iii)\\
$T^D=\Hom(A,\Zl(1))\simeq A(1)$ & Kummer dual, via the trace form & \Cref{prop:reps:trace-form}\\
$W=A\otimes\Ql/\Zl$, $W^D=A(1)\otimes\Ql/\Zl$ & discrete coefficients & \Cref{prop:reps:twist-length}(iii), (iv)\\
$H^0(\Q,W^D)$, $h_\ell(E)$ & twisted self-isomorphisms of $\rho\bmod\ell^k$ & \Cref{prop:reps:twist-length}\\
$\ad M_{f,\lambda}=\ad T$ & adjoint of a dual, twisted lattice & \Cref{prop:reps:adjoint-invariance}\\
\bottomrule
\end{tabular}
\end{center}
\end{remark}

\begin{theorem}[Dickson]\label{thm:reps:dickson}
Let $\ell$ be an odd prime and $G\le\GL_2(\F_\ell)$, with image $P$ in $\PGL_2(\F_\ell)$. If $\ell\mid\#G$, then $G\supseteq\SL_2(\F_\ell)$ or $G$ lies in a Borel
subgroup. If $\ell\nmid\#G$, then $P$ is cyclic (and $G$ lies in a Cartan subgroup), or dihedral (and $G$ lies in the normaliser of a Cartan
subgroup, not in the Cartan subgroup), or $P$ is isomorphic to $A_4$, $S_4$ or $A_5$.
\end{theorem}

\begin{proof}
We do not prove this; see \cite[Section 2]{Serre72} and \Cref{ch:groups}, where the statement is the same and the case $\ell=3$ is checked
by enumeration.
\end{proof}

For an absolutely irreducible $\rhobar$ with projective image $P$, \Cref{prop:reps:dihedral} reads Dickson's list as follows: $P$ is never cyclic;
$P$ is dihedral iff $\rhobar\otimes\Fbar_\ell$ is induced from a subgroup of index $2$; the exceptional images $A_4$, $S_4$, $A_5$ and the large images
($P\supseteq\PSL_2(\F_\ell)$) are not induced (for $\ell\ge5$, $\PSL_2(\F_\ell)$ is simple and non-abelian, so it is not contained in a dihedral group, whose
subgroups are cyclic or dihedral). At $\ell=3$: $\GL_2(\F_3)\mapsto S_4$, $N_{ns}(\F_3)\mapsto D_4$, $N_s(\F_3)\mapsto V_4$ (\Cref{prop:reps:gl23-table}). The subgroup
$N_s(\F_3)$ is not a maximal subgroup of $\GL_2(\F_3)$: it has index $2$ in $N_{ns}(\F_3)$ (\Cref{ch:groups}).

\section{Exercises}\label{sec:reps:exercises}

\begin{exercise}\label{ex:reps:ex1}
Show that $C_4=\langle J\rangle$ acting on $\F_3^2$ is irreducible but not absolutely irreducible, that its commutant $\F_3[J]$ is a field with $9$ elements,
and decompose $\F_9^2$ into eigenlines of $J$. Which eigenline is spanned by a vector with coordinates in $\F_3$?
\end{exercise}

\begin{exercise}\label{ex:reps:ex2}
Prove directly (without \Cref{thm:reps:schur}) that every absolutely irreducible representation of a finite abelian group is one-dimensional.
Show that $C_{ns}(\F_3)=\langle a\rangle\simeq C_8$ gives an irreducible two-dimensional representation of $C_8$ over $\F_3$, and determine all irreducible
representations of $C_8$ over $\F_3$ (their degrees and their number).
\end{exercise}

\begin{exercise}\label{ex:reps:ex3}
Show that the standard representation of $\mathrm U=\langle u\rangle$ on $\F_3^2$ is indecomposable but not irreducible, and that $\F_p[C_p]$ is uniserial.
Why does this not contradict Maschke's theorem? Show that every representation of $C_p$ over $\F_p$ is a direct sum of the quotients
$\F_p[x]/(x-1)^j$, $1\le j\le p$.
\end{exercise}

\begin{exercise}\label{ex:reps:ex4}
Give two non-isomorphic semisimple representations of $C_2$ over $\F_3$ with the same trace function. Show that two two-dimensional
representations with the same determinant and the same trace function have isomorphic semisimplifications, in every characteristic.
\end{exercise}

\begin{exercise}\label{ex:reps:ex5}
Check the claim of Proposition prop:ns that the determinant-one subgroup $C_4$ of $N_s(\F_3)$ acts on $\F_9^2$ as $\psi\oplus\psi^{-1}$ with $\psi$ of order $4$, and
that the standard representation of $N_s(\F_3)$ over $\F_9$ is isomorphic to $\Ind_{C_4}^{N_s}\psi$: compute the matrices of $J$ and $\sigma$ in the basis
$v=(1,i)^t$, $\sigma v$. Show that over $\F_3$ it is not induced from a character of $C_4$, and that it is nevertheless isomorphic over $\F_3$ to its twist by $\det$.
\end{exercise}

\begin{exercise}\label{ex:reps:ex6}
Check the claim of Proposition prop:ns that $Q_8$ and $\SL_2(\F_3)$ act absolutely irreducibly on $\F_3^2$ in two ways: by Maschke's theorem
(\Cref{prop:reps:dim2-absirr}) and by the commutator argument (a subgroup of a Borel subgroup of $\GL_2(\Fbar_3)$ has abelian commutator subgroup of
exponent dividing $3$). Explain why ``non-abelian and faithful'' does not suffice, using $\mathrm B(\F_3)$.
\end{exercise}

\begin{exercise}\label{ex:reps:ex7}
Show that the standard representation of $N_s(\F_3)$ is isomorphic to its twist by each of the three nontrivial characters of $N_s(\F_3)/\{\pm1\}$, with
intertwiners $J$, $\sigma$, $w$, and that these intertwiners have trace $0$. Show that the standard representation of $N_{ns}(\F_3)$ is induced, over $\F_9$,
from $C_{ns}(\F_3)$, but is not isomorphic to its twist by $\det$. Which of the two satisfies the hypothesis ``absolutely irreducible on the
determinant-one subgroup''?
\end{exercise}

\begin{exercise}\label{ex:reps:ex8}
Check the claim of Lemma lem:h0 that an isomorphism $\phi\colon\rhobar\to\rhobar\otimes\chi_{-3}$ has trace $0$. More generally, show that
$(\ad V\otimes\chi)^G=\Hom_G(V,V\otimes\chi)$ for every nontrivial character $\chi$ when $\dim V$ is invertible in $k$, and give an example with $\chi=1$ where the
two spaces differ.
\end{exercise}

\begin{exercise}\label{ex:reps:ex9}
Let $H\le G$ have index $2$, $\operatorname{char}k\ne2$, and $V$ absolutely irreducible of dimension $2$. Prove that $V\simeq V\otimes\chi_H$ iff $\tr\rho(g)=0$ for all
$g\notin H$. (For ``if'', compare traces and determinants of $V$ and $V\otimes\chi_H$ and use \Cref{cor:reps:bn-consequences}.) Deduce again that the
standard representation of $\GL_2(\F_3)$ is not isomorphic to its twist by $\det$, and that that of $N_s(\F_3)$ is.
\end{exercise}

\begin{exercise}\label{ex:reps:ex10}
Prove $\ad V\simeq\Sym V\otimes\det^{-1}$ for $\dim V=2$ by the explicit map $\Xi$ of \Cref{prop:reps:sym2}, and show that over $\Z$ the image of $\Sym\Z^2$ has index
$4$ in $\ad\Z^2$. Deduce the eigenvalues $\alpha/\beta$, $1$, $\beta/\alpha$ of $\Ad g$ on $\ad V$ from the eigenvalues $\alpha^2,\alpha\beta,\beta^2$ of $\Sym g$.
\end{exercise}

\begin{exercise}\label{ex:reps:ex11}
Compute the Gram matrix of the trace form on $\ad\Zl^2$ in the basis $e,h,f$, and check the identification $T^D=\Hom(A,\Zl(1))\simeq A(1)$ of the paper's
Section 4.2 for $\ell$ odd. What goes wrong for $\ell=2$? Is the trace form on all of $\End(\Z_2^2)$ perfect?
\end{exercise}

\begin{exercise}\label{ex:reps:ex12}
Check the formulas in the proof of Lemma lem:tamfactor: for $t=1+n'e$, $(t-1)e=0$, $(t-1)h=-2n'e$, $(t-1)f=n'h-n'^2e$; and for
$F=\sm{\varepsilon q}{c'}0\varepsilon$ with $\varepsilon=\pm1$, $FeF^{-1}=qe$ and $FhF^{-1}=h-2c'\varepsilon e$. Deduce that for $\ell$ odd $(t-1)\ad T=n'(\Zl e\oplus\Zl h)$.
\end{exercise}

\begin{exercise}\label{ex:reps:ex13}
For $g=-1$, $g=\sm0{-1}1{-1}$ (order $3$), $g=J$ (order $4$), $g=\sm1{-1}10$ (order $6$), $g=u$ and $g=-u$ in $\SL_2(\Z)$, compute $\dim V^g$ and $\dim(\ad V)^g$ for
$V=\Q^2$ and for $V=\F_3^2$, and a generator of $(\ad V)^g$. Check the claims of Lemma lem:IV and Remark rem:types about $E[3]^{I_q}$.
\end{exercise}

\begin{exercise}\label{ex:reps:ex14}
Let $u\ne1$ be unipotent and $FuF^{-1}=u^q$. Show that $F$ acts on $(\ad V)^u$ by $q$. Explain why the factor $1-q^{-1-s}$ of $L_q(A_f,s)^{-1}$ at a
multiplicative prime requires the geometric Frobenius $F^{-1}$, and why at the primes of types $\kod{II}$, $\kod{III}$, $\kod{IV}$ and their duals the choice of
Frobenius does not matter.
\end{exercise}

\begin{exercise}\label{ex:reps:ex15}
Show that for a finite group of order prime to $\ell$ acting on a $\Zl$-lattice $T$, invariants commute with reduction modulo $\ell^m$. Show that this
fails for $g\in\SL_2(\Z_3)$ of order $3$ acting on $T=\Z_3^2$ (Lemma lem:IV), although it holds for the invariants of $g$ on $\ad T$. Does it hold for
$g$ of order $6$ on $T$? On $\ad T$?
\end{exercise}

\begin{exercise}\label{ex:reps:ex16}
Let $G\subseteq\GL_2(\Zl)$ be the closed subgroup generated by $t=\sm1101$. Show that the lattices $L_k=\Zl e_1\oplus\ell^k\Zl e_2$, $k\ge0$, are $G$-stable and pairwise
non-homothetic, that their reductions have the same semisimplification, and that the reduction of $L_k$ is split iff $k\ge1$. Explain why this
does not contradict \Cref{thm:reps:homothety}, and prove that theorem again in your own words.
\end{exercise}

\begin{exercise}\label{ex:reps:ex17}
Let $T$ be free of rank $2$ over $\Zl$, $\ell$ odd, with an action of a group $G$ such that $T/\ell T$ is absolutely irreducible, and let $\chi\colon G\to\Zl^\times$ be
surjective. Show that $\Hom_G(T/\ell^m,(T/\ell^m)(\chi))$ is cyclic for every $m$, that $H^0(G,\ad T(\chi)\otimes\Ql/\Zl)$ is finite of length at most $1$, and that
it is $0$ unless $\ell=3$ (the corrected Lemma lem:h0). For $\ell=3$, show that the length is $1$ iff $T/3T\simeq(T/3T)(\bar\chi)$.
\end{exercise}

\section{Notes and references}\label{sec:reps:notes}

The representation theory of finite groups over $\C$ in \Cref{sec:reps:field} to \Cref{sec:reps:characters} follows \cite{SerreLinearReps}: Maschke's
theorem is in its Section 1.3, Schur's lemma in Section 2.2, the orthogonality relations in Sections 2.3 to 2.5, the divisibility of
degrees in Section 6.5, induction, Frobenius reciprocity and Mackey's criterion in Sections 7.1 to 7.4, Clifford's theorem in Section 8.1,
and the Frobenius-Schur indicator in Section 13.2. These section numbers are recalled, not checked against a printed copy, and should be
read as pointers; no theorem number of \cite{SerreLinearReps} is relied upon. The proofs of Schur's lemma, Maschke's theorem and
Frobenius reciprocity given here are characteristic-free, as is the averaging proof of $\dim\Hom_G(V,W)=\langle\tr_W,\tr_V\rangle$. The treatment of
representations over arbitrary fields (absolute irreducibility, extension of scalars, the density theorem, the theorem of Brauer and
Nesbitt, the Noether-Deuring theorem) is classical and can be found in \cite{CurtisReiner}, Sections 27 to 30 and 29 for Noether-Deuring,
and in \cite[Chapter XVII]{Lang}, where semisimple modules and the density theorem are treated; again these are unverified pointers.
The proof of Brauer-Nesbitt in \Cref{sec:reps:brauer-nesbitt} is the standard one by reduction to an algebraically closed field, cancellation, and
the Frobenius $f\mapsto f^p$; the reduction step uses only the density theorem, so no perfectness of $k$ is needed. The criterion
\Cref{prop:reps:absirr}(iv), that an irreducible representation with commutant $k$ is absolutely irreducible over every field, is a
consequence of the same density theorem. Riley-Hobson-Bence \cite{RHB06} treats character tables over $\C$ in its chapter on representation
theory (Chapter 29 of the third edition; an unverified pointer); \cite[Chapter 18]{DummitFoote} is an algebraic introduction at the
same level.

Clifford theory for a subgroup of index $2$ (\Cref{thm:reps:clifford-index-two}) and the characterisation of dihedral representations
(\Cref{prop:reps:dihedral}) are standard; for the Galois representations of elliptic curves they go back to \cite[Section 2]{Serre72}, which
also contains the classification of subgroups of $\GL_2(\F_\ell)$ stated as \Cref{thm:reps:dickson} (Propositions 15 and 16 there, as in
\Cref{ch:groups}). Lattices in $\ell$-adic representations, their reductions and the independence of $\rhobar^{ss}$ from the lattice are in
\cite[Chapter I]{SerreAbelian}; the homothety lemma \Cref{thm:reps:homothety} is the elementary half of that circle of ideas, and the converse
construction is Ribet's lemma \cite[Section 2]{Ribet76}, where it appears as a proposition about reducible residual representations (we recall
it as Proposition 2.1 but have not checked the number). The use of Brauer-Nesbitt together with Chebotarev for Galois representations,
as in Lemma lem:old, is standard; see \cite{DDT} and \Cref{ch:galois}.

The table of \Cref{prop:reps:gl23-table} and the facts of \Cref{prop:reps:ns3-induced} and \Cref{prop:reps:nns3} were computed by hand and checked by
enumerating $\GL_2(\F_3)$ with the script of \Cref{sec:reps:gl23} and with an independent Python enumeration (stable lines over $\F_3$ and $\F_9$,
commutants, intertwiners, projective images). The formulas of \Cref{prop:reps:adjoint}(iii), the map $\Xi$ of \Cref{prop:reps:sym2} and the Gram
matrix of \Cref{prop:reps:trace-form} were checked symbolically in PARI/GP \cite{PARI}, and the fixed spaces of \Cref{sec:reps:cyclic} on the six
integral matrices of \Cref{ex:reps:ex13} over $\Q$ and over $\F_3$.

Two corrections to the paper in this chapter's area were made in the revised version and are only recalled here: the isomorphism
$\rhobar\simeq\Ind\psi$ of Proposition prop:ns holds after extending scalars to $\F_9$ (\Cref{prop:reps:ns3-induced}), and in Lemma lem:euler the invariant
line for the types $\kod{III}$, $\kod{III}^*$ is spanned by $\tau$ itself, since $\tau$ need not be diagonalisable over $\Ql$ (\Cref{lem:reps:commutant}). The step
``non-abelian and faithful'' of the original proof of Proposition prop:ns, which does not suffice by the example of $\mathrm B(\F_3)$, is replaced in
the revised version by the commutator argument; \Cref{prop:reps:dim2-absirr} gives a second proof.
